% ============================================================
%  Constraint and Transformation
%  An Integrated Introduction to Algebra, Geometry,
%  Trigonometry, and Stoichiometry
%
%  Author: Flyxion
%  Second Edition (Revised and Expanded)
% ============================================================

\documentclass[11pt, openany]{book}

% ---- Core packages ------------------------------------------
\usepackage[T1]{fontenc}
\usepackage[utf8]{inputenc}
\usepackage[margin=1.15in, top=1.25in, bottom=1.25in]{geometry}
\usepackage{amsmath, amssymb, amsthm}
\usepackage{mathtools}
\usepackage{xcolor}
\usepackage{tikz}
\usetikzlibrary{arrows.meta, calc, shapes.geometric, positioning, decorations.markings}
\usepackage{pgfplots}
\pgfplotsset{compat=1.18}
\usepackage{booktabs}
\usepackage{array}
\usepackage{enumitem}
\usepackage{tcolorbox}
\tcbuselibrary{theorems, skins, breakable}
\usepackage{fancyhdr}
\usepackage{sectsty}
\usepackage{hyperref}
\usepackage{graphicx}
\usepackage{multicol}
\usepackage{caption}
\usepackage[numbers, sort&compress]{natbib}

% ---- Color palette ------------------------------------------
\definecolor{ctBlue}{RGB}{25, 90, 150}
\definecolor{ctTeal}{RGB}{0, 130, 130}
\definecolor{ctGold}{RGB}{180, 130, 20}
\definecolor{ctRed}{RGB}{170, 40, 40}
\definecolor{ctGray}{RGB}{90, 90, 95}
\definecolor{ctLightBlue}{RGB}{220, 235, 248}
\definecolor{ctLightTeal}{RGB}{215, 240, 238}
\definecolor{ctLightGold}{RGB}{252, 245, 215}
\definecolor{ctLightRed}{RGB}{250, 225, 222}
\definecolor{ctLightGray}{RGB}{245, 245, 247}
\definecolor{ctPurple}{RGB}{90, 40, 120}
\definecolor{ctLightPurple}{RGB}{235, 220, 248}

% ---- Section fonts ------------------------------------------
\chapterfont{\color{ctBlue}}
\sectionfont{\color{ctBlue}}
\subsectionfont{\color{ctGray}}

% ---- Hyperref -----------------------------------------------
\hypersetup{
  colorlinks=true, linkcolor=ctBlue, urlcolor=ctTeal,
  citecolor=ctGray,
  pdftitle={Constraint and Transformation},
  pdfauthor={Flyxion},
}

% ---- Headers / footers --------------------------------------
\pagestyle{fancy}
\fancyhf{}
\fancyhead[LE]{\small\color{ctGray}\textit{Constraint and Transformation}}
\fancyhead[RO]{\small\color{ctGray}\textit{\nouppercase{\leftmark}}}
\fancyfoot[C]{\small\color{ctGray}\thepage}
\renewcommand{\headrulewidth}{0.4pt}

% ============================================================
% THEOREM / PROOF ENVIRONMENTS
% ============================================================
\theoremstyle{plain}
\newtheorem{theorem}{Theorem}[chapter]
\newtheorem{proposition}[theorem]{Proposition}
\newtheorem{lemma}[theorem]{Lemma}
\newtheorem{corollary}[theorem]{Corollary}

\theoremstyle{definition}
\newtheorem{defn}[theorem]{Definition}

\theoremstyle{remark}
\newtheorem{remark}[theorem]{Remark}
\newtheorem{note}[theorem]{Note}

\newenvironment{proofsketch}{\noindent\textit{Proof Sketch.}\;}{\hfill$\square$\par\medskip}

% ============================================================
% TCOLORBOX ENVIRONMENTS
% ============================================================
\tcbset{
  principleStyle/.style={
    enhanced, breakable,
    colback=ctLightBlue, colframe=ctBlue,
    fonttitle=\bfseries\color{white}, colbacktitle=ctBlue,
    attach boxed title to top left={yshift=-2mm, xshift=4mm},
    boxed title style={colback=ctBlue, rounded corners},
    left=4pt, right=4pt, top=8pt, bottom=4pt
  },
  exampleStyle/.style={
    enhanced, breakable,
    colback=ctLightTeal, colframe=ctTeal,
    fonttitle=\bfseries\color{white}, colbacktitle=ctTeal,
    attach boxed title to top left={yshift=-2mm, xshift=4mm},
    boxed title style={colback=ctTeal, rounded corners},
    left=4pt, right=4pt, top=8pt, bottom=4pt
  },
  definitionStyle/.style={
    enhanced, breakable,
    colback=ctLightGold, colframe=ctGold,
    fonttitle=\bfseries\color{white}, colbacktitle=ctGold,
    attach boxed title to top left={yshift=-2mm, xshift=4mm},
    boxed title style={colback=ctGold, rounded corners},
    left=4pt, right=4pt, top=8pt, bottom=4pt
  },
  historyStyle/.style={
    enhanced, breakable,
    colback=ctLightGray, colframe=ctGray,
    fonttitle=\bfseries\color{white}, colbacktitle=ctGray,
    attach boxed title to top left={yshift=-2mm, xshift=4mm},
    boxed title style={colback=ctGray, rounded corners},
    left=4pt, right=4pt, top=8pt, bottom=4pt
  },
  appendixStyle/.style={
    enhanced, breakable,
    colback=ctLightPurple, colframe=ctPurple,
    fonttitle=\bfseries\color{white}, colbacktitle=ctPurple,
    attach boxed title to top left={yshift=-2mm, xshift=4mm},
    boxed title style={colback=ctPurple, rounded corners},
    left=4pt, right=4pt, top=8pt, bottom=4pt
  },
}

\newtcolorbox{principle}[1][]{principleStyle, title={Guiding Principle #1}}
\newtcolorbox{myexample}[1][]{exampleStyle, title={Worked Example #1}}
\newtcolorbox{mydefinition}[1][]{definitionStyle, title={Definition #1}}
\newtcolorbox{history}[1][]{historyStyle, title={Historical Note #1}}
\newtcolorbox{deepnote}[1][]{appendixStyle, title={Deeper Structure #1}}

% ============================================================
% EXERCISE ENVIRONMENTS
% ============================================================
\newcounter{exercise}[chapter]
\renewcommand{\theexercise}{\thechapter.\arabic{exercise}}

\newenvironment{practice}{%
  \refstepcounter{exercise}%
  \par\medskip\noindent
  \textbf{\color{ctTeal}Practice \theexercise.}\enskip}{\par}

\newenvironment{exercise}{%
  \refstepcounter{exercise}%
  \par\medskip\noindent
  \textbf{\color{ctBlue}Exercise \theexercise.}\enskip}{\par}

\newenvironment{challenge}{%
  \refstepcounter{exercise}%
  \par\medskip\noindent
  \textbf{\color{ctGold}Challenge \theexercise.}\enskip}{\par}

\newenvironment{exercises}[1][Exercises]{%
  \section*{#1}\addcontentsline{toc}{section}{#1}}{}

% ============================================================
% TIKZ STYLES
% ============================================================
\tikzset{
  qty/.style={circle, draw=ctGray, fill=ctLightBlue,
              minimum size=7mm, font=\small},
  eqedge/.style={thick, ctBlue},
  arrow/.style={-{Stealth[length=2mm]}, ctBlue, thick},
  varrow/.style={-{Stealth[length=2mm]}, ctTeal, thick},
}

% ============================================================
% MACROS
% ============================================================
\newcommand{\R}{\mathbb{R}}
\newcommand{\N}{\mathbb{N}}
\newcommand{\Z}{\mathbb{Z}}
\newcommand{\Q}{\mathbb{Q}}
\newcommand{\abs}[1]{\left\lvert#1\right\rvert}
\newcommand{\norm}[1]{\left\lVert#1\right\rVert}
\newcommand{\pd}[2]{\frac{\partial #1}{\partial #2}}
\newcommand{\mol}{{\,\text{mol}}}
\newcommand{\gpm}{{\,\text{g/mol}}}
\DeclareMathOperator{\Hom}{Hom}
\DeclareMathOperator{\Ob}{Ob}
\DeclareMathOperator{\id}{id}

% ============================================================
\begin{document}
\frontmatter

% ---- Title page --------------------------------------------
\begin{titlepage}
\centering\vspace*{1.8cm}
{\color{ctBlue}\rule{\linewidth}{1.5pt}}\\[0.5cm]
{\Huge\bfseries\color{ctBlue}Constraint and Transformation}\\[0.4cm]
{\LARGE\color{ctGray}An Integrated Introduction to}\\[0.2cm]
{\LARGE\bfseries\color{ctTeal}
  Algebra, Geometry, Trigonometry,\\[0.1cm]
  and Chemical Stoichiometry}\\[0.5cm]
{\color{ctBlue}\rule{\linewidth}{1.5pt}}\\[1.8cm]

\begin{tikzpicture}[scale=2.0]
  % Unit circle
  \draw[ctBlue!70, thick] (0,0) circle (1);
  \draw[ctGray!40, thin] (-1.3,0)--(1.3,0);
  \draw[ctGray!40, thin] (0,-1.3)--(0,1.3);
  % Angle and point
  \def\ang{40}
  \coordinate (P) at ({cos(\ang)},{sin(\ang)});
  \draw[ctTeal, very thick] (0,0)--(P);
  \draw[ctRed, thick, dashed] ({cos(\ang)},0)--(P)
    node[midway, right, font=\scriptsize, ctRed]{$\sin\theta$};
  \draw[ctGold, thick, dashed] (0,0)--({cos(\ang)},0)
    node[midway, below, font=\scriptsize, ctGold]{$\cos\theta$};
  \draw[ctGray] (0.18,0) arc (0:\ang:0.18)
    node[midway, right=1pt, font=\scriptsize]{$\theta$};
  \filldraw[ctBlue] (P) circle (1.2pt);
  % Constraint label
  \node[font=\tiny, ctGray] at (0,-1.55)
    {$\sin^2\theta+\cos^2\theta=1$};
  \node[font=\tiny, ctGray, right] at (1.05,0) {$1$};
\end{tikzpicture}

\vfill
{\Large\scshape\color{ctGray}Flyxion}\\[0.25cm]
{\color{ctGold}\rule{5cm}{0.8pt}}
\end{titlepage}

% ---- Copyright ---------------------------------------------
\thispagestyle{empty}\vspace*{\fill}
\noindent\textit{Constraint and Transformation}\\
Copyright \textcopyright{} Flyxion. All rights reserved.\\
Typeset in \LaTeX\ with \texttt{pgfplots} and \texttt{tcolorbox}.
\vspace*{\fill}

% ---- ToC ---------------------------------------------------
\tableofcontents

% ============================================================
% NOTATION TABLE
% ============================================================
\chapter*{Notation and Symbols}
\addcontentsline{toc}{chapter}{Notation and Symbols}

The following symbols are used consistently throughout the text.
A change of meaning is always announced explicitly.

\begin{center}
\renewcommand{\arraystretch}{1.55}
\begin{tabular}{>{\bfseries}p{2.2cm} p{10cm}}
\toprule
Symbol & Meaning \\
\midrule
$\R$         & real number line \\
$\N$         & natural numbers $\{0,1,2,\dots\}$ \\
$x,y,z$      & generic real variables \\
$a,b,c$      & constants or geometric side lengths \\
$f,g,h$      & function names \\
$\theta$     & angle (radians unless noted) \\
$r$          & radius of a circle or sphere \\
$n$          & amount of substance (moles, mol) \\
$M$          & molar mass (g/mol) \\
$m$          & mass (grams or kilograms) \\
$L$          & luminosity of a star (watts) \\
$I$          & radiation intensity (W/m$^2$) \\
$d$          & distance \\
$A$          & area \\
$V$          & volume \\
$\Phi$       & scalar field (RSVP plenum density) \\
$\mathbf{v}$ & vector flow field \\
$S$          & entropy density field \\
$\gamma$     & stoichiometric coefficient vector \\
$N$          & element--species incidence matrix \\
$\mathcal{J}$ & variational functional \\
$\preceq$    & causal priority order on events \\
\bottomrule
\end{tabular}
\end{center}

% ============================================================
% PREFACE
% ============================================================
\chapter*{Preface}
\addcontentsline{toc}{chapter}{Preface}

This book is built on one conviction: algebra, geometry, trigonometry, and
chemistry are not separate disciplines accidentally packaged together by
tradition. They are four dialects of the same language --- the language of
\emph{constraint and transformation}.

Every equation encountered here is a claim about balance. Algebra encodes
symbolic balance among quantities. Geometry encodes spatial balance among
dimensions. Trigonometry encodes rotational and cyclic balance. Stoichiometry
encodes material balance in chemical reactions. To understand any one of these
is to recognize a pattern that recurs across all the others: a system remains
coherent because its constitutive relations must simultaneously be satisfied.

This orientation reflects a structural feature of scientific reasoning.
Quantitative systems are defined by relations that must remain internally
consistent as quantities change. A child discovering that a toy machine follows
a hidden rule, an astrophysicist reading the composition of a star from its
spectrum, a political theorist analyzing how collective power emerges from
coordination, a novelist constructing an ecologically stable imaginary world ---
each is solving, in different vocabulary, the same problem: \emph{what
constraints govern this system, and what transformations preserve them?}

\medskip\noindent\textbf{Intellectual influences.}
The book's perspective has been shaped by several intellectual traditions. The
developmental psychology of \citet{Gopnik2009} shows how causal inference in
children is structurally isomorphic with algebraic constraint satisfaction. The
astrophysics of \citet{PayneGaposchkin1925} demonstrates that geometric and
spectroscopic reasoning can reveal chemical composition at stellar distances.
The political philosophy of \citet{Arendt1958} offers a model in which social
systems are organized by relational structure rather than by fixed hierarchies.
The complexity science of \citet{Juarrero1999} formalizes how constraints
propagate through dynamical systems, while the speculative fiction of
\citet{LeGuin1974} explores the societal consequences of such constraint
structures with rare concreteness. The visual mathematical approach of
\citet{Needham1997,Needham2021} informs the book's emphasis on diagrammatic
reasoning.

\medskip\noindent\textbf{How to read this book.}
Each chapter follows a three-step rhythm: \emph{observation} (a phenomenon in
plain language), \emph{formalization} (the mathematical statement), and
\emph{extension} (the same idea in a new domain). The main text is
self-contained for an introductory reader. Each chapter ends with three tiers of
exercises: \textbf{Practice} (computation), \textbf{Exercise} (conceptual
reasoning), and \textbf{Challenge} (open-ended exploration). The appendices
develop the deeper formal architecture --- field theory, category theory, sheaf
semantics, variational principles, and recursive simulation --- for readers who
wish to see how elementary mathematics extends into contemporary research.

\medskip\noindent\textbf{Prerequisites.}
Arithmetic: fractions, decimals, percentages, signed numbers. No prior
algebra, geometry, or chemistry is assumed.

\mainmatter

% ============================================================
% INTRODUCTION
% ============================================================
\chapter{Mathematics as Structured Relation}

\section{The Central Thesis}

Scientific descriptions can be viewed as networks of relationships rather
than lists of isolated facts. This book proposes that the four subjects it
covers --- algebra, geometry, trigonometry, and chemistry --- are four
languages for a single idea: a system is defined not merely by its components
but by the \emph{relations} those components must satisfy.

\begin{mydefinition}[Constraint]
A \textbf{constraint} is any condition restricting the possible values of one
or more quantities. Formally, given variables $x_1,\dots,x_n$, a constraint is
a relation
\[
  F(x_1,\dots,x_n) = 0.
\]
The \textbf{solution set} is
$\mathcal{C}=\{(x_1,\dots,x_n)\in\R^n \mid F(x_1,\dots,x_n)=0\}.$
\end{mydefinition}

\begin{principle}
Every equation in this book is a claim that two descriptions of the same
situation must agree. Solving the equation means finding the values that live
inside the solution set $\mathcal{C}$.
\end{principle}

\section{Four Languages, One Pattern}

The table below introduces the central motif. Each row names a type of
constraint, the domain where it appears, and a representative equation.

\begin{center}
\renewcommand{\arraystretch}{1.6}
\begin{tabular}{>{\bfseries}p{2.8cm} p{3.2cm} p{5.8cm}}
\toprule
\color{ctBlue}Domain & \color{ctBlue}Constraint type & \color{ctBlue}Canonical equation \\
\midrule
Algebra       & symbolic / numerical   & $2x+3=11$ \\
Geometry      & spatial / dimensional  & $a^2+b^2=c^2$ \\
Trigonometry  & cyclic / angular       & $\sin^2\theta+\cos^2\theta=1$ \\
Stoichiometry & material / atomic      & $2\mathrm{H_2+O_2\to2H_2O}$ \\
\bottomrule
\end{tabular}
\end{center}

In every case the equation is a \emph{conservation statement}: it asserts that
some quantity --- a sum, a squared length, a trigonometric magnitude, or an
atom count --- is preserved under transformation.

\section{Constraint Networks}

A single equation imposes one constraint. Real systems impose many constraints
simultaneously. Representing them as a network reveals their mutual
dependencies.

\begin{center}
\begin{tikzpicture}[node distance=1.8cm]
  \node[qty] (x) {$x$};
  \node[qty, right=of x] (y) {$y$};
  \node[qty, below=of x] (z) {$z$};
  \node[qty, below=of y] (w) {$w$};
  \draw[eqedge] (x)--(y) node[midway,above,font=\scriptsize]{$F_1=0$};
  \draw[eqedge] (x)--(z) node[midway,left,font=\scriptsize]{$F_2=0$};
  \draw[eqedge] (y)--(w) node[midway,right,font=\scriptsize]{$F_3=0$};
  \draw[eqedge] (z)--(w) node[midway,below,font=\scriptsize]{$F_4=0$};
  \draw[eqedge] (y)--(z) node[midway,above right,font=\scriptsize]{$F_5=0$};
\end{tikzpicture}
\end{center}

Each node represents a quantity; each edge represents a constraint relating the
two adjacent quantities. A consistent system is one in which all constraints can
be satisfied simultaneously.

\begin{deepnote}
The general form of a linear constraint system is $A\mathbf{x}=\mathbf{b}$,
where $A$ is an $m\times n$ matrix. Appendix~A develops this idea formally,
including its connections to the RSVP field framework and Spherepop
event-history calculus.
\end{deepnote}

\section{The Book's Arc}

The book progresses from local to global. Part~I develops algebraic constraints
on individual quantities. Part~II extends this to spatial relationships. Part~III
describes cyclic and rotational constraints. Part~IV applies all of this to
chemical systems, which are constraint networks over atomic species. Part~V
synthesizes the four domains into a unified picture of structured systems.

\begin{exercises}
\begin{exercise}
Write down three everyday relationships and express each as an equation
$F(x,y)=0$. What does the solution set mean in each case?
\end{exercise}
\begin{challenge}
Draw a constraint network with five quantities and at least six constraints.
Choose a real system (a recipe, a budget, a simple machine) as your model.
Explain why some constraints are necessary and others may be redundant.
\end{challenge}
\end{exercises}

% ============================================================
% PART I --- ALGEBRA
% ============================================================
\part{Algebra: Relations and Constraint Systems}

\chapter{Numbers, Variables, and Representation}

\section*{From Counting to Representation}
\addcontentsline{toc}{section}{Opening Perspective}

Human beings began using numbers long before mathematics existed as a
formal discipline. Early societies counted livestock, measured land,
tracked seasons, and divided resources among households. In each case
numbers served a practical role: they were tools for describing the
quantitative structure of the world.

At first, numbers were attached to specific objects. Five sheep meant
five actual animals standing in a field. Three days meant three rotations
of the Earth between sunrise and sunrise. Gradually, however, a profound
conceptual shift occurred. People began to treat numbers not merely as
labels for physical things but as abstract entities that could be
manipulated independently of the objects they described.

This transition—from counting objects to reasoning about numbers
themselves—marks the birth of mathematics.

\medskip

\textbf{Representation as the foundation of mathematics.}

The central idea of this chapter is that numbers allow us to represent
quantitative relationships. Representation means that a symbolic
structure stands in for a physical or conceptual situation. A number can
represent the length of a road, the temperature of the air, the number of
molecules in a gas sample, or the brightness of a distant star.

Once a quantity is represented symbolically, we can manipulate the
symbols according to consistent rules. Arithmetic operations such as
addition and multiplication are not merely computational procedures.
They encode structural facts about quantities themselves. If two lengths
are combined end to end, their numerical representations add. If an
object is scaled by a factor of three, the number describing its length
is multiplied by three. Arithmetic therefore reflects the structure of
the physical relationships that numbers represent.

In this way mathematics becomes a language for describing patterns of
constraint.

\medskip

\textbf{Variables and the power of abstraction.}

The next conceptual step beyond numbers is the introduction of variables.
A variable is a symbol that stands for a quantity whose value may change
or may be unknown. Instead of describing a single specific situation, a
variable allows us to describe an entire family of situations at once.

For example, consider the distance travelled by a car moving at constant
speed. If the speed is 60 kilometres per hour, then the distance
travelled depends on the amount of time that passes. Writing the
relationship as

\[
d = 60t
\]

expresses not just one measurement but infinitely many possible
measurements. For every possible value of the time variable \(t\), the
distance \(d\) is determined by the same rule.

This is the essence of algebra: relationships between quantities are
expressed symbolically so that the structure of the relationship becomes
visible.

\medskip

\textbf{From particular facts to general laws.}

Before algebraic notation was developed, mathematicians described
relationships in words. Babylonian and Greek texts contain lengthy
verbal descriptions of geometric and arithmetic procedures. The
introduction of symbolic algebra in the early modern period transformed
mathematics because it allowed general laws to be expressed concisely.

A single equation such as

\[
y = kx
\]

encodes an entire class of proportional relationships. The same equation
can describe the distance travelled by a moving object, the amount of
chemical produced in a reaction, or the intensity of light spreading
through space. What changes from one application to another is the
interpretation of the variables.

In other words, algebra separates structure from context.

\medskip

\textbf{Numbers as coordinates in a space of possibilities.}

One of the most powerful ways to visualize numbers is to place them on a
line. The number line transforms arithmetic into geometry: addition
becomes movement to the right, subtraction becomes movement to the
left, and distance between numbers corresponds to magnitude.

This geometric viewpoint will become increasingly important as the book
progresses. Algebraic equations describe sets of points on the number
line or in the plane. Systems of equations describe intersections of
these sets. Later chapters will extend this idea further, showing how
geometric, trigonometric, and chemical relationships all arise from the
same underlying structure of constraints among quantities.

\medskip

\textbf{Constraint and transformation.}

Throughout this book we will repeatedly encounter two complementary
ideas.

A \emph{constraint} is a rule that restricts which combinations of
quantities are possible. An equation such as \(d = 60t\) is a constraint
because it specifies the relationship that distance and time must
satisfy for an object moving at constant speed.

A \emph{transformation} is a change in the quantities of a system that
preserves the constraint. When time increases, distance increases in
exactly the way required by the equation. The equation itself remains
true throughout the transformation.

Seen from this perspective, mathematics is the study of structures that
remain consistent under change. Numbers provide the language in which
those structures can be expressed.

\medskip

The sections that follow introduce the basic elements of this language.
We begin with numbers as representations of quantities, then introduce
variables that allow relationships to be expressed in general form. The
number line provides a geometric picture of these ideas, and the concept
of proportionality gives our first example of a mathematical law linking
two quantities together.

From these simple ingredients the entire machinery of algebra will grow.

\section{Numbers as Quantities}

A \textbf{number} is a symbolic stand-in for a measurable quantity: a count, a
length, a temperature. Numbers obey arithmetic: addition, subtraction,
multiplication, division.

\begin{mydefinition}[Variable]
A \textbf{variable} is a letter that stands for an unknown or varying quantity.
An \textbf{expression} is a combination of numbers and variables using
arithmetic operations. An expression does not assert equality; it simply
describes a quantity.
\end{mydefinition}

The expression $3x+5$ has a specific numerical value for each specific value
of $x$. It is not an equation and does not claim anything --- it is simply a
description.

\section{The Real Number Line}

Every real number corresponds to a unique point on an infinite line. Distance
from the origin measures magnitude; direction indicates sign.

\begin{center}
\begin{tikzpicture}[scale=1.1]
  \draw[arrow] (-3.5,0)--(3.8,0) node[right,font=\small]{$\R$};
  \foreach \x/\l in {-3/-3,-2/-2,-1/-1,0/0,1/1,2/2,3/3}
    \draw (\x,0.08)--(\x,-0.08) node[below,font=\scriptsize]{\l};
  \filldraw[ctRed] (2.3,0) circle (2.2pt) node[above,font=\scriptsize,ctRed]{$2.3$};
  \filldraw[ctTeal] (-1.6,0) circle (2.2pt) node[above,font=\scriptsize,ctTeal]{$-1.6$};
\end{tikzpicture}
\end{center}

\section{Proportionality: The Simplest Relation}

\begin{theorem}[Direct Proportionality]
If two quantities $x$ and $y$ satisfy $y=kx$ for a constant $k\ne0$, we say
$y$ is \textbf{directly proportional} to $x$ with proportionality constant $k$.
The graph of this relation is a straight line through the origin.
\end{theorem}

\begin{myexample}
A car travels at a constant speed of 60~km/h. Distance $d$ (km) and time $t$
(h) satisfy
\[
  d = 60\,t.
\]
Checking units: $\text{km} = (\text{km/h})\cdot\text{h}$. \checkmark

The slope $k=60$ is the rate of change: each additional hour yields 60 more
kilometres. For $t=2.5$~h, $d=150$~km.
\end{myexample}

\begin{history}[Alison Gopnik and Causal Inference]
In her research on how children understand the world, Alison Gopnik showed
that even toddlers spontaneously construct algebraic-style rules from
observation \citep{Gopnik2009}. A child who notices that a machine lights up
whenever two specific blocks are inserted together is, in effect, discovering a
proportional or Boolean rule governing the system. Elementary algebra makes
this intuition precise.
\end{history}

\begin{exercises}
\begin{practice}
A recipe uses 3 cups of flour per 2 cups of sugar. Write an equation relating
flour $f$ to sugar $s$. How much sugar is needed for 9 cups of flour?
\end{practice}
\begin{exercise}
A toy machine activates according to $A = 3B + 2$, where $B$ is the number of
blue blocks inserted. If the observed activation value is $A=14$, determine $B$
and interpret the result.
\end{exercise}
\begin{challenge}
Collect five proportional relationships from daily life (speed and distance,
ingredients in a recipe, etc.) and estimate each proportionality constant.
Which relationships are only approximately proportional? What causes the
deviation?
\end{challenge}
\end{exercises}

\chapter{Equations and Balance}

\section*{Balance as the Logic of Equations}
\addcontentsline{toc}{section}{Opening Perspective}

In the previous chapter we introduced numbers as representations of
quantities and variables as symbols that allow relationships between
quantities to be written in general form. Once variables enter the
picture, mathematics acquires a new kind of expressive power. We can
describe not only specific measurements but entire classes of situations.
The central tool that makes this possible is the equation.

An equation is a statement that two expressions represent the same
quantity. At first glance this seems almost trivial. Writing $2+3=5$
merely confirms something already known. Yet the conceptual importance
of the equals sign cannot be overstated. It asserts that two different
ways of describing a situation must ultimately agree.

Seen in this light, an equation is not merely a piece of notation but a
claim about the structure of a system.

\medskip

\textbf{Equality as balance.}

One of the most helpful ways to understand an equation is to imagine a
balance scale. A balance scale compares two masses placed on opposite
sides. When the scale is level, the masses must be equal. If one side is
heavier, the scale tilts until the imbalance is corrected.

An equation behaves in exactly the same way. The two sides of the
equation correspond to the two sides of the scale. The equality sign
indicates that the system is in equilibrium: whatever quantity appears
on the left must match the quantity on the right.

If we alter one side of the equation without making the same alteration
to the other, the balance is broken. To preserve equality we must perform
the same operation on both sides simultaneously. This simple principle
lies at the heart of all algebraic manipulation.

\medskip

\textbf{Solving equations as restoring equilibrium.}

Most equations encountered in algebra contain unknown quantities. The
purpose of solving an equation is to determine which value of the
variable restores balance between the two sides.

For example, suppose the equation

\[
2x + 3 = 11
\]

represents a relationship between two quantities. The value of $x$ is not
immediately obvious, but the equation tells us that whatever number $x$
represents must make the left-hand side equal to eleven.

Solving the equation means systematically transforming it until the
variable stands alone. Each transformation preserves the balance of the
equation while making its structure clearer. Eventually the value of the
unknown becomes visible.

In this sense algebraic problem solving is not a matter of guessing
numbers but of applying operations that preserve equality.

\medskip

\textbf{Equations as constraints.}

Within the broader framework of this book, equations play the role of
constraints. A constraint is a rule restricting which combinations of
quantities are possible. The equation $2x+3=11$ does not permit $x$ to
take arbitrary values. Only one value satisfies the condition imposed by
the equation.

More complicated systems may contain many equations simultaneously. Each
equation contributes its own constraint, and the solution to the system
must satisfy all of them at once. Later chapters will show that geometry,
trigonometry, and chemistry all rely on exactly this idea. The
Pythagorean theorem, the identity $\sin^2\theta+\cos^2\theta=1$, and the
conservation of atoms in a chemical reaction are all equations imposing
constraints on the quantities involved.

Thus the humble algebraic equation provides a prototype for the laws of
science.

\medskip

\textbf{Units and dimensional consistency.}

Equations describing real physical systems must obey an additional
principle: the units on both sides must match. If one side of an equation
is measured in metres, the other side must also represent a length. If
one side is measured in joules, the other side must represent energy.

This requirement is known as dimensional consistency. It acts as a
powerful error-checking mechanism. If the units do not agree, something
in the reasoning must be wrong.

Dimensional analysis therefore complements algebraic manipulation. While
algebra preserves equality between numerical expressions, dimensional
analysis preserves consistency between the physical quantities those
numbers represent.

\medskip

\textbf{Equations as the language of science.}

Nearly every scientific law can be expressed as an equation. Newton’s
laws of motion, the ideal gas law, the conservation of energy, and the
chemical equations governing reactions are all statements that certain
quantities must remain in balance.

The algebra developed in this chapter provides the basic techniques
needed to work with such relationships. By learning how to manipulate
equations and isolate unknown quantities, we gain the ability to predict
how a system will behave under changing conditions.

In the sections that follow we begin with the formal definition of an
equation and examine how its structure reflects the idea of balance.
From there we develop systematic methods for solving linear equations and
apply them to problems drawn from everyday reasoning and scientific
measurement.

What begins as a simple principle of equality will ultimately become the
foundation for describing the structure of the physical world.

\section{Equations as Statements of Equality}

\begin{mydefinition}[Equation and Solution]
An \textbf{equation} is a statement $A=B$ where $A$ and $B$ are expressions.
A \textbf{solution} is any assignment of values to the variables that makes
the statement true.
\end{mydefinition}

An equation is therefore a constraint: it asserts that exactly one value (or
set of values) satisfies the balance condition.

\subsection{The Balance Model}

\begin{center}
\begin{tikzpicture}[scale=0.9]
  \draw[ctBlue, very thick] (-3.1,0.9)--(3.1,0.9);
  \draw[ctBlue, thick] (0,0.9)--(0,-0.2);
  \draw[ctGray] (-0.3,-0.2)--(0.3,-0.2);
  \draw[ctTeal,thick] (-3.1,0.9)--(-3.1,0.4);
  \draw[ctTeal,thick] (3.1,0.9)--(3.1,0.4);
  \draw[ctTeal,thick] (-3.5,0.4)--(-2.7,0.4);
  \draw[ctTeal,thick] (2.7,0.4)--(3.5,0.4);
  \node at (-3.1,0.1){\small$2x+3$};
  \node at (3.1,0.1){\small$11$};
  \node[ctGray,font=\small] at (0,-0.55){balance};
\end{tikzpicture}
\end{center}

Whatever operation is applied to one side of the equation must be applied
identically to the other --- the scale must remain balanced.

\section{Solving Linear Equations}

\begin{theorem}[Isolation Principle]
A linear equation $ax+b=c$ (with $a\ne0$) has the unique solution
\[
  x = \frac{c-b}{a}.
\]
\end{theorem}

\begin{proofsketch}
Subtract $b$ from both sides: $ax=c-b$. Divide both sides by $a\ne0$.
\end{proofsketch}

\begin{myexample}[Balancing a simple equation]
Solve $2x+3=11$.
\begin{align*}
  2x+3 &= 11 && \text{original constraint}\\
  2x   &= 8  && \text{subtract }3\\
  x    &= 4  && \text{divide by }2
\end{align*}
\textit{Verification:} $2(4)+3=11$. \checkmark
\end{myexample}

\begin{myexample}[Gopnik: Hypothesis testing as equation solving]
A child hypothesises that a toy machine follows the rule $A=3G+Y$, where $G$
and $Y$ are the numbers of green and yellow blocks. She observes $A=8$ and
inserts $Y=2$. Solving for $G$:
\[
  3G+2=8 \implies 3G=6 \implies G=2.
\]
\end{myexample}

\section{Dimensional Analysis}

\begin{principle}
Every valid equation involving physical quantities must be \textbf{dimensionally
consistent}: the units on both sides must be identical.
\end{principle}

\begin{myexample}[Converting units]
Convert 72~km/h to m/s:
\[
  72\;\frac{\text{km}}{\text{h}}
  \times\frac{1000\;\text{m}}{1\;\text{km}}
  \times\frac{1\;\text{h}}{3600\;\text{s}}
  = 20\;\frac{\text{m}}{\text{s}}.
\]
\end{myexample}

\begin{exercises}
\begin{practice}
Solve each equation.\\
(a) $5x-7=18$\quad (b) $\dfrac{x}{3}+4=9$\quad
(c) $-2x+10=0$\quad (d) $3(x-2)=15$.
\end{practice}
\begin{exercise}
A reaction produces 0.5~mol of gas per minute. Write the constraint equation
and determine how many moles are produced in 2.5 hours. Include units at every
step.
\end{exercise}
\begin{challenge}
Arendt \citep{Arendt1958} argues that political power $P$ emerges from
coordination, not mere numbers. Model $P=cn$ where $c$ is coordination
strength and $n$ is the number of participants. If $c=0.2$ initially and the
group triples its coordination to $c=0.6$, by what factor does $P$ change?
What does this say about collective action?
\end{challenge}
\end{exercises}

\chapter{Proportions and Scaling}

\section*{Patterns That Scale}
\addcontentsline{toc}{section}{Opening Perspective}

Many relationships in nature do not depend on the absolute size of a
system but on the ratios between its parts. If two ingredients in a
recipe must be mixed in a ratio of three to two, the taste of the dish
remains the same whether we cook for two people or for two hundred. If a
map is drawn at a certain scale, the shapes of mountains and rivers are
preserved even though their distances are reduced. When astronomers
analyze the light of distant stars, they compare the relative strengths
of spectral lines rather than their absolute brightness.

In each case the essential information lies in a proportion.

\medskip

\textbf{Ratios as comparative measurements.}

A ratio compares two quantities by expressing how large one is relative
to the other. When we write a ratio such as $3:2$ or $\frac{3}{2}$, we are
not describing the absolute size of either quantity individually. We are
describing their relationship.

Ratios therefore capture the structure of a system rather than its
magnitude. Doubling both quantities in the ratio $3:2$ produces $6:4$,
yet the relationship remains unchanged because the proportion between
the quantities is preserved.

This invariance under scaling makes ratios an especially powerful
mathematical tool.

\medskip

\textbf{Proportions as constraints.}

A proportion is a statement that two ratios are equal. Writing

\[
\frac{a}{b}=\frac{c}{d}
\]

asserts that the relationship between $a$ and $b$ is identical to the
relationship between $c$ and $d$. Such equations appear constantly in
science and everyday reasoning. The gears of a machine rotate in fixed
ratios determined by their teeth. The concentration of a chemical
solution expresses the ratio between solute and solvent. The brightness
of a star measured through different spectral lines reveals the
proportional abundance of its elements.

Proportions therefore act as constraints linking multiple quantities
together.

\medskip

\textbf{Scaling and the structure of space.}

Ratios become even more powerful when we consider how systems change
size. Imagine enlarging a geometric figure such as a triangle. All of
its side lengths increase by the same factor, but the angles remain the
same. The figure preserves its shape even though its size changes.

This property is known as similarity, and it arises because the
relationships between the lengths of the sides remain proportional.
Geometry therefore provides a visual illustration of proportional
reasoning: scaling a system preserves its structure when ratios remain
constant.

Scaling relationships also govern many physical processes. The area of a
square grows with the square of its side length. The volume of a cube
grows with the cube of its side length. The brightness of a star
observed from Earth decreases with the square of the distance from the
observer. In each case a simple change in size produces a predictable
change in another quantity.

Understanding these scaling laws is one of the most important goals of
mathematics in science.

\medskip

\textbf{Proportional reasoning in scientific discovery.}

Historically, many major scientific discoveries relied on recognizing
proportional relationships. In astronomy, the intensity of spectral
lines in starlight reveals the relative abundance of chemical elements.
Cecilia Payne-Gaposchkin's groundbreaking work showed that the Sun and
most stars consist primarily of hydrogen and helium. Her conclusion
followed from proportional comparisons between observed spectral lines
and theoretical predictions.

Similar reasoning appears in physics, chemistry, biology, and economics.
When scientists compare measurements across systems of different sizes,
they often find that the same proportional pattern repeats again and
again.

Such patterns hint at underlying laws governing the system.

\medskip

\textbf{From ratios to scaling laws.}

The concept of proportionality can be extended even further. Sometimes
one quantity does not increase in direct proportion to another but
instead follows a power relationship such as

\[
y = kx^p .
\]

When the input quantity $x$ changes by some factor, the output quantity
changes by that factor raised to a power. This phenomenon, known as
power scaling, appears across an extraordinary range of natural
processes.

For example, doubling the radius of a circle quadruples its area. The
relationship between radius and area therefore follows a square law.
Similarly, the energy emitted by a star depends on the fourth power of
its temperature, a relationship known as the Stefan–Boltzmann law.

Power scaling reveals how quantities transform as systems grow or
shrink.

\medskip

\textbf{Scaling in human systems.}

Scaling relationships do not apply only to natural phenomena. They also
shape the organization of societies and ecosystems. The carrying
capacity of an island, the distribution of resources in a community, and
the growth of cities all depend on how space and resources scale with
population.

Ursula K.~Le Guin's speculative fiction explores such questions in
remarkably concrete ways. In her imagined societies, ecological balance
often depends on proportional relationships between land area, resource
availability, and population size. Mathematical reasoning about scaling
therefore becomes a tool for understanding how complex systems remain
sustainable.

\medskip

In this chapter we develop the mathematics needed to work with ratios,
proportions, and scaling laws. We begin by examining how ratios express
relationships between quantities and how proportions allow unknown
quantities to be determined. We then extend these ideas to power laws
that describe how systems change when their size changes.

These concepts will reappear throughout the book. Geometry relies on
similarity and scale. Trigonometry describes cyclic relationships that
repeat proportionally. Chemistry uses stoichiometric ratios to balance
reactions. In each domain the same principle applies: understanding the
relative relationships between quantities often reveals more about a
system than measuring their absolute values.

\section{Ratios and Proportions}

A \textbf{ratio} compares two quantities. A \textbf{proportion} asserts that
two ratios are equal:
\[
  \frac{a}{b}=\frac{c}{d}.
\]

\begin{myexample}[Stellar mass ratios]
In the hydrogen fusion reaction $4\mathrm{H}\to\mathrm{He}$, the mass ratio of
hydrogen consumed to helium produced is approximately $4:1$ (by nucleon count).
If a stellar region contains $10^{12}$ hydrogen nuclei per second, the
production rate of helium nuclei satisfies
\[
  \frac{n_\mathrm{He}}{1}=\frac{n_\mathrm{H}}{4}
  \implies n_\mathrm{He}=\frac{10^{12}}{4}=2.5\times10^{11}\text{ nuclei/s}.
\]
\end{myexample}

\begin{history}[Cecilia Payne-Gaposchkin]
Cecilia Payne-Gaposchkin (1900--1979) discovered in her 1925 doctoral
dissertation that stars are composed primarily of hydrogen and helium --- a
result initially rejected by the astronomical establishment but later
recognised as one of the most important discoveries in astrophysics
\citep{PayneGaposchkin1925}. Her analysis relied on proportional reasoning
connecting spectral line intensities to elemental abundances.
\end{history}

\section{Scaling Laws}

\begin{theorem}[Power Scaling]
If quantity $y$ scales as the $p$-th power of quantity $x$, then $y=kx^p$.
When $x$ is multiplied by a factor $\lambda$, $y$ is multiplied by $\lambda^p$.
\end{theorem}

\begin{myexample}[Le Guin: island resource capacity]
In a circular island settlement of radius $r=2$~km, the total area is
$A=\pi r^2=4\pi\approx12.6\;\text{km}^2$. If agricultural land occupies
half the island and each km$^2$ supports 500 people:
\[
  \text{capacity} = 500\times\frac{4\pi}{2}\times100
  = 500\times200\pi\approx314{,}159\text{ people}.
\]
Doubling the radius to $r=4$~km quadruples the area and thus quadruples the
capacity. Scaling is not linear \citep{LeGuin1974}.
\end{myexample}

\begin{exercises}
\begin{practice}
A car uses 8~L per 100~km. How far can it travel on 22~L?
\end{practice}
\begin{exercise}
A solution is prepared at 5~g of salt per 200~mL. How much salt is needed
for 500~mL? Write the proportion and solve.
\end{exercise}
\begin{challenge}
Le Guin's utopian novel \textit{The Dispossessed} \citep{LeGuin1974} depicts
a society on an arid moon that must balance population against scarce
resources. Suppose the moon's habitable area is $A$~km$^2$, each km$^2$
supports at most $p$ people, and the society aims to keep population at $80\%$
of capacity. Express the target population as a function of $A$ and $p$.
How does the target change if a drought reduces $A$ by $15\%$?
\end{challenge}
\end{exercises}

\chapter{Systems of Equations}

\section*{Networks of Constraints}
\addcontentsline{toc}{section}{Opening Perspective}

In the previous chapters we studied equations involving a single
unknown quantity. Such equations impose a constraint on the possible
values that the variable may take. Solving the equation means finding
the value that satisfies the condition imposed by the relationship.

Many real systems, however, cannot be described by a single equation.
They involve several quantities that influence one another
simultaneously. The price of goods in a market depends on supply and
demand. The motion of planets depends on the gravitational forces of
many bodies at once. The composition of a chemical mixture depends on
the quantities of all substances present.

In situations like these, the quantities involved form a web of
relationships rather than a single isolated rule.

\medskip

\textbf{From individual equations to systems.}

A system of equations is a collection of constraints that must all be
satisfied at the same time. Each equation expresses one relationship
between the variables. The full solution must satisfy every equation in
the system simultaneously.

This requirement dramatically narrows the set of possible solutions. A
single equation in two variables usually describes infinitely many
points. But when two independent equations are imposed together, only
the points that satisfy both relations remain possible.

In this sense a system of equations acts like a set of intersecting
filters, progressively restricting the set of admissible states.

\medskip

\textbf{The geometric viewpoint.}

One of the most illuminating ways to understand systems of equations is
to interpret them geometrically. A linear equation such as

\[
ax + by = c
\]

represents a straight line in the coordinate plane. Every point on that
line satisfies the equation.

If we introduce a second equation involving the same variables, it
defines a second line. The point where the two lines intersect is the
unique pair of values that satisfies both equations at once.

This geometric picture reveals the logic of the system immediately. If
the lines intersect at one point, the system has a unique solution. If
the lines are parallel, they never meet, and the system has no solution.
If the lines coincide, every point on the line satisfies both equations,
producing infinitely many solutions.

Algebra and geometry are therefore describing the same structure in
different languages.

\medskip

\textbf{Constraint propagation in complex systems.}

Systems of equations also provide a simple model of how constraints
propagate through larger networks. When one relationship in a system
changes, the other relationships must adjust in order to maintain
consistency.

Consider a system in which food production depends on labor, and energy
production depends on both labor and food supply. A reduction in labor
does not merely affect the first equation; it ripples through the entire
system. The quantities connected by the equations must reorganize
themselves until all constraints are again satisfied.

Philosopher of complexity Alicia Juarrero describes this phenomenon as
the propagation of constraints through a network of relations. In a
well-structured system, the relationships between components determine
how the system responds to change.

Mathematically, systems of equations provide one of the simplest
representations of this idea.

\medskip

\textbf{From small systems to scientific models.}

Although the examples in this chapter involve only two variables and two
equations, the same principles apply to much larger systems. Economists
model markets with dozens of simultaneous equations. Engineers analyze
electrical circuits by solving systems relating currents and voltages.
Chemists balance reactions by writing conservation equations for each
element involved.

Even modern physical theories often begin as systems of equations
relating multiple interacting quantities.

The methods developed in this chapter therefore form the foundation of
many areas of scientific reasoning.

\medskip

In the sections that follow we first define what it means for several
equations to form a system. We then explore the geometric interpretation
of such systems and learn algebraic techniques for solving them. By
viewing equations as interconnected constraints, we begin to see how
mathematical relationships organize entire networks of quantities rather
than isolated variables.

\section{Multiple Constraints, Multiple Unknowns}

When a situation involves two or more unknowns, a single equation constrains
but does not determine. A system of $n$ equations in $n$ unknowns typically
yields a unique solution.

\begin{mydefinition}[System of Equations]
A \textbf{system of $m$ equations in $n$ unknowns} is a collection
$\{F_1=0,\dots,F_m=0\}$ that must all hold simultaneously. The
\textbf{solution} is any tuple satisfying every equation at once.
\end{mydefinition}

\section{Geometric Interpretation}

Each linear equation $ax+by=c$ defines a \textbf{line} in the $xy$-plane. A
system of two such equations corresponds to two lines; the solution is their
intersection.

\begin{center}
\begin{tikzpicture}
\begin{axis}[
  width=7cm, height=6cm, xlabel={$x$}, ylabel={$y$},
  xmin=0,xmax=8, ymin=-1,ymax=10,
  axis lines=center, grid=both, grid style={ctGray!20},
  tick label style={font=\small},
]
  \addplot[ctBlue,thick,domain=0:8]{10-x}
    node[above right,font=\small]{$x+y=10$};
  \addplot[ctTeal,thick,domain=0:8]{3*x-2}
    node[below right,font=\small]{$3x-y=2$};
  \addplot[ctRed, only marks, mark=*, mark size=2.5pt]
    coordinates{(3,7)};
  \node[font=\small,ctRed,above right] at (axis cs:3,7){$(3,7)$};
\end{axis}
\end{tikzpicture}
\end{center}

\begin{myexample}[Juarrero: constraint propagation through a system]
Juarrero \citep{Juarrero1999} emphasises that constraints propagate through
systems: a change in one relation alters the whole. Consider a simple model:
\[
  \begin{cases}
    F = 2L \\
    E = F + L,
  \end{cases}
\]
with $L=100$. Then $F=200$, $E=300$. If a drought reduces $L$ to 80, then
$F=160$ and $E=240$. The constraint cascade propagates through the entire
system.
\end{myexample}

\begin{myexample}[Solving by elimination]
Solve $\{x+y=10,\;3x-y=2\}$.

Adding the two equations: $4x=12$, so $x=3$. Then $3+y=10$, so $y=7$.

\textit{Verification:} $3+7=10$ \checkmark;\; $9-7=2$ \checkmark.
\end{myexample}

\begin{exercises}
\begin{practice}
Solve the system $\{2x+3y=12,\; x-y=1\}$ by both substitution and elimination.
\end{practice}
\begin{exercise}
Interpret geometrically: (a) a system with no solution; (b) a system with
infinitely many solutions. What does each case mean for the constraint
network?
\end{exercise}
\begin{challenge}
A chemist mixes Solution~A (20\% acid) and Solution~B (50\% acid) to produce
300~mL of 35\% acid. Set up and solve the system. Connect this to Juarrero's
idea of constraints propagating through a mixture network.
\end{challenge}
\end{exercises}

\chapter{Matrices and Linear Transformations}

\section*{Transforming Systems}
\addcontentsline{toc}{section}{Opening Perspective}

In the previous chapter we studied systems of equations as collections of
constraints that must be satisfied simultaneously. Solving such systems
allowed us to determine the values of several unknown quantities at once.
Each equation imposed a restriction, and the intersection of those
restrictions determined the solution.

Although the algebraic techniques for solving systems are useful in
their own right, they also point toward a deeper mathematical idea.
Behind every system of linear equations lies a transformation that
relates one set of quantities to another.

Matrices provide a compact language for describing these transformations.

\medskip

\textbf{From equations to structure.}

Consider again a pair of linear equations such as

\[
a_1x + b_1y = c_1,
\qquad
a_2x + b_2y = c_2 .
\]

At first glance these appear to be two independent statements. Yet both
equations involve the same variables, and both share the same
coefficients arranged in a particular pattern. Writing these coefficients
as a rectangular array reveals the underlying structure that the
equations have in common.

This array is called a matrix.

A matrix collects numerical relationships into a single mathematical
object. Instead of working with several equations individually, we can
manipulate the entire system at once.

\medskip

\textbf{Matrices as operators on space.}

One of the most remarkable discoveries of nineteenth-century mathematics
was that matrices do more than organize equations. They also describe
transformations of geometric space.

When a matrix multiplies a vector, it sends that vector to a new
location. If we apply the same matrix to every point in the plane, the
entire plane is transformed. Lines may rotate, stretch, or shear, but
the relationships between points follow precise rules determined by the
matrix.

In this way matrices bridge algebra and geometry. A system of equations
can be interpreted either as a set of algebraic constraints or as a
geometric transformation of space.

\medskip

\textbf{Linear transformations.}

Transformations represented by matrices have a special property called
linearity. A transformation is linear if it preserves two fundamental
operations: addition and scalar multiplication.

If two vectors are added together before a linear transformation is
applied, the result is the same as transforming each vector first and
then adding the results. Similarly, multiplying a vector by a number
before the transformation produces the same outcome as transforming the
vector and then scaling the result.

These properties ensure that straight lines remain straight lines and
that the structure of space is preserved in a predictable way.

Because of this consistency, linear transformations form the foundation
of many areas of mathematics and science.

\medskip

\textbf{Matrices across scientific disciplines.}

Matrices appear wherever systems of relationships must be analyzed
simultaneously. In physics they describe rotations, reflections, and the
dynamics of coupled systems. In computer graphics they determine how
objects are rotated and scaled on a screen. In statistics they encode
relationships among many variables at once.

Chemistry provides another striking example. The balancing of chemical
reactions can be expressed using a stoichiometric matrix whose rows
represent chemical elements and whose columns represent molecular
species. Conservation of atoms then becomes the statement that the
matrix multiplied by the reaction coefficients equals zero.

This is exactly the same algebraic structure we encountered when solving
systems of equations.

\medskip

\textbf{From small matrices to large systems.}

Although the examples in this chapter involve small matrices, the same
principles apply to systems with hundreds or thousands of variables.
Modern scientific models often consist of enormous matrices describing
networks of interactions among components.

Climate models, neural networks, and economic simulations all rely on
linear algebra as a foundational tool. The ability to represent and
manipulate systems using matrices makes it possible to analyze complex
relationships that would otherwise be impossible to manage.

\medskip

The sections that follow introduce matrices formally and show how they
allow systems of equations to be written in a compact and powerful
notation. We then interpret matrices as transformations of geometric
space and examine how the same mathematical structure appears in the
analysis of chemical reactions.

By the end of the chapter, systems of equations will no longer appear as
isolated algebraic puzzles. Instead they will reveal themselves as
instances of a broader idea: linear transformations acting on spaces of
quantities.

\section{Systems in Matrix Form}

A system $\{a_1x+b_1y=c_1,\;a_2x+b_2y=c_2\}$ may be written as
\[
  A\mathbf{x}=\mathbf{b},
  \qquad
  A=\begin{pmatrix}a_1&b_1\\a_2&b_2\end{pmatrix},\quad
  \mathbf{x}=\begin{pmatrix}x\\y\end{pmatrix},\quad
  \mathbf{b}=\begin{pmatrix}c_1\\c_2\end{pmatrix}.
\]

\begin{mydefinition}[Matrix]
An $m\times n$ \textbf{matrix} is a rectangular array of $mn$ numbers
arranged in $m$ rows and $n$ columns. It encodes a linear transformation
$T:\R^n\to\R^m$ by $T(\mathbf{x})=A\mathbf{x}$.
\end{mydefinition}

\section{Matrices as Transformations of Space}

A $2\times2$ matrix $A$ transforms every point $(x,y)$ in the plane. The
unit square is sent to a parallelogram.

\begin{center}
\begin{tikzpicture}[scale=1.1]
  % Original square
  \fill[ctLightBlue] (0,0)--(1,0)--(1,1)--(0,1)--cycle;
  \draw[ctBlue,thick] (0,0)--(1,0)--(1,1)--(0,1)--cycle;
  \node[font=\scriptsize,ctBlue] at (0.5,0.5){unit};
  % Transformed (shear A = [[1,1],[0,1]])
  \begin{scope}[xshift=3.5cm]
    \fill[ctLightTeal] (0,0)--(1,0)--(2,1)--(1,1)--cycle;
    \draw[ctTeal,thick] (0,0)--(1,0)--(2,1)--(1,1)--cycle;
    \node[font=\scriptsize,ctTeal] at (1,0.5){$A\cdot\square$};
  \end{scope}
  % Arrow
  \draw[arrow] (1.3,0.5)--(3.2,0.5)
    node[midway,above,font=\scriptsize]{$A=\!\scriptsize\begin{pmatrix}1&1\\0&1\end{pmatrix}$};
\end{tikzpicture}
\end{center}

\section{The Stoichiometric Matrix}

Chemical reactions are encoded by a stoichiometric matrix $N$ whose rows
index elements and whose columns index species.

\begin{myexample}[Water formation as matrix equation]
For $2\mathrm{H_2}+\mathrm{O_2}\to2\mathrm{H_2O}$, writing
$\gamma=(-2,-1,2)^\top$ (negative for reactants):
\[
  N=\begin{pmatrix}2&0&2\\0&2&1\end{pmatrix},\qquad
  N\gamma=\begin{pmatrix}2(-2)+0(-1)+2(2)\\0(-2)+2(-1)+1(2)\end{pmatrix}
  =\begin{pmatrix}0\\0\end{pmatrix}.
\]
Conservation is encoded by $N\gamma=\mathbf{0}$.
\end{myexample}

\begin{deepnote}
Appendix~D develops stoichiometry as conservation algebra in full generality,
including the reaction rate equation $\dot{c}=Sr(c)$.
\end{deepnote}

\begin{exercises}
\begin{practice}
Multiply $\begin{pmatrix}2&1\\0&3\end{pmatrix}\begin{pmatrix}4\\-1\end{pmatrix}$.
\end{practice}
\begin{exercise}
Write the system $\{x+2y=7,\;3x-y=5\}$ as $A\mathbf{x}=\mathbf{b}$ and solve
using Cramer's rule or row reduction.
\end{exercise}
\begin{challenge}
Model Arendt's power equation $P=C\mathbf{n}$ as a matrix transformation,
where $\mathbf{n}$ is a vector of participants across three groups and $C$
is a $1\times3$ coordination matrix. What does the rank of $C$ represent?
\end{challenge}
\end{exercises}

% ============================================================
% PART II --- GEOMETRY
% ============================================================
\part{Geometry: Spatial Structure and Invariance}

\chapter{Points, Lines, and Distance}

\section*{Space as Structured Relation}
\addcontentsline{toc}{section}{Opening Perspective}

Up to this point the relationships we have studied have been expressed
primarily through algebraic symbols. Variables represented quantities,
equations expressed constraints among those quantities, and systems of
equations revealed how multiple constraints interact simultaneously.
Although these ideas may seem purely symbolic, they already contain the
seeds of geometry.

Geometry arises when quantities are interpreted as positions in space.

\medskip

\textbf{From numbers to locations.}

A single number can represent a position along a line. When we place
numbers on a number line, addition corresponds to movement along the
line and subtraction corresponds to movement in the opposite direction.
The numerical structure of arithmetic therefore acquires a spatial
interpretation.

Extending this idea produces the coordinate plane. Instead of describing
position with a single number, we describe it using an ordered pair of
numbers. One coordinate measures horizontal position and the other
measures vertical position. Together they determine a unique point in
the plane.

This simple idea, introduced by René Descartes in the seventeenth
century, unified two previously separate branches of mathematics:
algebra and geometry.

\medskip

\textbf{Geometry as the study of spatial constraints.}

Once points in space can be represented by coordinates, geometric
relationships can be expressed as equations. A line becomes the set of
all points satisfying a linear equation. A circle becomes the set of
points satisfying a quadratic relation. More complex curves correspond
to more complicated equations.

In this way geometry becomes another language for describing
constraints.

Instead of writing an equation and solving for a number, we can visualize
the entire set of solutions as a geometric figure. The shape of the
figure reflects the structure of the equation.

For example, a linear equation in two variables produces a straight
line, while a quadratic equation produces a curved shape such as a
circle or parabola. Algebra and geometry are therefore two perspectives
on the same underlying mathematical structure.

\medskip

\textbf{Distance and the geometry of space.}

Among the most fundamental questions in geometry is the concept of
distance. Given two points in space, how far apart are they? Measuring
distance allows us to compare lengths, describe shapes, and determine
the size of geometric figures.

Remarkably, the formula for distance in the coordinate plane follows
directly from the Pythagorean theorem. A pair of coordinates defines the
legs of a right triangle whose hypotenuse represents the distance
between the points. The relationship between algebra and geometry
therefore appears again: an equation describing squares of lengths
becomes a geometric measurement.

Distance is one of the key invariants of geometry. Certain
transformations of space—such as translations and rotations—preserve
distance exactly. Understanding these invariants will become central in
later chapters.

\medskip

\textbf{Angles and rotation.}

Another fundamental geometric concept is the angle. While distance
measures separation between points, angles measure the amount of
rotation between directions.

Angles appear in many contexts: the orientation of a line, the turning
of a wheel, the orbit of a planet, or the motion of waves and periodic
phenomena. Because of this wide applicability, mathematicians developed
a natural unit for measuring angles based on the geometry of circles.

This unit is the radian.

Unlike degrees, which divide a circle into an arbitrary number of
parts, the radian measures angles directly in terms of arc length. The
angle subtended by an arc equal to the radius of the circle is defined
to be one radian. This definition links angles to the geometry of
circles in a natural and mathematically convenient way.

\medskip

\textbf{Geometry as the bridge between algebra and physics.}

The coordinate methods introduced in this chapter form the foundation of
analytic geometry, a subject that plays a crucial role in nearly every
scientific discipline. Physicists describe motion using coordinates and
distances. Astronomers locate stars on celestial coordinate systems.
Engineers analyze forces and structures using geometric relationships
expressed algebraically.

By translating spatial relationships into equations, analytic geometry
allows the methods of algebra to be applied directly to problems
involving space.

\medskip

In the sections that follow we introduce the coordinate plane and show
how points, lines, and geometric relationships can be expressed using
algebraic formulas. We then examine the concept of angle and introduce
the radian as the natural unit for measuring rotation.

These ideas provide the foundation for the geometric viewpoint that will
support the rest of the book. Just as equations describe numerical
constraints among quantities, geometric figures describe spatial
constraints among points. Understanding this connection will allow us to
move freely between algebraic reasoning and geometric intuition.

\section{The Coordinate Plane}

René Descartes' insight was that geometric objects can be described
algebraically. A point is an ordered pair $(x,y)$; a line is the solution
set of a linear equation; a circle is the solution set of a quadratic relation.

\begin{theorem}[Distance Formula]
The distance between points $(x_1,y_1)$ and $(x_2,y_2)$ in the plane is
\[
  d = \sqrt{(x_2-x_1)^2+(y_2-y_1)^2}.
\]
This formula is a consequence of the Pythagorean theorem, which we derive in
the next chapter.
\end{theorem}

\section{Angles and Their Measure}

\begin{mydefinition}[Radian]
One \textbf{radian} is the angle subtended at the centre of a circle by an
arc equal in length to the radius. Since the full circumference is $2\pi r$,
a full rotation equals $2\pi$ radians.
\end{mydefinition}

Conversion: $\theta_\text{rad}=\frac{\pi}{180}\theta_\text{deg}$.

\begin{exercises}
\begin{practice}
Convert: (a) $60°$ to radians; (b) $\frac{3\pi}{4}$ to degrees.
\end{practice}
\begin{exercise}
Show algebraically that the midpoint of the segment from $(x_1,y_1)$ to
$(x_2,y_2)$ is $\bigl(\frac{x_1+x_2}{2},\frac{y_1+y_2}{2}\bigr)$.
\end{exercise}
\end{exercises}

\chapter{Triangles and the Pythagorean Relation}

\section*{Geometry of Orthogonality}
\addcontentsline{toc}{section}{Opening Perspective}

In the previous chapter we introduced the coordinate plane and showed
how geometric objects can be described algebraically. Points became
ordered pairs of numbers, lines became the solution sets of linear
equations, and circles emerged from quadratic relations. Geometry and
algebra were revealed as two different descriptions of the same
mathematical structure.

The next step is to understand how distances in the plane are related to
one another.

Among all geometric figures, the triangle plays a uniquely fundamental
role. Any polygon can be decomposed into triangles, and the geometry of
curved surfaces can often be approximated by triangulating them into
many small triangular regions. Because of this universality, triangles
serve as the basic building blocks of much of geometry.

Within the family of triangles, right triangles occupy a special
position.

\medskip

\textbf{Orthogonality and right angles.}

A right triangle contains an angle of exactly ninety degrees. The two
sides forming that angle are called the legs of the triangle, and the
side opposite the right angle is called the hypotenuse.

This configuration introduces a fundamental geometric relationship:
the directions of the two legs are perpendicular, or orthogonal, to one
another.

Orthogonality appears in many areas of mathematics and science. The
horizontal and vertical axes of the coordinate plane are orthogonal.
The directions north and east used in navigation are orthogonal. In
physics, independent components of motion are often resolved along
orthogonal directions.

Right triangles therefore provide the simplest geometric model of
independent directions in space.

\medskip

\textbf{The Pythagorean relation.}

The central result governing right triangles is the Pythagorean theorem,
one of the oldest and most celebrated results in mathematics. It states
that the square of the hypotenuse equals the sum of the squares of the
two legs.

\[
a^2 + b^2 = c^2 .
\]

Although the theorem is ancient, its significance extends far beyond the
geometry of triangles. The equation expresses a deeper relationship
between orthogonal components of a system.

If two quantities represent independent directions, the magnitude of
their combined effect is determined by adding the squares of their
components. This idea reappears throughout mathematics and physics.

In analytic geometry it becomes the formula for the distance between two
points in the plane. In physics it appears in vector addition and in the
energy relations of mechanical systems. In statistics it emerges in the
definition of variance and Euclidean distance between data points.

The simple geometric picture of a right triangle therefore encodes a
general rule for combining independent contributions.

\medskip

\textbf{Squares and geometric measurement.}

The presence of squares in the Pythagorean theorem is not accidental.
Area plays a central role in the geometric interpretation of the
relationship.

The classical proof, illustrated in the figure that follows, shows that
the areas of the squares constructed on the legs of a right triangle
combine exactly to equal the area of the square constructed on the
hypotenuse. In other words, the theorem is a statement about how areas
transform when geometric figures are rearranged.

Many different proofs of the Pythagorean theorem exist, ranging from
purely geometric arguments to algebraic demonstrations. Each proof
reveals a different perspective on the same invariant relationship.

\medskip

\textbf{From triangles to spatial laws.}

Although the Pythagorean theorem arises from elementary geometry, its
influence extends into many branches of science. One important example
is the inverse-square law governing the spread of radiation, sound, and
gravitational influence.

When energy radiates outward from a source, it spreads over the surface
of an expanding sphere. The surface area of that sphere grows in
proportion to the square of its radius. As a result, the intensity of
the radiation decreases with the square of the distance from the source.

This geometric principle explains phenomena ranging from the brightness
of stars to the strength of gravitational attraction.

Astronomers such as Cecilia Payne-Gaposchkin relied on these geometric
relations to interpret measurements of stellar radiation and to infer
the physical properties of distant stars.

\medskip

In the sections that follow we examine right triangles in detail and
develop the Pythagorean theorem as a fundamental geometric relation.
From this theorem we derive practical methods for computing distances
and analyzing spatial relationships.

Although the theorem concerns a simple triangular figure, it expresses a
deep principle about the structure of space: independent directions
combine according to a precise geometric rule. Understanding this rule
will allow us to analyze a wide variety of problems involving distance,
measurement, and the geometry of physical systems.

\section{Right Triangles}

\begin{theorem}[Pythagorean Theorem]
In any right triangle with legs $a,b$ and hypotenuse $c$,
\[
  a^2+b^2=c^2.
\]
The length of the hypotenuse is completely determined by the two legs.
\end{theorem}

\begin{proofsketch}
Place four congruent right triangles inside a square of side $a+b$.
The uncovered central region is a square of side $c$.

\begin{center}
\begin{tikzpicture}[scale=1.5]
  \def\a{1.1}\def\b{0.7}
  \pgfmathsetmacro{\s}{\a+\b}
  \draw[ctGray,thick] (0,0)--(\s,0)--(\s,\s)--(0,\s)--cycle;
  \filldraw[ctLightBlue,draw=ctBlue]
    (0,0)--(\a,0)--(0,\b)--cycle;
  \filldraw[ctLightBlue,draw=ctBlue]
    (\s,0)--(\s,\a)--(\b,0)--cycle;
  \filldraw[ctLightBlue,draw=ctBlue]
    (\s,\s)--(0,\s)--(\s,\s-\b)--cycle;
  \filldraw[ctLightBlue,draw=ctBlue]
    (0,\s)--(0,\s-\a)--(\a,\s)--cycle;
  \draw[ctRed,very thick]
    (\a,0)--(\s,\b)--(\b,\s)--(0,\s-\a)--cycle;
  \node[font=\small,ctRed] at ({\s/2},{\s/2}){$c^2$};
  \node[font=\scriptsize,below] at (\a/2,0){$a$};
  \node[font=\scriptsize,below] at ({\a+\b/2},0){$b$};
\end{tikzpicture}
\end{center}
Area of outer square: $(a+b)^2=a^2+2ab+b^2$. Four triangles: $4\cdot\tfrac12ab=2ab$.
Inner square: $c^2=(a+b)^2-2ab=a^2+b^2$.
\end{proofsketch}

\section{Applications}

\begin{myexample}[Payne-Gaposchkin: radiation and distance]
Light from a star of luminosity $L$ (W) spreads uniformly over a sphere of
radius $r$. Surface area $A=4\pi r^2$ (m$^2$), so intensity is
\[
  I = \frac{L}{4\pi r^2}\;\text{W/m}^2.
\]
For the Sun: $L=3.8\times10^{26}$~W at $r=1.5\times10^{11}$~m:
\[
  4\pi r^2\approx 2.83\times10^{23}\text{ m}^2,\qquad
  I\approx1340\;\text{W/m}^2.
\]
The inverse-square law is a constraint on spatial geometry.
\end{myexample}

\begin{exercises}
\begin{practice}
Find the diagonal of a $7\times24$~cm rectangle.
\end{practice}
\begin{exercise}
Two ships leave the same port. One travels 40~km due north; the other 30~km
due east. How far apart are they? Draw and label the constraint diagram.
\end{exercise}
\begin{challenge}
Star~A has luminosity $2L$ and is at distance $3d$. Star~B has luminosity $L$
and is at distance $d$. Which appears brighter to an observer, and by what
factor? Generalise: express the apparent brightness ratio as a function of
the luminosity and distance ratios.
\end{challenge}
\end{exercises}

\chapter{Similarity and Scaling in Geometry}

\section*{Shape, Scale, and Invariance}
\addcontentsline{toc}{section}{Opening Perspective}

In the previous chapter we examined right triangles and the Pythagorean
relation, discovering how distances combine when directions are
orthogonal. The geometry of space revealed a precise rule connecting the
lengths of sides in a triangle. Yet another question naturally arises
once distance is understood: how do geometric figures behave when they
change size?

The study of similarity answers this question.

Similarity concerns the preservation of shape under changes of scale.
When two figures are similar, one can be obtained from the other by
enlarging or shrinking it while preserving the relative relationships
among its parts. Angles remain the same, and the ratios between
corresponding lengths remain constant.

This property reveals one of the most important ideas in mathematics:
structure can remain invariant even when scale changes.

\medskip

\textbf{Shape versus size.}

Two geometric figures may differ greatly in size yet still share exactly
the same shape. A small triangle drawn on paper and a large triangular
roof on a building may have identical angles and proportional sides.
Although their absolute dimensions differ, the pattern of relationships
between their sides is preserved.

Similarity therefore separates the concept of form from the concept of
magnitude. Geometry is not only concerned with measuring objects but
also with understanding the structural patterns that remain unchanged
when objects are enlarged or reduced.

This distinction between shape and size is central to many scientific
disciplines.

\medskip

\textbf{Scaling and dimensionality.}

When a figure is scaled by a factor $k$, every linear dimension becomes
$k$ times larger. Yet quantities that depend on two or three dimensions
change more dramatically.

Areas scale with the square of the scale factor, and volumes scale with
the cube of the scale factor. A square whose sides double in length has
four times the area. A cube whose edges double in length has eight times
the volume.

These relationships arise because area measures two-dimensional extent
and volume measures three-dimensional extent. Each independent spatial
dimension contributes an additional factor of the scale change.

The simple act of enlarging a figure therefore produces predictable
changes in geometric measurements.

\medskip

\textbf{Scaling laws in nature.}

The same principles appear throughout the natural world. The surface
area of a sphere grows with the square of its radius, while its volume
grows with the cube. These geometric relationships influence many
physical processes.

The brightness of a star decreases with the square of the distance from
the observer because the star's radiation spreads across the surface of
an expanding sphere. The strength of gravitational attraction follows a
similar inverse-square relationship.

Even biological systems are shaped by scaling laws. The surface area of
an organism determines how it exchanges heat with its environment,
while its volume determines its mass and metabolic needs. As organisms
increase in size, these quantities scale differently, placing
constraints on how large biological structures can become.

Geometry therefore provides a mathematical language for describing how
structure changes with size.

\medskip

\textbf{Similarity as an invariant transformation.}

From a geometric viewpoint, similarity transformations belong to a
family of operations that preserve certain properties of figures while
altering others. Rotations and reflections preserve both shape and
size. Translations move figures without altering their geometry at all.
Scaling transformations preserve shape but change size.

Understanding which quantities remain invariant under each type of
transformation allows mathematicians to classify geometric structures
according to their fundamental properties.

Similarity transformations preserve angles and proportional
relationships among lengths. These invariants make it possible to study
entire families of figures at once rather than analyzing each individual
instance separately.

\medskip

\textbf{From geometric similarity to mathematical models.}

The idea of similarity also plays a crucial role in modeling physical
systems. Engineers use scaled models of buildings, bridges, and aircraft
to study their behavior before constructing the full structures.
Astronomers interpret observations of distant objects by comparing them
to known physical models scaled to cosmic dimensions.

In each case the assumption of similarity allows small systems to reveal
the behavior of much larger ones.

Mathematically, similarity provides the bridge between geometry and the
scaling laws introduced earlier in the study of proportions. What began
as a simple ratio between numbers now becomes a transformation of
entire geometric figures.

\medskip

In the sections that follow we formalize the concept of similarity and
examine how scaling transformations affect geometric measurements such
as length, area, and volume. These ideas will allow us to analyze how
shapes change under enlargement and reduction while preserving their
essential structure.

By understanding similarity, we gain a powerful perspective on how
patterns persist across scales—from small geometric constructions to the
vast structures of the physical universe.

\section{Similar Figures}

\begin{mydefinition}[Similarity]
Two figures are \textbf{similar} (written $\sim$) if one can be obtained
from the other by scaling, rotation, or reflection. Corresponding angles are
equal; corresponding lengths are proportional with scale factor $k$.
\end{mydefinition}

\begin{theorem}[Scaling Laws]
If two similar figures have linear scale factor $k$, then areas scale as
$k^2$ and volumes scale as $k^3$.
\end{theorem}

\begin{proofsketch}
Area involves two length dimensions; volume involves three. Each scales as a
power of $k$ equal to the number of dimensions.
\end{proofsketch}

\begin{exercises}
\begin{practice}
Two similar triangles have scale factor $3:5$. If the smaller has area
$27\;\text{cm}^2$, find the area of the larger.
\end{practice}
\begin{exercise}
A spherical star with radius $r$ has surface area $4\pi r^2$. If the star
expands to radius $2r$, by what factor does its surface area increase? Its
volume?
\end{exercise}
\end{exercises}

\chapter{Transformations of the Plane}

\section*{Motion and Invariance in the Plane}
\addcontentsline{toc}{section}{Opening Perspective}

Up to this point our study of geometry has focused on the relationships
between static figures. Points, lines, triangles, and circles were
described by equations and measurements that determined their shape and
size. Yet geometry also concerns the ways figures may move within space
while preserving their fundamental structure.

A geometric transformation describes such motion.

Instead of asking what a figure is, we ask how it changes when every
point of the plane is mapped to a new position according to a rule.
This viewpoint shifts attention from individual shapes to the operations
that act upon them.

\medskip

\textbf{Geometry as transformation.}

In classical geometry figures were often studied as fixed objects. A
triangle had certain side lengths and angles, and the goal was to
discover relationships among those measurements. Modern geometry,
however, emphasizes the transformations that carry one figure to
another.

A transformation is a function that assigns a new location in the plane
to every point of the original figure. When the same rule is applied to
all points simultaneously, the entire geometric configuration moves in a
coherent way.

This perspective reveals an important idea: many geometric properties
remain unchanged under certain transformations. These unchanged
quantities are called invariants.

\medskip

\textbf{Distance as a geometric invariant.}

Among the most fundamental invariants of Euclidean geometry is distance.
If the distance between two points remains unchanged after a
transformation, then the transformation preserves the shape and size of
all figures in the plane.

Transformations with this property are called isometries.

When an isometry is applied to a geometric figure, the figure may move
to a new position or orientation, but its internal structure remains
exactly the same. Angles remain equal, lengths remain unchanged, and the
overall shape is preserved.

Isometries therefore describe motions of the plane that do not distort
geometry.

\medskip

\textbf{The basic motions of the plane.}

There are three fundamental ways a figure can move while preserving
distance. A figure may slide from one location to another without
changing its orientation. This motion is called a translation. A figure
may rotate about a point, turning through some angle while preserving
its shape. Finally, a figure may be reflected across a line, producing a
mirror image.

Every distance-preserving transformation of the plane can be constructed
from combinations of these three basic motions.

These operations form the foundation of the symmetry structure of
Euclidean space.

\medskip

\textbf{Transformations and matrices.}

The study of transformations connects geometry back to the algebraic
methods introduced earlier in the book. Many geometric transformations
can be represented by matrices acting on coordinate vectors. When a
matrix multiplies a coordinate pair $(x,y)$, it produces a new pair
representing the transformed point.

Rotations provide a particularly elegant example. A rotation through an
angle $\theta$ can be expressed using a matrix involving the trigonometric
functions $\cos\theta$ and $\sin\theta$. This representation allows
geometric motion to be analyzed using the tools of linear algebra.

In this way algebra and geometry once again converge.

\medskip

\textbf{Symmetry and invariance.}

The concept of invariance under transformation extends far beyond
geometry. In physics, symmetries of space and time determine the laws
governing motion and conservation. In chemistry, molecular structures
are classified according to their symmetry groups. In art and
architecture, patterns and designs rely on repeating transformations of
basic shapes.

Mathematically, the collection of transformations that preserve a given
set of properties forms a group, a structure that captures the internal
symmetry of a system.

The isometries of the plane therefore provide one of the simplest
examples of symmetry in mathematics.

\medskip

In the sections that follow we define transformations of the plane and
study the special class of isometries that preserve distance. We examine
translations, rotations, and reflections and explore how these
transformations can be combined to produce more complex motions.

Through this study we will see that geometry is not only the study of
shapes but also the study of the transformations that leave those shapes
unchanged. Understanding these invariances will prepare us for the next
stage of the book, where rotational motion leads naturally to the
periodic structures studied in trigonometry.

\section{Isometries: Transformations Preserving Distance}

A \textbf{transformation} is a function mapping the plane to itself. An
\textbf{isometry} preserves distances:
$d(T(P),T(Q))=d(P,Q)$ for all points $P,Q$.

The three basic isometries are:
\begin{itemize}
\item \textbf{Translation} by vector $\mathbf{v}$: $T(P)=P+\mathbf{v}$.
\item \textbf{Rotation} by angle $\theta$ around origin: encoded by $R(\theta)$.
\item \textbf{Reflection} across a line $\ell$.
\end{itemize}

\begin{center}
\begin{tikzpicture}[scale=0.85]
  % Original triangle
  \coordinate (A) at (0,0); \coordinate (B) at (1.2,0); \coordinate (C) at (0.6,1);
  \filldraw[ctLightBlue,draw=ctBlue] (A)--(B)--(C)--cycle;
  \node[font=\scriptsize,ctBlue] at (0.6,0.3){$T$};
  % Translated
  \begin{scope}[xshift=2.5cm]
    \coordinate (A') at (0,0); \coordinate (B') at (1.2,0); \coordinate (C') at (0.6,1);
    \filldraw[ctLightTeal,draw=ctTeal] (A')--(B')--(C')--cycle;
    \node[font=\scriptsize,ctTeal] at (0.6,0.3){$T'$};
  \end{scope}
  \draw[arrow] (1.3,0.5)--(2.4,0.5) node[midway,above,font=\scriptsize]{translate};
  % Rotated
  \begin{scope}[xshift=5.5cm, rotate=35]
    \coordinate (A'') at (0,0); \coordinate (B'') at (1.2,0); \coordinate (C'') at (0.6,1);
    \filldraw[ctLightGold,draw=ctGold] (A'')--(B'')--(C'')--cycle;
  \end{scope}
  \draw[arrow] (3.9,0.5)--(5.2,0.5) node[midway,above,font=\scriptsize]{rotate};
\end{tikzpicture}
\end{center}

\begin{principle}
Every isometry preserves shape and size. It is a constraint-preserving
transformation: the distance relation $d(P,Q)=c$ becomes $d(T(P),T(Q))=c$.
\end{principle}

\begin{history}[Hannah Arendt and Invariance]
Arendt's analysis of political action distinguishes between acts that alter
the social space and those that maintain it. An isometry is the spatial analogue
of a norm-preserving act: the shape of relations is maintained even as positions
change. The rotation matrix $R(\theta)$ is the canonical mathematical example.
\end{history}

\begin{exercises}
\begin{exercise}
Show that the composition of two rotations around the origin is another
rotation. What is the resulting angle?
\end{exercise}
\begin{challenge}
A rotation by $\theta$ is represented by the matrix
$R(\theta)=\begin{pmatrix}\cos\theta&-\sin\theta\\\sin\theta&\cos\theta\end{pmatrix}$.
Compute $R(\theta)^\top R(\theta)$ and explain why this confirms that $R(\theta)$
is an isometry.
\end{challenge}
\end{exercises}

\chapter{Area, Volume, and Spatial Distribution}

\section*{Measuring Space}
\addcontentsline{toc}{section}{Opening Perspective}

Up to this point we have examined the structure of geometric figures and
the transformations that preserve them. Points were located by
coordinates, triangles revealed relations among lengths, and
transformations showed how shapes could move through space without
changing their essential structure. Yet geometry also concerns a more
practical question: how much space does a figure occupy?

The measurement of space gives rise to the concepts of area and volume.

Area measures the extent of a figure in two dimensions, while volume
measures the extent of a solid body in three dimensions. These
quantities allow us to compare regions of space, determine capacities,
and analyze the physical constraints imposed by spatial structure.

\medskip

\textbf{From shape to magnitude.}

Earlier chapters emphasized relationships among lengths and angles. A
triangle could be classified by its angles or by the proportional
relationships among its sides. But when we ask how large a triangular
region is, we must introduce a new type of measurement.

Area quantifies the size of a two–dimensional region.

The formulas used to compute area often reflect the geometric structure
of the figure itself. A rectangle has area equal to the product of its
length and width because its interior can be tiled by a grid of equal
squares. A triangle has half the area of a rectangle with the same base
and height because two congruent triangles fill that rectangle exactly.

Thus area formulas arise from simple geometric decompositions.

\medskip

\textbf{Curved shapes and geometric constants.}

Not all geometric figures are composed of straight edges. Circles and
spheres introduce curvature, and their measurement reveals one of the
most famous constants in mathematics: $\pi$.

The area of a circle depends on the square of its radius, while the
surface area of a sphere depends on the square of its radius multiplied
by a constant factor. These relationships illustrate again how scaling
laws operate in geometry. Doubling the radius of a circle quadruples its
area, and doubling the radius of a sphere quadruples its surface area
while increasing its volume by a factor of eight.

Such scaling relationships are direct consequences of the dimensional
structure of space.

\medskip

\textbf{Volume and the geometry of solids.}

While area measures two–dimensional extent, volume measures the capacity
of three–dimensional objects. Many volume formulas can be understood as
extensions of area formulas. For example, the volume of a cylinder is
obtained by multiplying the area of its circular base by its height.
The same principle applies to prisms and many other solids: the volume
is the cross–sectional area multiplied by the length of the object.

This method reflects an important idea in geometry: complex spatial
quantities can often be understood by analyzing simpler lower–dimensional
structures.

\medskip

\textbf{Spatial constraints in natural and social systems.}

The measurement of area and volume is not merely an abstract exercise.
These quantities determine the capacity of physical systems and the
limits within which many processes must operate.

Agricultural productivity depends on the available land area. The amount
of water a reservoir can hold depends on its volume. The intensity of
radiation emitted by a source decreases as it spreads over larger and
larger surfaces.

In ecological and social systems, spatial distribution often imposes
fundamental constraints. A settlement must balance its population
against the productive capacity of the land that supports it. Urban
planning, resource management, and environmental sustainability all
require careful reasoning about how area and volume translate into
available resources.

The speculative societies described in the work of Ursula K.~Le Guin
often illustrate such constraints vividly. Communities organized around
limited ecological resources must constantly consider how population,
land area, and productivity interact.

\medskip

\textbf{Geometry as the mathematics of spatial limits.}

Area and volume formulas therefore serve as algebraic expressions of
spatial constraints. They translate geometric structure into numerical
relationships that can be used to analyze real systems.

When we write $A=\pi r^2$ or $V=\pi r^2 h$, we are not merely describing
shapes; we are expressing the rules governing how space is distributed
within those shapes.

In the sections that follow we examine several common geometric figures
and the formulas used to measure their area and volume. By interpreting
these formulas as constraints relating geometric quantities, we gain a
deeper understanding of how spatial structure determines the limits of
physical and ecological systems.

\section{Area Formulas as Algebraic Constraints}

\begin{center}
\renewcommand{\arraystretch}{1.5}
\begin{tabular}{llp{5cm}}
\toprule
Shape & Formula & Constraint expressed \\
\midrule
Rectangle & $A=lw$           & product of perpendicular dimensions \\
Triangle   & $A=\tfrac12 bh$ & half the enclosing rectangle \\
Circle     & $A=\pi r^2$     & area-radius relation \\
Sphere     & $A=4\pi r^2$   & surface of constant curvature \\
Cylinder   & $V=\pi r^2 h$  & circular cross-section times height \\
\bottomrule
\end{tabular}
\end{center}

\begin{myexample}[Le Guin: sustainable population]
A circular settlement of radius $r=3$~km devotes $60\%$ of its area to crops.
Each km$^2$ supports 400 people.
\[
  A = \pi(3)^2 = 9\pi\;\text{km}^2,\qquad
  A_c = 0.6\times9\pi = 5.4\pi\;\text{km}^2.
\]
Population capacity: $400\times5.4\pi\times100\approx679{,}000$.
(Factor of 100 converts km$^2$ to hectares.)
Le Guin's ecological societies are governed by exactly such spatial constraints
\citep{LeGuin1974}.
\end{myexample}

\begin{exercises}
\begin{practice}
A circle has area $50\pi\;\text{cm}^2$. Find its circumference.
\end{practice}
\begin{exercise}
A cylindrical grain silo has radius 2~m and height 6~m. Find its volume and
total surface area. Interpret each quantity in the context of grain storage.
\end{exercise}
\end{exercises}

% ============================================================
% PART III --- TRIGONOMETRY
% ============================================================
\part{Trigonometry: Cycles and Rotational Systems}

\chapter{Angles and Circular Geometry}

\section*{Rotation and Periodicity}
\addcontentsline{toc}{section}{Opening Perspective}

The preceding chapters explored geometry as the study of spatial
structure. We examined points and distances, the relationships within
triangles, the invariance of shapes under transformations, and the ways
in which area and volume quantify the extent of regions in space.
Throughout these topics one geometric figure appeared repeatedly: the
circle.

The circle occupies a unique place in mathematics because it naturally
connects geometry with motion.

\medskip

\textbf{Rotation as a geometric process.}

A point moving along a straight line changes position without changing
direction. A point moving along a circle, however, continually changes
its direction while maintaining a constant distance from a central
point. This motion is called rotation.

Rotation introduces a new type of measurement: the angle. Angles measure
how far a point has rotated around a center. Unlike distance, which
measures separation between points, angles measure change in direction.

Understanding rotation therefore requires a mathematical language that
describes how positions change as a point moves around a circle.

\medskip

\textbf{The circle as a generator of periodic motion.}

When a point travels around a circle at constant speed, its coordinates
do not change randomly. Instead they follow a repeating pattern.

The horizontal coordinate increases and decreases smoothly as the point
moves from the right side of the circle to the left and back again. The
vertical coordinate follows a similar oscillation as the point moves
from the top of the circle to the bottom and back.

These oscillations repeat every full revolution.

Trigonometry arises from the mathematical description of this periodic
behavior. The trigonometric functions $\sin\theta$ and $\cos\theta$
record the vertical and horizontal coordinates of a point rotating
around the unit circle. As the angle $\theta$ increases, these functions
trace smooth waves that repeat indefinitely.

Thus trigonometry translates circular motion into algebraic functions.

\medskip

\textbf{Angles and the radian measure.}

Earlier chapters introduced the radian as the natural unit for measuring
angles. The radian arises directly from the geometry of circles: it is
defined so that the length of an arc on a circle equals the radius of
the circle when the subtended angle is one radian.

This definition connects angular measurement with linear distance.
Because of this connection, many formulas in trigonometry and calculus
take their simplest form when angles are measured in radians rather than
degrees.

For example, the arc length of a circular segment is given simply by
$s=r\theta$ when $\theta$ is measured in radians.

\medskip

\textbf{Cycles in the natural world.}

The periodic relationships described by trigonometry appear throughout
nature. Planetary orbits, the oscillation of waves, the vibration of
strings, and the rotation of mechanical systems all exhibit cyclic
behavior. Even biological and ecological processes often follow rhythms
that repeat over time.

Astronomers have long relied on such periodic patterns to understand the
structure of the universe. The regular oscillations of certain stars,
known as Cepheid variables, provide one of the most important tools for
measuring cosmic distances. By observing the period of a star's
brightness variation, astronomers can infer its intrinsic luminosity and
therefore estimate how far away it must be.

Cecilia Payne-Gaposchkin contributed significantly to the study of such
stellar phenomena, helping to establish the role of periodic processes
in astrophysical measurement.

\medskip

\textbf{Trigonometry as the mathematics of cycles.}

Trigonometry therefore extends geometry into the study of rotational and
periodic systems. Where earlier chapters analyzed static shapes,
trigonometry analyzes continuous motion along circular paths.

The circle becomes a generator of oscillations, and angles become the
parameter describing position within a cycle. By understanding the
relationships between angles, arc lengths, and coordinates on the
circle, we gain the mathematical tools needed to model waves,
rotations, and repeating phenomena.

\medskip

In the sections that follow we examine the geometry of the circle and
derive the basic relationships connecting angles, arc length, and
coordinates. These ideas will lead naturally to the trigonometric
functions, which describe the oscillatory patterns generated by circular
motion.

Through this development the circle will emerge not merely as a geometric
figure but as a fundamental source of periodic structure in mathematics
and science.

\section{The Circle as Generator of Cycles}

A circle of radius $r$ centred at the origin is the solution set of
$x^2+y^2=r^2$. As a point moves around the circle, its $x$- and
$y$-coordinates each trace out a periodic oscillation. Trigonometry is the
mathematics of this periodicity.

\begin{history}[Cyclic patterns in astrophysics]
Cecilia Payne-Gaposchkin also studied Cepheid variable stars, whose
brightness oscillates with remarkable regularity. The period of oscillation
encodes the star's intrinsic luminosity --- a fact used to measure cosmic
distances. Trigonometric modelling of such cycles was central to her work
\citep{PayneGaposchkin1925}.
\end{history}

\begin{exercises}
\begin{practice}
State the arc length formula $s=r\theta$ (where $\theta$ is in radians).
Find the arc length subtended by an angle of $\pi/3$ on a circle of radius 5~cm.
\end{practice}
\end{exercises}

\chapter{Sine, Cosine, and the Unit Circle}

\section*{Coordinates of Rotation}
\addcontentsline{toc}{section}{Opening Perspective}

In the previous chapter we saw that circular motion naturally produces
periodic patterns. A point moving around a circle generates oscillations
in its horizontal and vertical coordinates. These repeating patterns
form the basis of trigonometry.

To describe this motion precisely, mathematicians introduced two
fundamental functions: sine and cosine.

These functions record the position of a rotating point on a circle.

\medskip

\textbf{The unit circle as a geometric reference system.}

Although circles may have any radius, trigonometry becomes especially
simple when the radius is chosen to be one. Such a circle is called the
unit circle.

When a point moves around the unit circle, its distance from the origin
remains exactly one. The coordinates of the point therefore describe the
direction of the radius rather than its length.

This geometric situation leads directly to the definitions of sine and
cosine. The horizontal coordinate of the rotating point is called the
cosine of the angle, and the vertical coordinate is called the sine.

Thus the trigonometric functions arise naturally from the geometry of
rotation.

\medskip

\textbf{From triangles to circular motion.}

Historically, trigonometric ratios were first defined using right
triangles. In a right triangle with angle $\theta$, the ratio between the
length of the opposite side and the hypotenuse defines $\sin\theta$,
while the ratio between the adjacent side and the hypotenuse defines
$\cos\theta$.

The unit circle interpretation extends these definitions beyond the
limitations of triangles. Instead of restricting $\theta$ to acute
angles, the circular definition allows angles to take any value,
including angles larger than $90^\circ$ and even negative angles.

In this way trigonometry becomes a theory of rotational position rather
than merely a set of triangle ratios.

\medskip

\textbf{Coordinates and vectors.}

The pair $(\cos\theta,\sin\theta)$ can also be interpreted as a vector of
length one pointing in direction $\theta$. Because the vector lies on the
unit circle, its components satisfy a special relationship: the square
of the horizontal component plus the square of the vertical component
equals one.

This relationship is the trigonometric form of the Pythagorean theorem.

\[
\sin^2\theta + \cos^2\theta = 1 .
\]

The identity expresses the geometric fact that the coordinates of any
point on the unit circle must satisfy the equation of that circle.

Thus the most fundamental trigonometric identity emerges directly from
geometry.

\medskip

\textbf{Angles, symmetry, and periodic structure.}

As the angle $\theta$ increases, the point on the unit circle moves
continuously around the circle, and the values of $\sin\theta$ and
$\cos\theta$ vary smoothly between $-1$ and $1$. After a full rotation of
$2\pi$ radians, the point returns to its starting position and the cycle
repeats.

This periodic behavior makes trigonometric functions ideal tools for
describing oscillatory systems. Waves, vibrations, rotating machines,
alternating electrical currents, and many astronomical motions all
follow patterns that repeat with regular periods.

Because sine and cosine describe the geometry of rotation, they provide
a natural mathematical language for analyzing such phenomena.

\medskip

\textbf{Trigonometry and linear transformations.}

Trigonometric functions also appear naturally in the study of geometric
transformations. A rotation of the plane through an angle $\theta$ can be
represented by a matrix whose entries involve $\sin\theta$ and
$\cos\theta$.

This matrix rotates every vector while preserving its length, a property
closely related to the Pythagorean identity. In this way trigonometry
connects circular geometry with the linear algebra of transformations
studied earlier.

\medskip

In the sections that follow we formalize the definitions of sine and
cosine using the unit circle and examine the fundamental identities that
govern their behavior. We also compute the exact values of these
functions for several special angles, which serve as important reference
points for trigonometric reasoning.

Through these ideas the unit circle becomes a geometric framework for
understanding rotation, periodicity, and the oscillatory patterns that
appear throughout mathematics and science.

\section{The Fundamental Definitions}

\begin{mydefinition}[Sine and Cosine]
For a point on the unit circle at angle $\theta$ from the positive $x$-axis,
\[
  \cos\theta = x\text{-coordinate},\qquad
  \sin\theta = y\text{-coordinate}.
\]
Equivalently, in a right triangle with adjacent side $a$, opposite
side $o$, and hypotenuse $h$:
\[
  \cos\theta=\frac{a}{h},\qquad\sin\theta=\frac{o}{h},\qquad\tan\theta=\frac{o}{a}.
\]
\end{mydefinition}

\begin{theorem}[Pythagorean Identity]
For all $\theta$,
\[
  \sin^2\theta+\cos^2\theta=1.
\]
The squared components of a unit vector always sum to one.
\end{theorem}

\begin{proofsketch}
A point $(\cos\theta,\sin\theta)$ lies on the unit circle by definition.
The circle's equation $x^2+y^2=1$ gives the identity directly.
\end{proofsketch}

\section{Exact Values}

\begin{center}
\renewcommand{\arraystretch}{1.55}
\begin{tabular}{ccccc}
\toprule
$\theta$ & radians & $\sin\theta$ & $\cos\theta$ & $\tan\theta$ \\
\midrule
$0°$   & $0$        & $0$            & $1$            & $0$ \\
$30°$  & $\pi/6$    & $1/2$          & $\sqrt{3}/2$   & $1/\sqrt{3}$ \\
$45°$  & $\pi/4$    & $\sqrt{2}/2$   & $\sqrt{2}/2$   & $1$ \\
$60°$  & $\pi/3$    & $\sqrt{3}/2$   & $1/2$          & $\sqrt{3}$ \\
$90°$  & $\pi/2$    & $1$            & $0$            & undef \\
\bottomrule
\end{tabular}
\end{center}

\begin{exercises}
\begin{practice}
Without a calculator, compute $\sin(120°)$ and $\cos(135°)$.
\end{practice}
\begin{exercise}
Prove that $\sin(90°-\theta)=\cos\theta$ using the unit circle definitions.
\end{exercise}
\begin{challenge}
The rotation matrix $R(\theta)$ preserves the norm of every vector. Express
this statement as an identity involving $\sin\theta$ and $\cos\theta$, and
relate it to the Pythagorean identity.
\end{challenge}
\end{exercises}

\chapter{Trigonometric Functions and Graphs}

\section*{Waves, Oscillations, and Repeating Patterns}
\addcontentsline{toc}{section}{Opening Perspective}

In the previous chapter we introduced the sine and cosine functions as
coordinates of a point moving around the unit circle. As the point
rotates, its horizontal and vertical coordinates vary smoothly between
$-1$ and $1$, tracing a repeating pattern. This geometric description
reveals the essential nature of trigonometric functions: they encode
periodic motion.

A periodic function is one whose values repeat after a fixed interval.
Such functions are fundamental in mathematics because they provide a
natural way to describe cycles.

\medskip

\textbf{From circular motion to waves.}

When a point travels around a circle at constant angular speed, its
coordinates oscillate in a regular rhythm. If we plot the vertical
coordinate of that point as a function of time, the resulting curve is a
wave.

This wave is the graph of the sine function.

Similarly, the horizontal coordinate produces the cosine function.
Although these graphs appear different at first glance, they represent
the same underlying motion viewed from different starting positions
along the circle.

In this way circular motion and wave motion become two perspectives on
the same mathematical phenomenon.

\medskip

\textbf{Parameters of oscillation.}

Real oscillatory systems rarely follow the simple form $\sin t$ or
$\cos t$ exactly. Instead they may oscillate with larger or smaller
amplitudes, repeat their cycles more quickly or more slowly, or begin
their motion at different points in the cycle.

To capture these variations, trigonometric functions are often written
in a general form:

\[
f(t) = A \sin(\omega t + \phi) + D .
\]

Each parameter has a clear physical interpretation. The amplitude $A$
determines how large the oscillation is. The angular frequency $\omega$
determines how rapidly the system repeats its cycle. The phase $\phi$
shifts the entire pattern forward or backward in time, and the constant
$D$ raises or lowers the equilibrium level around which the oscillation
occurs.

Together these parameters allow trigonometric functions to model a wide
range of periodic phenomena.

\medskip

\textbf{Oscillations in natural systems.}

Periodic behavior appears throughout the natural world. Ocean tides rise
and fall with predictable rhythms. Sound waves propagate through air as
vibrations of pressure. Mechanical systems such as springs and pendulums
oscillate around equilibrium positions. Even biological systems often
exhibit regular cycles, including daily circadian rhythms and seasonal
patterns of growth.

Astronomy provides especially striking examples of periodic behavior.
Certain stars, known as variable stars, change their brightness in
regular cycles. By analyzing these oscillations, astronomers can infer
properties of the stars and even estimate cosmic distances.

The work of Cecilia Payne-Gaposchkin contributed significantly to the
study of such variable stars. Her research helped establish the
connection between stellar structure and observable brightness
oscillations, demonstrating how periodic patterns reveal information
about distant astrophysical systems.

\medskip

\textbf{Graphs as visual representations of cycles.}

Graphing trigonometric functions provides a powerful way to visualize
periodic behavior. When sine or cosine is plotted against time or angle,
the resulting curve displays the repeating structure of the oscillation.
Peaks correspond to maximum values, troughs correspond to minimum
values, and the midpoint represents the equilibrium level of the
system.

Because these graphs repeat indefinitely, they offer a concise visual
representation of cycles that continue without bound.

\medskip

In the sections that follow we formalize the concept of periodic
functions and examine the graphs of sine and cosine. We then extend
these ideas to general sinusoidal functions that include parameters for
amplitude, frequency, phase, and vertical shift.

By studying these graphs we gain the mathematical tools needed to
analyze oscillatory systems, from the motion of waves to the rhythmic
patterns observed throughout nature and science.

\section{Periodic Functions}

\begin{mydefinition}[Period]
A function $f$ is \textbf{periodic} with period $T>0$ if $f(t+T)=f(t)$ for
all $t$. The sine and cosine functions have period $2\pi$.
\end{mydefinition}

The general sinusoidal function
\[
  f(t)=A\sin(\omega t+\phi)+D
\]
has \textbf{amplitude} $A$, \textbf{angular frequency} $\omega$,
\textbf{phase} $\phi$, and \textbf{vertical shift} $D$. The period is
$T=2\pi/\omega$.

\begin{center}
\begin{tikzpicture}
\begin{axis}[
  width=10cm, height=4cm,
  xlabel={$t$}, ylabel={$f(t)$},
  xmin=0, xmax=4*pi, ymin=-1.6, ymax=1.6,
  xtick={0,pi,2*pi,3*pi,4*pi},
  xticklabels={$0$,$\pi$,$2\pi$,$3\pi$,$4\pi$},
  axis lines=center, grid=both, grid style={ctGray!20},
  tick label style={font=\small},
]
  \addplot[ctBlue,thick,smooth,samples=120,domain=0:4*pi]{sin(deg(x))};
  \addplot[ctTeal,thick,smooth,dashed,samples=120,domain=0:4*pi]{cos(deg(x))};
  \legend{$\sin t$, $\cos t$}
\end{axis}
\end{tikzpicture}
\end{center}

\section{Variable Stars: Trigonometric Modelling}

\begin{myexample}[Payne-Gaposchkin: stellar brightness cycles]
A Cepheid variable star's brightness is modelled by
\[
  B(t)=10+2\sin\!\left(\frac{2\pi}{5}t\right),
\]
with amplitude 2 and equilibrium brightness $=10$.

At $t=2$ days:
\[
  B(2)=10+2\sin\!\left(\frac{4\pi}{5}\right)
      =10+2\sin(144°)\approx10+2(0.588)=11.18.
\]
This model was used by Payne-Gaposchkin to link oscillation period to
distance measurement.
\end{myexample}

\begin{exercises}
\begin{practice}
Identify the amplitude, period, and phase of $f(t)=3\sin(2t-\pi/4)$.
\end{practice}
\begin{exercise}
Gopnik \citep{Gopnik2009} notes that children alternate between exploration
and exploitation in a roughly periodic pattern. Model the exploration
intensity as $E(t)=0.5+0.5\sin(\omega t)$. When is exploration at maximum?
When is it at minimum? What is the period if the cycle is one week?
\end{exercise}
\end{exercises}

\chapter{Modeling Cyclic Phenomena}

\section*{Oscillation as a Universal Pattern}
\addcontentsline{toc}{section}{Opening Perspective}

In the previous chapter we examined trigonometric functions as the
mathematical description of periodic patterns. The sine and cosine
functions arise naturally from the geometry of circular motion, and
their graphs reveal repeating waves that extend indefinitely in time.
Such functions provide a compact language for describing cycles.

Many physical and biological systems exhibit precisely this type of
repeating behavior.

A vibrating string moves back and forth around an equilibrium position.
A pendulum swings periodically under the influence of gravity. Electrical
currents in alternating circuits oscillate with predictable frequency.
Even astronomical systems often display rhythmic patterns, from the
rotation of planets to the pulsation of certain stars.

The mathematics that describes these phenomena is known as harmonic
motion.

\medskip

\textbf{Restoring forces and equilibrium.}

Harmonic motion arises when a system is displaced from an equilibrium
position and experiences a restoring influence that pulls it back toward
that equilibrium. The simplest example occurs in a mass attached to a
spring. If the mass is pulled away from its resting position, the spring
exerts a force proportional to the displacement, pulling the mass back
toward the center.

Because the restoring force depends linearly on displacement, the
resulting motion follows a remarkably regular pattern: the system
oscillates back and forth in a repeating cycle.

Mathematically, this behavior leads to a differential equation whose
solutions are sine and cosine functions.

\medskip

\textbf{Trigonometric functions as solutions of motion.}

The equation governing simple harmonic motion relates the acceleration
of a system to its displacement. When the restoring force is
proportional to the displacement, the acceleration becomes proportional
to the negative of the position. This relationship produces an equation
of the form

\[
\frac{d^2x}{dt^2} + \omega^2 x = 0 .
\]

Although this expression may appear abstract, its solutions describe the
same oscillatory patterns already encountered in trigonometry. The
functions $\sin(\omega t)$ and $\cos(\omega t)$ provide the fundamental
building blocks of these solutions.

Thus trigonometry, which began as the study of circular geometry,
becomes the mathematical language of oscillatory motion.

\medskip

\textbf{Superposition and interference.}

Many real systems involve not just a single oscillation but several
oscillations interacting simultaneously. When two waves overlap, their
effects combine according to the principle of superposition. The
resulting motion is simply the sum of the individual waves.

Depending on their relative phases, waves may reinforce or cancel each
other. When their peaks align, the amplitudes add and produce a larger
wave. When a peak meets a trough, the waves partially or completely
cancel.

This phenomenon, known as interference, plays a central role in the
behavior of sound, light, and other wave-like systems.

\medskip

\textbf{Cycles within cycles.}

Oscillatory behavior often appears in nested or interacting layers.
Planetary systems exhibit orbital cycles within larger galactic
rotations. Biological organisms maintain rhythms within rhythms,
including heartbeats, breathing cycles, and circadian patterns.

Philosopher Alicia Juarrero describes such systems as networks of
constraints in which cycles interact to produce coherent behavior. The
mathematical study of oscillators offers a simplified model of this idea.
By analyzing how periodic processes interact, we gain insight into how
complex systems maintain stability through recurring patterns.

\medskip

\textbf{From geometry to dynamics.}

The development of trigonometry in this part of the book has gradually
shifted attention from static geometric structures to dynamic processes.
Where geometry focused on shapes and spatial relationships, trigonometry
reveals how those relationships evolve through time.

The circle, originally studied as a geometric figure, now becomes the
source of oscillations that describe motion, vibration, and repeating
phenomena.

\medskip

In the sections that follow we introduce the mathematical equation
governing simple harmonic motion and explore how combinations of
trigonometric functions model interacting waves. These ideas will
illustrate how the geometry of the circle provides a powerful framework
for understanding cyclic processes across many areas of science.

\section{Harmonic Motion}

Simple harmonic motion arises whenever a restoring force is proportional to
displacement. The governing equation is
\[
  \frac{d^2x}{dt^2}+\omega^2 x=0,
\]
with general solution $x(t)=A\cos(\omega t)+B\sin(\omega t)$.

\begin{myexample}[Wave interference]
Two waves $f_1=A\sin(\omega t)$ and $f_2=A\sin(\omega t+\delta)$ combine to
\[
  f_1+f_2=2A\cos\!\left(\frac{\delta}{2}\right)\sin\!\left(\omega t+\frac{\delta}{2}\right).
\]
When $\delta=\pi$, amplitude vanishes (destructive interference). The constraint
$\delta=\pi$ is exactly the condition for cancellation.
\end{myexample}

\begin{exercises}
\begin{exercise}
Two sides of a triangle are 7 and 9; the included angle is $40°$. Use the
Law of Cosines to find the third side. Show all unit checks.
\end{exercise}
\begin{challenge}
Juarrero \citep{Juarrero1999} argues that organisms maintain coherent
behaviour through nested cycles of constraint. Model two interacting
oscillators $x(t)=\sin(t)$, $y(t)=\sin(t+\phi)$. For what value of $\phi$
do they cancel? For what value do they reinforce? Discuss the biological
analogy.
\end{challenge}
\end{exercises}

\chapter{Rotations and Harmonic Motion}

\section*{Circular Motion as Hidden Oscillation}
\addcontentsline{toc}{section}{Opening Perspective}

In the previous chapter we studied harmonic motion as a model of
oscillation. A system displaced from equilibrium and subjected to a
restoring force was shown to move in a regular, repeating cycle, and the
solutions to its governing equation were expressed in terms of sine and
cosine. Those functions described the motion algebraically, but a deeper
geometric question remains: why do trigonometric functions appear so
naturally in the study of oscillations?

The answer lies in circular motion.

\medskip

\textbf{A rotating point and its projections.}

Imagine a point moving around a circle at constant angular speed. Its
distance from the center never changes, yet its horizontal and vertical
coordinates continually vary. If we observe only the horizontal motion
of the point, ignoring its vertical position, the point appears to move
back and forth in a smooth oscillation. The same is true for the
vertical component.

In other words, uniform circular motion contains within it two hidden
oscillations: one along the $x$-axis and one along the $y$-axis.

This fact reveals one of the deepest connections in elementary
mathematics. Circular motion and harmonic motion are not separate
phenomena described by coincidentally similar formulas. They are two
different views of the same underlying geometric process.

\medskip

\textbf{The geometry behind sine and cosine.}

When a point rotates around a circle of radius $R$, its position can be
described in terms of angle. The horizontal coordinate is
$R\cos(\omega t)$ and the vertical coordinate is $R\sin(\omega t)$.
These expressions are not arbitrary formulas imposed from outside. They
arise directly from the geometry of the circle and the definitions of
sine and cosine.

Thus trigonometric functions become the natural coordinate language of
rotation.

If we project the rotating point onto one axis, we obtain a sinusoidal
motion. If we project onto the other axis, we obtain another sinusoidal
motion shifted in phase by one quarter of a cycle. Harmonic motion is
therefore the shadow of circular motion cast onto a line.

\medskip

\textbf{Velocity, change, and invariance.}

Circular motion also provides a natural setting for understanding how
position and velocity are related. As a point moves around the circle,
its velocity changes direction continuously even if its speed remains
constant. The motion is therefore governed by both variation and
invariance: the direction changes at every instant, but the distance
from the center and the magnitude of the speed remain fixed.

This combination of change and preservation is central to the broader
theme of the book. A system may transform continuously while still
maintaining invariant structure. In the case of circular motion, the
invariants are radius and speed; the changing quantities are the
coordinates and velocity components.

\medskip

\textbf{From geometric motion to physical systems.}

The relation between rotation and oscillation is not merely a geometric
curiosity. It appears throughout science. Alternating current in an
electrical circuit can be represented by rotating vectors in the complex
plane. Wave motion in physics is often analyzed as the projection of
rotational or phase-based structure. Planetary motion, mechanical
vibration, and signal analysis all depend on the same trigonometric
framework.

By recognizing sinusoidal motion as the projection of circular motion,
we gain a unified picture of periodic systems.

\medskip

\textbf{A bridge between geometry and dynamics.}

This chapter serves as a culmination of the trigonometric part of the
book. Earlier chapters introduced angles, the unit circle, and the
graphs of sine and cosine. Then we studied harmonic motion as a dynamic
system described by those functions. Here the two strands come together:
geometry and dynamics are revealed as different expressions of a single
rotational structure.

The circle is not only a static figure and not only the source of
trigonometric ratios. It is also the hidden geometric engine of
oscillation.

\medskip

In the sections that follow we examine a point moving around a circle at
constant angular velocity and derive the sinusoidal forms of its
coordinate projections. We then analyze the corresponding velocities and
use the Pythagorean identity to show that the speed remains constant.

Through this analysis, circular motion and harmonic motion will appear
not as separate topics but as two mathematically equivalent ways of
describing periodic change.

\section{Circular Motion and Projections}

A point moving at constant angular velocity $\omega$ on a circle of radius
$R$ has position
\[
  (x,y)=(R\cos(\omega t), R\sin(\omega t)).
\]
The $x$- and $y$-projections are independent sinusoidal oscillations, which is
why circular motion and wave motion share the same mathematics.

\begin{center}
\begin{tikzpicture}[scale=1.2]
  \draw[ctGray,thin] (-1.5,0)--(1.5,0);
  \draw[ctGray,thin] (0,-1.5)--(0,1.5);
  \draw[ctBlue!70,thick] (0,0) circle (1);
  \def\t{50}
  \coordinate (P) at ({cos(\t)},{sin(\t)});
  \draw[ctTeal,very thick] (0,0)--(P);
  \draw[ctRed,thick,dashed] ({cos(\t)},0)--(P);
  \draw[ctGold,thick,dashed] (0,0)--({cos(\t)},0);
  \filldraw[ctBlue] (P) circle (1.5pt);
  \node[font=\scriptsize,right] at (P){$P=(\cos\omega t,\sin\omega t)$};
\end{tikzpicture}
\end{center}

\begin{exercises}
\begin{exercise}
Show that if a point traces the unit circle at constant speed, its $x$- and
$y$-velocities are $-\sin(\omega t)$ and $\cos(\omega t)$ respectively.
Verify using the Pythagorean identity that the speed is constant.
\end{exercise}
\end{exercises}

% ============================================================
% PART IV --- STOICHIOMETRY
% ============================================================
\part{Stoichiometry: Conservation and Transformation in Chemistry}

\chapter{Atoms, Molecules, and Chemical Formulas}

\section*{Matter as Conserved Structure}
\addcontentsline{toc}{section}{Opening Perspective}

The earlier parts of this book developed a sequence of mathematical ideas
that at first may have seemed abstract. Algebra introduced the language
of variables and constraints. Geometry described the spatial relations
among points and figures. Trigonometry revealed the cyclic patterns
arising from rotation and oscillation. Despite their different forms,
all of these topics shared a common theme: systems are governed by
relationships that remain invariant under transformation.

Chemistry provides a striking physical realization of this principle.

\medskip

\textbf{Transformation without loss.}

Chemical reactions often appear dramatic. Substances change color, gases
form, heat is released, and entirely new materials emerge. Yet beneath
these visible transformations lies a simple rule: the atoms themselves
are not destroyed. They are merely rearranged.

This principle is known as the law of conservation of mass.

If a chemical reaction begins with a certain number of atoms of each
element, the same number must appear in the products. The atoms change
partners, forming new molecules, but the total count of each element
remains constant.

In this sense a chemical reaction resembles the balance equations
studied earlier in algebra. The quantities involved may change form, but
the underlying totals must remain consistent.

\medskip

\textbf{Atoms as discrete units of matter.}

The atomic hypothesis provides the conceptual framework for this
conservation principle. According to this idea, all matter is composed of
discrete particles called atoms. Atoms of different elements combine in
specific numerical ratios to form molecules.

For example, a molecule of water contains two hydrogen atoms and one
oxygen atom. Carbon dioxide contains one carbon atom and two oxygen
atoms. These numerical patterns define the chemical identity of each
substance.

Chemical formulas therefore function as symbolic descriptions of atomic
structure.

\medskip

\textbf{Counting atoms with the mole.}

Because atoms are extraordinarily small, chemists rarely count them
individually. Instead they measure quantities of substances using the
mole, a unit that represents a fixed number of particles. One mole
contains approximately $6.022\times10^{23}$ entities, a quantity known
as Avogadro's number.

The mole serves as a bridge between the microscopic world of atoms and
the macroscopic quantities measured in the laboratory. By relating mass
to the number of particles present, chemists can determine how much of a
substance participates in a reaction.

This relationship is expressed through the concept of molar mass.

\medskip

\textbf{Chemical formulas as algebraic relations.}

Just as equations in algebra describe relationships between variables,
chemical formulas describe relationships between atoms within a
molecule. The formula $\mathrm{H_2O}$ states that two hydrogen atoms are
combined with one oxygen atom. The formula $\mathrm{CO_2}$ expresses a
ratio of one carbon atom to two oxygen atoms.

When reactions occur, these ratios must be preserved across the entire
system.

Balancing a chemical equation therefore resembles solving a system of
algebraic constraints. Each element imposes its own conservation rule,
and the final reaction must satisfy all of them simultaneously.

\medskip

\textbf{Chemical composition in the cosmos.}

The importance of stoichiometric reasoning extends far beyond laboratory
chemistry. Astronomers use similar principles to interpret the spectra
of stars. By analyzing the patterns of light absorbed and emitted by
stellar atmospheres, scientists can determine the relative abundances of
elements present.

In her pioneering work, Cecilia Payne-Gaposchkin demonstrated that the
Sun and other stars consist primarily of hydrogen and helium. Her
analysis relied on comparing observed spectral intensities with
theoretical predictions, effectively performing a kind of cosmic
stoichiometry.

The proportions of elements in a star determine its structure,
temperature, and energy production.

\medskip

\textbf{Conservation as a unifying theme.}

The transition from mathematics to chemistry in this part of the book is
therefore not abrupt but natural. The balance principles introduced in
algebra, the invariants studied in geometry, and the periodic structures
described by trigonometry all reflect the broader idea that certain
quantities remain constant even as systems transform.

Chemical reactions provide a concrete example of this principle. Matter
changes form, molecules rearrange, and energy flows through the system,
but the total number of atoms of each element remains unchanged.

Understanding this conservation principle is the first step in learning
how to analyze chemical reactions quantitatively.

\medskip

In the sections that follow we examine the atomic hypothesis in more
detail and introduce the mole as the standard unit for measuring the
amount of a substance. We then learn how chemical formulas encode
stoichiometric relationships among atoms and how these relationships
allow us to compute the quantities involved in chemical processes.

Through this approach, chemistry becomes another domain in which
mathematical reasoning reveals the hidden structure of transformation.

\section{The Atomic Hypothesis and Conservation}

In ordinary chemical reactions, atoms are neither created nor destroyed ---
they are rearranged. This is the \textbf{law of conservation of mass}, the
chemical analogue of every balance principle studied so far.

\begin{mydefinition}[Mole and Molar Mass]
The \textbf{mole} (mol) is the SI unit of chemical amount: one mole contains
$N_A=6.022\times10^{23}$ particles (Avogadro's number). The
\textbf{molar mass} $M$ (g/mol) is the mass of one mole of a substance.
The relationship
\[
  n = \frac{m}{M}
\]
converts between mass and amount of substance.
\end{mydefinition}

\begin{history}[Payne-Gaposchkin: hydrogen and helium in stars]
Payne-Gaposchkin's 1925 discovery that stars are overwhelmingly hydrogen and
helium was resisted because it contradicted the then-prevailing view. Her
argument was essentially stoichiometric: matching observed spectral absorption
strengths to theoretical abundances. The molar proportions of H to He in a
star constrain its entire energy budget \citep{PayneGaposchkin1925}.
\end{history}

\begin{myexample}[Moles of water]
How many moles are present in 36~g of $\mathrm{H_2O}$?
\[
  M(\mathrm{H_2O})=2(1.008)+16.00=18.016\gpm.
  \qquad
  n=\frac{36}{18.016}\approx2.0\mol.
\]
\end{myexample}

\begin{exercises}
\begin{practice}
Find the molar mass of (a)~NaCl; (b)~CO$_2$; (c)~C$_6$H$_{12}$O$_6$.
\end{practice}
\begin{exercise}
How many molecules are in 44~g of CO$_2$? Express in scientific notation.
\end{exercise}
\begin{challenge}
Payne-Gaposchkin computed that the Sun's atmosphere is approximately 92\%
hydrogen by number of atoms. If the Sun's atmosphere contains $N$ total atoms,
write an expression for the number of hydrogen atoms and the number of
helium atoms (assuming the remaining 8\% is helium). What mole ratio does
this represent?
\end{challenge}
\end{exercises}

\chapter{Balancing Chemical Equations}

\section*{Chemical Reactions as Constraint Systems}
\addcontentsline{toc}{section}{Opening Perspective}

The previous chapter introduced the atomic hypothesis and the mole as
the fundamental unit for measuring chemical quantities. Molecules were
described as combinations of atoms arranged in fixed numerical ratios,
and chemical formulas were shown to encode these structural patterns.

When a chemical reaction occurs, those molecular structures change.
Atoms break old bonds and form new ones, producing different substances.
Yet the atoms themselves remain present throughout the process. No atom
disappears, and none is created out of nothing.

This requirement leads directly to one of the central quantitative tasks
in chemistry: balancing chemical equations.

\medskip

\textbf{Reactions as rearrangements of atoms.}

A chemical equation represents a transformation of matter. The
substances written on the left-hand side of the equation are called
reactants, and those on the right-hand side are called products. The
equation asserts that the reactants rearrange into the products through
a specific pattern of molecular changes.

However, the equation must also satisfy a conservation constraint.

Because atoms are conserved during ordinary chemical reactions, the
total number of atoms of each element must be the same before and after
the reaction. If a reaction begins with two carbon atoms, six hydrogen
atoms, and several oxygen atoms, the products must contain exactly the
same numbers of those atoms.

Balancing a chemical equation therefore ensures that these conservation
conditions are satisfied.

\medskip

\textbf{From counting atoms to solving equations.}

At first glance balancing reactions might appear to be a matter of trial
and error: adjusting coefficients until the numbers of atoms match on
both sides. While this approach can work for simple reactions, it
quickly becomes cumbersome when many elements or molecules are involved.

A more systematic approach emerges when the problem is interpreted
mathematically.

Each chemical element imposes a conservation rule. The number of atoms
of that element on the reactant side must equal the number on the
product side. If a reaction involves several elements, each one
contributes an equation expressing its conservation.

The coefficients multiplying each chemical species then become unknown
variables in a system of equations.

Balancing the reaction is equivalent to solving that system.

\medskip

\textbf{Stoichiometric matrices and linear algebra.}

Earlier chapters introduced matrices as tools for representing systems
of linear relationships. The same idea can be applied to chemical
reactions. A matrix can be constructed whose rows correspond to
elements and whose columns correspond to chemical species. The entries
of the matrix record how many atoms of each element appear in each
molecule.

When this matrix multiplies the vector of stoichiometric coefficients,
the result must be zero if the reaction is balanced. This condition
expresses the conservation of every element simultaneously.

Thus the balancing of chemical equations becomes a problem in linear
algebra.

The mathematical techniques developed earlier in the book—systems of
equations, matrices, and constraint networks—now find a direct
application in chemical reasoning.

\medskip

\textbf{Stoichiometry as quantitative chemistry.}

The systematic study of these numerical relationships is called
stoichiometry. Once the coefficients of a reaction are determined, they
reveal how much of each substance participates in the reaction. The
balanced equation tells us how many molecules or moles of reactants are
required to produce a given quantity of products.

In industrial chemistry and laboratory practice, these ratios determine
how reactions are carried out efficiently. They guide the calculation of
reactant quantities, product yields, and energy requirements.

\medskip

\textbf{A bridge between mathematics and chemistry.}

Balancing chemical equations therefore illustrates a deeper theme of
this book: the same mathematical structures appear across seemingly
different domains. The conservation rules of chemistry correspond to
constraint systems in algebra. The stoichiometric matrix mirrors the
matrix representations studied earlier in linear transformations. The
solutions of these systems describe how matter transforms while
preserving its fundamental components.

Chemistry, in this sense, becomes another arena in which mathematical
structure governs physical processes.

\medskip

In the sections that follow we formalize the conservation principles
underlying chemical reactions and express them using systems of linear
equations. We then develop practical methods for balancing reactions and
interpreting the resulting coefficients.

Through this approach the balancing of chemical equations emerges not as
a mechanical procedure but as a clear example of conservation and
constraint operating within a physical system.

\section{Conservation as a Linear System}

\begin{theorem}[Stoichiometric Conservation]
A chemical equation is balanced if and only if, for each element, the total
atom count on the left side equals the total atom count on the right side.
Equivalently, the stoichiometric coefficient vector $\gamma$ satisfies
$N\gamma=\mathbf{0}$, where $N$ is the element--species incidence matrix.
\end{theorem}

\begin{myexample}[Combustion of ethane]
Balance $\mathrm{C_2H_6+O_2\to CO_2+H_2O}$.

Assign coefficients $a,b,c,d$ and write conservation equations:
\begin{align*}
  \text{C}: & \quad 2a=c \\
  \text{H}: & \quad 6a=2d \\
  \text{O}: & \quad 2b=2c+d
\end{align*}
Set $a=1$: $c=2$, $d=3$, $b=7/2$. Multiply by 2:
\[
  2\mathrm{C_2H_6}+7\mathrm{O_2}\to4\mathrm{CO_2}+6\mathrm{H_2O}.
\]
\end{myexample}

\begin{center}
\begin{tikzpicture}[node distance=1.4cm, scale=0.9]
  % Reaction network
  \node[qty,fill=ctLightGold] (CH4) {CH$_4$};
  \node[qty,fill=ctLightBlue,right=2.2cm of CH4] (CO2) {CO$_2$};
  \node[qty,fill=ctLightTeal,below=of CH4] (O2) {O$_2$};
  \node[qty,fill=ctLightTeal,below=of CO2] (H2O) {H$_2$O};
  \draw[arrow] (CH4) -- (CO2) node[midway,above,font=\scriptsize]{reaction};
  \draw[arrow] (O2) -- (CO2);
  \draw[arrow] (O2) -- (H2O);
  \draw[arrow] (CH4) -- (H2O);
  \node[font=\scriptsize,ctGray,below=1.8cm of O2]{$\mathrm{CH_4+2O_2\to CO_2+2H_2O}$};
\end{tikzpicture}
\end{center}

\begin{exercises}
\begin{practice}
Balance (a)~$\mathrm{Fe+O_2\to Fe_2O_3}$;\;
(b)~$\mathrm{C_3H_8+O_2\to CO_2+H_2O}$.
\end{practice}
\begin{exercise}
Write the stoichiometric matrix $N$ for the reaction
$\mathrm{CH_4+2O_2\to CO_2+2H_2O}$ and verify $N\gamma=\mathbf{0}$.
\end{exercise}
\begin{challenge}
Balance $\mathrm{KMnO_4+HCl\to KCl+MnCl_2+H_2O+Cl_2}$ by writing and
solving the linear system. This is a non-trivial redox reaction with 7
unknowns.
\end{challenge}
\end{exercises}

\chapter{Moles and Quantitative Chemistry}

\section*{Counting Matter in Chemical Reactions}
\addcontentsline{toc}{section}{Opening Perspective}

The previous chapter showed that balancing a chemical equation is
equivalent to satisfying conservation constraints for every element
involved in a reaction. Once a reaction is balanced, the coefficients
in the equation express the ratios in which molecules participate in
the transformation. These coefficients therefore encode the structural
logic of the reaction.

To apply this information in practice, however, we must connect the
symbolic language of chemical equations with measurable quantities of
matter.

This connection is provided by the mole.

\medskip

\textbf{From molecules to measurable quantities.}

Chemical equations describe reactions in terms of molecules or atoms.
For example, the equation

\[
2\mathrm{H_2}+ \mathrm{O_2} \rightarrow 2\mathrm{H_2O}
\]

states that two molecules of hydrogen combine with one molecule of
oxygen to produce two molecules of water. While this description is
precise at the microscopic level, real chemical experiments involve
enormous numbers of molecules. Even a tiny droplet of water contains
trillions upon trillions of particles.

Because such numbers are impractical to count individually, chemists
use the mole as a standard counting unit.

\medskip

\textbf{The mole as a bridge between scales.}

One mole represents a fixed number of particles, known as Avogadro's
number:

\[
N_A = 6.022 \times 10^{23}.
\]

This constant plays a role similar to the idea of a dozen in everyday
counting. Just as a dozen always means twelve objects, one mole always
means $6.022\times10^{23}$ entities, whether those entities are atoms,
molecules, or ions.

The mole therefore links the microscopic world of particles with the
macroscopic measurements of mass used in the laboratory.

Molar mass provides the conversion between these two scales. Because
each substance has a characteristic molar mass, dividing the measured
mass of a sample by its molar mass reveals the number of moles present.

\medskip

\textbf{Stoichiometric ratios as quantitative rules.}

Once quantities are expressed in moles, the coefficients of a balanced
chemical equation can be interpreted directly as ratios between amounts
of substances. In the reaction

\[
2\mathrm{H_2}+ \mathrm{O_2} \rightarrow 2\mathrm{H_2O},
\]

two moles of hydrogen react with one mole of oxygen to produce two
moles of water. If four moles of hydrogen react, two moles of oxygen
are required and four moles of water will be formed.

The balanced equation therefore functions as a quantitative rule
governing the transformation of matter.

Stoichiometric calculations simply apply these ratios in combination
with the mole–mass conversion.

\medskip

\textbf{Conservation revisited.}

The ability to compute quantities in chemical reactions reflects the
same conservation principles that have appeared throughout this book.
Earlier chapters described balance conditions in algebra, invariants
under geometric transformations, and conserved relationships in cyclic
systems. In chemistry the invariant is the number of atoms of each
element.

Stoichiometry translates this conservation law into numerical
relationships that allow chemists to predict the outcomes of reactions.

\medskip

\textbf{Quantitative reasoning in chemistry.}

These calculations are not merely classroom exercises. In laboratory
practice and industrial chemistry, stoichiometric reasoning determines
how much reactant must be supplied to produce a desired quantity of
product. It also reveals how efficiently a reaction proceeds and how
much material might remain unused.

Whether producing pharmaceuticals, refining fuels, or analyzing
environmental samples, chemists rely on these quantitative relationships
to control chemical processes.

\medskip

In the sections that follow we develop systematic methods for converting
between mass, moles, and molecular quantities. Using balanced chemical
equations, we then compute the amounts of substances produced or
consumed in reactions.

Through these calculations the abstract conservation principles of
chemistry become practical tools for measuring and predicting the
transformation of matter.

\section{Stoichiometric Calculations}

\begin{myexample}[Mass of water produced]
From $2\mathrm{H_2+O_2\to2H_2O}$, how many grams of water form from 4~g of
$\mathrm{H_2}$ (excess $\mathrm{O_2}$)?
\begin{align*}
  n(\mathrm{H_2}) &= \frac{4\;\text{g}}{2.016\gpm} \approx1.984\mol\\
  n(\mathrm{H_2O}) &= 1.984\mol\times\frac{2\;\text{mol H}_2\text{O}}{2\;\text{mol H}_2}
    = 1.984\mol\\
  m(\mathrm{H_2O}) &= 1.984\times18.016 \approx 35.7\;\text{g}
\end{align*}
\end{myexample}

\begin{exercises}
\begin{practice}
From $\mathrm{C_3H_8+5O_2\to3CO_2+4H_2O}$, how many grams of CO$_2$ are
produced by burning 11~g of propane ($M=44.1\gpm$)?
\end{practice}
\end{exercises}

\chapter{Reaction Ratios and Limiting Reagents}

\section*{Scarcity and Constraint in Chemical Systems}
\addcontentsline{toc}{section}{Opening Perspective}

In the previous chapter we learned how balanced chemical equations
translate the conservation of atoms into quantitative relationships
between substances. By converting masses to moles and applying the
stoichiometric ratios encoded in a reaction, we can determine how much
product a reaction could theoretically produce.

However, real chemical reactions rarely proceed under perfectly balanced
conditions. In practice, the reactants present at the beginning of a
reaction are not always available in the exact ratios required by the
balanced equation.

This imbalance introduces a new constraint.

\medskip

\textbf{Reactions as resource transformations.}

A chemical reaction may be viewed as a transformation of resources. Each
reactant provides a supply of atoms that can participate in the
formation of products. For the reaction to proceed, the required atoms
must be available in the proportions specified by the balanced
equation.

If one reactant runs out before the others, the reaction can no longer
continue, even if other substances remain present.

The reactant that is consumed first therefore determines how far the
reaction can proceed.

\medskip

\textbf{The limiting reagent principle.}

The reactant that restricts the extent of the reaction is called the
limiting reagent. Because it is used up first, it sets the maximum
possible quantity of product that can form. All other reactants are
present in excess and remain partially unused.

This idea is analogous to many everyday situations involving limited
resources. If a baker has flour, sugar, and eggs but only enough eggs to
make ten cakes, then the eggs determine how many cakes can be baked,
regardless of how much flour or sugar remains.

In chemistry the same logic applies. The scarcest required component
propagates its limitation through the entire reaction.

\medskip

\textbf{Constraint propagation in reaction networks.}

When a reaction involves multiple substances, each reactant imposes a
potential constraint on the system. The limiting reagent is the one that
becomes active first, restricting the overall outcome.

Philosopher Alicia Juarrero has described similar processes in complex
systems as cascades of constraints. When one constraint becomes active,
it influences the behavior of the entire system, shaping what outcomes
are possible.

Chemical reactions provide a clear quantitative example of this idea.
The availability of each reactant determines how far the reaction can
proceed and how much product can be formed.

\medskip

\textbf{Theoretical yield and real outcomes.}

Once the limiting reagent is identified, it determines the theoretical
yield of the reaction: the maximum amount of product that could form if
the reaction proceeded perfectly according to stoichiometric ratios.

In real laboratory conditions, however, reactions rarely achieve this
ideal. Side reactions, incomplete conversion, and practical losses
during purification often reduce the amount of product actually
obtained.

To measure this difference, chemists compare the actual yield obtained
in an experiment with the theoretical yield predicted by stoichiometric
calculations.

The ratio between these two quantities, expressed as a percentage, is
called the percent yield.

\medskip

\textbf{Stoichiometry as system analysis.}

The concepts of limiting reagents and percent yield reveal how chemical
reactions can be analyzed as constrained systems. Each reactant
contributes a resource, and the reaction proceeds only as far as the
most restrictive resource allows.

By identifying which component limits the system, chemists can predict
the maximum output of a reaction and design experiments that use
materials efficiently.

\medskip

In the sections that follow we formalize the definition of the limiting
reagent and develop methods for identifying it in chemical reactions.
We then introduce the concept of percent yield and show how it compares
theoretical predictions with experimental outcomes.

Through these ideas the stoichiometric analysis of reactions becomes a
study of how conservation laws and resource constraints determine the
possible transformations of matter.

\section{The Limiting Reagent}

\begin{mydefinition}[Limiting Reagent]
The \textbf{limiting reagent} is the reactant consumed first, setting the
maximum possible yield. All other reactants are \textbf{excess reagents}.
\end{mydefinition}

\begin{myexample}[Juarrero: constraint propagation in a reaction]
Consider $\mathrm{N_2+3H_2\to2NH_3}$. Begin with 2~mol $\mathrm{N_2}$ and
5~mol $\mathrm{H_2}$.

$\mathrm{N_2}$ alone could produce $4\;\text{mol NH}_3$; $\mathrm{H_2}$ alone
could produce $5\times\tfrac{2}{3}\approx3.33\;\text{mol NH}_3$.

$\mathrm{H_2}$ is limiting. Yield $=3.33\;\text{mol NH}_3$.

Juarrero \citep{Juarrero1999} calls this \textit{constraint cascade}: the
scarcest resource propagates its limitation through the entire network.
\end{myexample}

\section{Percent Yield}

\[
  \%\text{yield}
  = \frac{\text{actual yield}}{\text{theoretical yield}}\times100\%.
\]

\begin{exercises}
\begin{exercise}
How many moles of excess $\mathrm{N_2}$ remain after the reaction in the
example above?
\end{exercise}
\begin{challenge}
Model a three-step reaction cascade $A+B\to C$, $C+D\to E$, $E+F\to G$.
Starting with 4~mol each of $A,B,D,F$ and none of the intermediates,
determine the theoretical yield of $G$. Identify which reagent is limiting
at each stage.
\end{challenge}
\end{exercises}

\chapter{Reaction Networks and Conservation Laws}

\section*{Networks of Transformation}
\addcontentsline{toc}{section}{Opening Perspective}

The previous chapters introduced stoichiometry as the quantitative
analysis of chemical reactions. Individual reactions were balanced by
enforcing conservation of atoms, and stoichiometric ratios allowed us to
compute how much product could be formed from given reactants. These
methods treated reactions largely as isolated events.

In reality, however, most chemical processes do not occur in isolation.
They form networks of interacting transformations.

\medskip

\textbf{From single reactions to reaction networks.}

In many natural and industrial systems, the products of one reaction
become the reactants of another. A sequence of reactions may convert raw
materials into intermediate compounds, which are then transformed into
final products. Biological metabolism, atmospheric chemistry, and
industrial synthesis all involve such chains of reactions.

When many reactions occur simultaneously, the system becomes a reaction
network. Each reaction alters the concentrations of multiple chemical
species, and these changes propagate through the network as the system
evolves over time.

The analysis of such systems requires a mathematical framework capable
of describing how many interacting transformations occur at once.

\medskip

\textbf{Stoichiometric structure of networks.}

The key mathematical object in the study of reaction networks is the
stoichiometric matrix. Each column of this matrix corresponds to a
reaction, and each row corresponds to a chemical species. The entries of
the matrix record how many molecules of each species are consumed or
produced by each reaction.

In this representation, the entire network of chemical transformations
is encoded within a single matrix.

When this matrix acts on a vector of reaction rates, it determines how
the concentrations of all species change over time. The resulting
equations describe the dynamic evolution of the chemical system.

\medskip

\textbf{Conservation laws and invariant quantities.}

Even within complex reaction networks, certain quantities remain
constant. Because atoms are conserved, combinations of species that
contain the same atoms must satisfy conservation relations. These
relations correspond mathematically to vectors in the kernel of the
stoichiometric matrix.

If a vector lies in this kernel, it represents a quantity that does not
change as reactions proceed.

Thus conservation laws appear as invariants within the mathematical
structure of the reaction network.

\medskip

\textbf{Flows of matter and mathematical continuity.}

The equations governing reaction networks resemble other conservation
equations encountered in physics. For example, fluid dynamics describes
the motion of matter using continuity equations that express the
principle that matter cannot disappear or appear spontaneously. The rate
of change of density within a region must equal the net flow of material
into or out of that region.

Although the mathematical forms differ, the underlying principle is the
same. Whether we are studying molecules reacting in a chemical system or
mass flowing through a fluid, the equations ensure that conserved
quantities remain balanced.

\medskip

\textbf{From discrete reactions to continuous fields.}

Reaction networks therefore provide a bridge between discrete chemical
processes and continuous physical systems. The same conservation
principles appear in algebraic stoichiometry, in differential equations
describing reaction kinetics, and in field theories describing the flow
of matter or energy through space.

Seen from this perspective, chemistry becomes one example of a broader
mathematical pattern: systems evolve through transformations that
preserve certain underlying quantities.

\medskip

In the sections that follow we formalize the mathematical description of
reaction networks using the stoichiometric matrix and reaction rate
functions. These tools allow us to describe how concentrations change
over time and how conservation laws emerge naturally from the structure
of the network.

Through this framework we will see that the same balance principles
encountered throughout this book—algebraic constraints, geometric
invariants, and periodic structures—also govern the dynamic evolution of
chemical systems.

\section{Dynamic Stoichiometry}

In a reaction network, concentrations $\mathbf{c}$ evolve by
\[
  \frac{d\mathbf{c}}{dt} = S\,\mathbf{r}(\mathbf{c}),
\]
where $S$ is the stoichiometric matrix and $\mathbf{r}(\mathbf{c})$ is the
vector of reaction rates. Conservation is encoded in $\ker(S)$.

\begin{myexample}[Continuity in disguise]
The equation above is the discrete analogue of the field-theoretic continuity
equation
\[
  \pd{\rho}{t} + \nabla\cdot(\rho\mathbf{v}) = 0,
\]
describing conservation of density $\rho$ in a flowing medium with velocity
$\mathbf{v}$. Algebra, stoichiometry, and field theory all encode the same
structural idea: \emph{what flows in must flow out}.
\end{myexample}

\begin{deepnote}
Appendix~E derives the RSVP field equations as the continuum limit of this
conservation structure.
\end{deepnote}

% ============================================================
% PART V --- INTEGRATION
% ============================================================
\part{Integration: Mathematics of Structured Systems}

\chapter{Constraint Networks}

\section*{Systems of Relations}
\addcontentsline{toc}{section}{Opening Perspective}

The preceding parts of this book developed several distinct mathematical
languages. Algebra described relationships among quantities through
equations. Geometry revealed invariants under spatial transformation.
Trigonometry analyzed cyclic behavior generated by rotation. Chemistry
introduced conservation laws governing the transformation of matter.
Although these topics appear to belong to different domains, they share
a common structural theme.

Each system can be understood as a collection of interconnected
constraints.

\medskip

\textbf{From isolated equations to systems of relations.}

In early chapters we studied individual equations and small systems of
equations. Each equation imposed a restriction on the possible values of
variables. When multiple equations were combined, the set of allowable
solutions became progressively narrower until a consistent solution
emerged.

This logic extends naturally to larger systems.

In many scientific and social contexts, relationships do not occur in
pairs but within complex webs of interaction. Variables influence one
another through chains of dependencies. A change in one component
propagates through the entire structure of relations.

To represent such systems clearly, mathematics introduces the concept of
a network.

\medskip

\textbf{Nodes, edges, and structure.}

A network, or graph, consists of nodes connected by edges. The nodes
represent entities within a system, while the edges represent
relationships or constraints linking those entities.

Depending on the context, nodes might represent physical quantities,
chemical species, individuals in a social system, or locations in an
ecological landscape. Edges may represent equations, flows of material,
communication links, or causal influences.

By representing systems as networks, we focus attention not only on the
individual components but on the structure of their interactions.

\medskip

\textbf{Propagation of constraints.}

One of the most important features of a constraint network is that
changes propagate through the structure. When the value of one node
changes, the constraints connecting it to other nodes require those
nodes to adjust as well. This chain reaction continues until the system
reaches a new consistent configuration.

Earlier chapters provided several examples of this phenomenon. In a
system of linear equations, changing one variable affects the values of
others. In a chemical reaction network, altering the concentration of
one species influences the entire set of reaction rates. In ecological
models, changing the availability of a resource affects every process
that depends on it.

Constraint propagation therefore provides a unified perspective on many
types of systems.

\medskip

\textbf{Networks across disciplines.}

The language of networks appears in many areas of science and social
analysis. Physicists describe interacting particles as networks of
forces. Biologists model metabolic pathways as networks of chemical
reactions. Computer scientists analyze information flow through networks
of processors or communication channels.

Philosopher Hannah Arendt described political life as a web of
relationships among individuals whose actions and communications create
a shared public space. Similarly, speculative social systems imagined in
the work of Ursula K.~Le Guin often revolve around carefully balanced
networks of ecological and social constraints.

Although the contexts differ, the underlying mathematical idea remains
the same: complex systems are defined not only by their components but
by the relationships connecting those components.

\medskip

\textbf{Networks as mathematical objects.}

Modern mathematics studies networks using the theory of graphs and
related structures. These tools allow us to analyze connectivity,
stability, and flow within complex systems. Techniques from linear
algebra, probability, and dynamical systems all contribute to the study
of networks.

By representing a system as a network of nodes and edges, we obtain a
framework capable of describing interactions that would be difficult to
analyze in isolation.

\medskip

In the sections that follow we formalize the concept of a network and
explore how systems of equations, ecological relationships, and social
interactions can all be represented as constraint networks. Through this
approach we see that many of the mathematical ideas developed throughout
this book are different manifestations of a single underlying
principle: structure emerges from the relationships that link the parts
of a system together.

\section{Networks as Mathematical Objects}

A \textbf{network} (or graph) consists of \textbf{nodes} (vertices) and
\textbf{edges} (links). When nodes represent quantities and edges represent
constraints, the network encodes the structure of a system.

\begin{myexample}[Social constraint network after Arendt]
Arendt's analysis of political space can be formalised as a constraint network
in which nodes are actors and directed edges are communicative acts
\citep{Arendt1958}. The system is stable when the network satisfies certain
consistency conditions --- analogous to the solvability of a linear system.
\end{myexample}

\begin{myexample}[Ecological resource balance after Le Guin]
The allocation system $F=2L$, $E=F+L$ from Chapter~4 is a two-node constraint
network. Le Guin's utopian societies must maintain exactly such networks at
equilibrium \citep{LeGuin1974}.
\end{myexample}

\chapter{Cycles, Flows, and Equilibrium}

\section*{Balance in Mathematical Systems}
\addcontentsline{toc}{section}{Opening Perspective}

Throughout this book we have encountered many different mathematical
structures: algebraic equations, geometric relations, trigonometric
identities, chemical reaction balances, and differential equations
describing physical flows. At first glance these subjects appear to
belong to separate domains.

Yet each of them expresses a common structural idea: balance.

In algebra, balance appears as the equality of two expressions. In
geometry it appears as the invariant relationship between the sides of a
right triangle. In trigonometry it appears as the identity linking sine
and cosine on the unit circle. In chemistry it appears as the
conservation of atoms in a reaction. In physics it appears as the
continuity equation governing the flow of matter.

Although the mathematical symbols differ, the underlying principle is
the same. A conserved quantity constrains how a system can change.

\medskip

\textbf{Balance as a structural law.}

A balanced relation states that certain quantities cannot change
independently. If one part of the relation varies, another part must vary
in a compensating way.

For example, if the sides of a right triangle change, the relationship
$a^2+b^2=c^2$ must still hold. If the concentrations of species in a
chemical reaction evolve, the total number of atoms of each element
remains constant. If mass flows through a region of space, the change in
density must equal the net flow entering or leaving that region.

Such constraints define the allowable states of the system.

\medskip

\textbf{Dynamic systems and equilibrium.}

Many systems evolve over time. Chemical concentrations change, physical
fields move through space, and populations fluctuate as resources are
consumed and replenished. The evolution of these systems is described by
rules that transform one state into another.

Sometimes these dynamics reach a configuration in which further change
ceases. The system has reached equilibrium.

Mathematically, equilibrium corresponds to a fixed point of the system's
transformation rule. Applying the system dynamics to that state produces
no further change.

\medskip

\textbf{Unifying patterns.}

The purpose of this chapter is to place the different mathematical ideas
studied throughout the book within a single conceptual framework. We
will see that algebraic equations, geometric identities, trigonometric
relations, chemical balances, and physical continuity equations all
express the same fundamental idea: systems evolve under constraints that
preserve certain quantities.

Recognizing these shared structures allows us to understand seemingly
different phenomena as manifestations of the same underlying principles
of balance, flow, and equilibrium.

\section{Three Kinds of Balance}

All the subjects in this book converge on three structural ideas:

\begin{center}
\renewcommand{\arraystretch}{1.6}
\begin{tabular}{>{\bfseries}p{3cm} p{4cm} p{5cm}}
\toprule
\color{ctBlue}Type & \color{ctBlue}Domain & \color{ctBlue}Equation \\
\midrule
Algebraic balance  & variable equations     & $ax=b$ \\
Spatial balance    & Pythagorean geometry   & $a^2+b^2=c^2$ \\
Rotational balance & unit circle            & $\sin^2+\cos^2=1$ \\
Material balance   & stoichiometry          & $N\gamma=\mathbf{0}$ \\
Dynamic balance    & field continuity       & $\partial_t\rho+\nabla\cdot(\rho\mathbf{v})=0$ \\
\bottomrule
\end{tabular}
\end{center}

\section{Equilibrium as Fixed Point}

\begin{mydefinition}[Equilibrium]
A system state $\mathbf{x}^*$ is an \textbf{equilibrium} if applying the
system's dynamics to $\mathbf{x}^*$ returns $\mathbf{x}^*$:
$T(\mathbf{x}^*)=\mathbf{x}^*$.
\end{mydefinition}

\chapter{Scaling and Multilevel Structure}

\section*{Structure Across Scales}
\addcontentsline{toc}{section}{Opening Perspective}

One of the most striking features of natural systems is that similar
patterns often appear at very different scales. Structures that govern
microscopic processes frequently reappear in macroscopic phenomena.
Mathematics provides the language that allows us to recognize these
recurring patterns.

In earlier chapters we studied scaling relationships such as
$y=kx^p$. These relations describe how one quantity changes when another
is enlarged or reduced. The exponent $p$ determines how rapidly the
dependent quantity grows relative to the independent one.

Many important physical laws are scaling laws. The surface area of a
sphere grows with the square of its radius, while its volume grows with
the cube of the radius. The intensity of light from a star decreases
with the square of the distance from the source. Such relationships
arise because geometric dimensions change in predictable ways when
objects are enlarged.

\medskip

\textbf{Self-similar patterns.}

When the same mathematical relationship appears at multiple scales, the
system is said to exhibit scale invariance or self-similarity. A pattern
that governs a small system also governs larger versions of the system,
provided that the relevant variables are scaled appropriately.

This property occurs in many scientific contexts. Physical laws remain
valid across enormous ranges of size, from laboratory experiments to
astronomical systems. Ecological networks exhibit similar structural
principles whether they involve microorganisms in a soil community or
entire planetary biospheres.

The recurrence of such patterns suggests that certain structural
relationships are fundamental features of the systems being studied.

\medskip

\textbf{Multilevel systems.}

Many complex systems consist of multiple interacting levels. Chemical
interactions determine the properties of minerals. Geological processes
determine the distribution of those minerals within planetary crusts.
Biological activity can in turn alter mineral diversity by introducing
new chemical environments.

Each level of organization operates according to its own mechanisms, yet
the same constraint logic often governs the emergence of structure.

\medskip

\textbf{Scaling as a unifying idea.}

Understanding how relationships behave under scaling allows us to
connect phenomena that occur at very different magnitudes. The same
mathematical principles that describe laboratory reactions, planetary
processes, or stellar radiation fields may be applied across many orders
of magnitude.

In this chapter we examine how patterns repeat across scales and how
mathematical scaling laws provide a framework for understanding
multilevel structure in complex systems.

\section{Patterns That Recur Across Scales}

\begin{theorem}[Self-Similarity Under Scaling]
If a relationship $y=kx^p$ holds at one scale, and the same exponent $p$
governs the relationship at larger or smaller scales, the pattern is
\textbf{scale-invariant}.
\end{theorem}

The inverse-square law $I=L/(4\pi r^2)$ is scale-invariant: it holds whether
$r$ is one metre or one light-year.

\begin{myexample}[Hazen: mineral evolution]
Hazen et al.\ \citep{Hazen2008} trace how the diversity of mineral species on
Earth grew in proportion to the complexity of planetary processes. At each
scale --- atomic bonding, crystal structure, tectonic activity, biological
metabolism --- the same constraint logic applies: new mineral species form
precisely when new constraints (temperature, pressure, chemical environment)
are satisfied.
\end{myexample}

\chapter{From Elementary Systems to Scientific Models}

\section*{From Simple Relations to Scientific Models}
\addcontentsline{toc}{section}{Opening Perspective}

The mathematical ideas introduced in this book may appear elementary when
considered individually. Ratios, triangles, trigonometric identities, and
balanced chemical equations are often presented as separate topics within
secondary education. Yet these concepts form the essential building blocks
of nearly every quantitative scientific model.

Scientific modeling begins with the identification of relationships between
quantities. These relationships may describe spatial structure, temporal
variation, material transformation, or flows of energy and matter. Once
such relationships are identified, mathematics provides the language
needed to express them precisely.

\medskip

\textbf{From relations to systems.}

Early chapters focused on simple equations and geometric relations. A
single equation constrains the possible values of variables, while a
geometric relation restricts the shape of a figure. When multiple
relations are combined, they form a system whose behavior is determined
by the interaction of its constraints.

More advanced scientific models follow exactly the same pattern.
Physical laws are expressed as systems of equations whose solutions
describe the evolution of a system over time or the configuration of a
system in space. Whether one studies planetary motion, ecological
dynamics, or chemical reactions, the mathematical structure always
consists of interacting constraints.

\medskip

\textbf{Transformation and invariance.}

Two complementary ideas organize scientific reasoning. The first is
transformation: the rules that describe how a system changes. The second
is invariance: the quantities that remain constant despite those
changes.

For example, geometric transformations may rotate or translate a shape
while preserving distances. Chemical reactions rearrange atoms while
preserving elemental counts. Physical fields may flow through space
while conserving mass or energy.

Recognizing invariant quantities often reveals the fundamental laws
governing a system.

\medskip

\textbf{Layered mathematical structures.}

The diagram in this chapter illustrates how the topics studied in this
book connect to one another. Algebra provides the symbolic language used
to express relations. Geometry describes spatial structure. Trigonometry
captures periodic behavior arising from rotation. Stoichiometry
expresses conservation in chemical transformations.

These ideas can be extended to more advanced frameworks. Field theories
describe how quantities vary continuously through space and time, while
modern mathematical structures such as category theory provide abstract
ways of representing relationships between entire systems.

\medskip

\textbf{Scientific models as structured networks.}

At its core, a scientific model is a structured network of mathematical
relations. Each component of the model represents a quantity or process,
and the equations linking those components encode the constraints that
govern their interactions.

Although the models used in advanced science may appear complex, their
architecture rests on the same fundamental ideas encountered throughout
this book: equations, transformations, invariants, and conservation
laws.

\medskip

The goal of this final chapter is to place the mathematical tools we have
studied into a broader scientific perspective. By recognizing the common
structure underlying algebra, geometry, trigonometry, and stoichiometry,
we can see how these elementary systems grow into the sophisticated
models used to understand the natural world.

\section{The Architecture of Science}

The tools introduced in this book form the foundation of all quantitative
science. Every advanced model is an elaboration of the same constraint
structure encountered in Chapter~1.

\begin{center}
\begin{tikzpicture}[node distance=1.2cm and 2.8cm, scale=0.9]
  \node[qty,fill=ctLightBlue] (A) {Algebra};
  \node[qty,fill=ctLightBlue,right=of A] (G) {Geometry};
  \node[qty,fill=ctLightTeal,below=of A] (T) {Trig};
  \node[qty,fill=ctLightTeal,right=of T] (S) {Stoich};
  \node[qty,fill=ctLightGold,below=2cm of T] (F) {Fields};
  \node[qty,fill=ctLightGold,right=of F] (C) {Cats};
  \draw[eqedge] (A)--(G); \draw[eqedge] (A)--(T); \draw[eqedge] (A)--(S);
  \draw[eqedge] (G)--(T); \draw[eqedge] (G)--(S);
  \draw[eqedge] (T)--(F); \draw[eqedge] (S)--(F);
  \draw[eqedge] (F)--(C); \draw[eqedge] (A)--(F);
\end{tikzpicture}
\end{center}

\begin{principle}
Scientific understanding grows by identifying patterns of constraint and
transformation across domains. Transformation reveals structure; invariance
reveals law.
\end{principle}


% ============================================================
% APPENDICES
% ============================================================
\appendix

\chapter{Algebra as Constraint Before Content}

\section{Constraint Systems}

Let $x=(x_1,\dots,x_n)\in\mathbb{R}^n$ denote a vector of system variables.
A \textbf{constraint system} is a collection of equations
\[
F_\alpha(x)=0,\qquad \alpha=1,\dots,m,
\]
where each function $F_\alpha:\mathbb{R}^n\to\mathbb{R}$ specifies a
restriction on the admissible configurations of the variables.

The set of all points satisfying every constraint is the
\textbf{solution set} (or constraint locus)
\[
\mathcal{C}=\{x\in\mathbb{R}^n \mid F_\alpha(x)=0\;\text{for all }\alpha\}.
\]
Geometrically, $\mathcal{C}$ is the intersection of the level sets
$F_\alpha^{-1}(0)$. Each additional constraint restricts the admissible
region of the state space.

\medskip

In the special case where each $F_\alpha$ is linear,
\[
F_\alpha(x)=a_{\alpha1}x_1+\cdots+a_{\alpha n}x_n-b_\alpha,
\]
the system may be written compactly in matrix form as
\[
A\mathbf{x}=\mathbf{b},
\]
where $A\in\mathbb{R}^{m\times n}$ is the coefficient matrix,
$\mathbf{x}\in\mathbb{R}^n$ is the vector of unknowns, and
$\mathbf{b}\in\mathbb{R}^m$ is the vector of constants.

The solution set of a consistent linear system is an affine subspace of
$\mathbb{R}^n$.

\section{Admissibility and Relaxation}

A constraint system may fail to admit an exact solution.
If the number of constraints exceeds the number of variables ($m>n$),
the system is \textbf{overdetermined}. In this case the equations may be
mutually incompatible.

To measure the degree to which a configuration violates the constraints,
define the \textbf{residual functional}
\[
\mathcal{R}(\mathbf{x})
=\sum_{\alpha=1}^{m}\bigl|F_\alpha(\mathbf{x})\bigr|^{2}.
\]

Configurations that minimise $\mathcal{R}$ provide the closest possible
approximate solution. For linear systems this leads to the classical
least–squares formulation.

\begin{proposition}[Least–Squares Solution]
For the linear system $A\mathbf{x}=\mathbf{b}$, the minimiser of
\[
\mathcal{R}(\mathbf{x})=\|A\mathbf{x}-\mathbf{b}\|^{2}
\]
satisfies the \textbf{normal equations}
\[
A^{\top}A\mathbf{x}=A^{\top}\mathbf{b}.
\]
\end{proposition}

\begin{proof}
Let
\[
\mathcal{R}(\mathbf{x})=(A\mathbf{x}-\mathbf{b})^\top(A\mathbf{x}-\mathbf{b}).
\]
Expanding gives
\[
\mathcal{R}(\mathbf{x})
=\mathbf{x}^\top A^\top A\mathbf{x}
-2\mathbf{b}^\top A\mathbf{x}
+\mathbf{b}^\top\mathbf{b}.
\]

Differentiating with respect to $\mathbf{x}$ and setting the gradient to
zero yields
\[
\nabla\mathcal{R}
=2A^\top A\mathbf{x}-2A^\top\mathbf{b}=0,
\]
which simplifies to
\[
A^\top A\mathbf{x}=A^\top\mathbf{b}.
\]
\end{proof}

The least–squares solution therefore represents the configuration that
minimises the total squared violation of the constraints.

\section{Degrees of Freedom}

When the number of variables exceeds the number of independent
constraints ($m<n$), the system is \textbf{underdetermined}. In this case
multiple solutions exist.

Let $r=\mathrm{rank}(A)$. The nullspace
\[
\ker(A)=\{\mathbf{v}\in\mathbb{R}^n\mid A\mathbf{v}=0\}
\]
has dimension $n-r$ by the rank–nullity theorem.

If $\mathbf{x}_0$ is a particular solution of $A\mathbf{x}=\mathbf{b}$,
then the complete set of solutions is
\[
\mathbf{x}=\mathbf{x}_0+\mathbf{v},
\qquad
\mathbf{v}\in\ker(A).
\]

Thus the solution set forms an affine space of dimension
\[
n-\mathrm{rank}(A).
\]

These remaining dimensions represent the \textbf{degrees of freedom} of
the system: directions in which the variables may vary while preserving
all constraints.

In many scientific contexts these degrees of freedom correspond to the
internal flexibility of a system. They determine how the system can
reconfigure itself while remaining consistent with the governing
relations.

\chapter{Geometry as Structured Invariance}

\section{Euclidean Space and the Inner Product}

Let $\mathbb{R}^n$ denote $n$–dimensional Euclidean space.  A point
$x\in\mathbb{R}^n$ is represented by a coordinate vector
$x=(x_1,\dots,x_n)$.

The fundamental geometric structure on $\mathbb{R}^n$ is the
\textbf{inner product}
\[
\langle x,y\rangle=\sum_{i=1}^{n} x_i y_i ,
\]
which assigns to each pair of vectors a real number measuring their
mutual alignment.

From the inner product we obtain the \textbf{norm}
\[
\|x\|=\sqrt{\langle x,x\rangle},
\]
which represents the Euclidean length of a vector.

\medskip

\textbf{Distance.}

The distance between two points $x,y\in\mathbb{R}^n$ is defined by
\[
d(x,y)=\|x-y\|.
\]

Expanding the square of the norm gives
\[
\|x-y\|^2
=\langle x-y,x-y\rangle
=\langle x,x\rangle
-2\langle x,y\rangle
+\langle y,y\rangle .
\]

In two dimensions this identity yields the familiar
\[
(x_1-y_1)^2+(x_2-y_2)^2,
\]
which is the algebraic form of the Pythagorean theorem.

\medskip

\textbf{Angles.}

The inner product also determines the angle between vectors.  For
nonzero vectors $x$ and $y$, define
\[
\cos\theta=\frac{\langle x,y\rangle}{\|x\|\|y\|}.
\]

Orthogonality corresponds to the condition
\[
\langle x,y\rangle=0 .
\]

Thus Euclidean geometry can be understood as the study of invariants
arising from the inner product structure.

\section{Linear Transformations and Invariance}

A linear map $T:\mathbb{R}^n\to\mathbb{R}^n$ preserves Euclidean
geometry if it preserves the inner product:
\[
\langle Tx,Ty\rangle=\langle x,y\rangle
\quad\text{for all }x,y.
\]

Such transformations form the \textbf{orthogonal group}
\[
O(n)=\{A\in\mathbb{R}^{n\times n}\mid A^\top A=I\}.
\]

Matrices in $O(n)$ preserve lengths and angles:
\[
\|Ax\|=\|x\|.
\]

Rotations and reflections therefore represent the fundamental symmetry
operations of Euclidean space.

\section{Configuration Spaces}

Many physical and mathematical systems cannot be described by a finite
number of coordinates.  Instead, their states are represented by
functions defined over a spatial region $\Omega$.

In such cases the state space becomes infinite–dimensional.

For the RSVP framework, the system state consists of three fields:
a scalar density $\Phi$, a vector flow field $\mathbf{v}$, and an
entropy field $S$.  The configuration space is therefore

\[
\mathcal{M}=
\left\{
(\Phi,\mathbf{v},S)
\;\middle|\;
\Phi:\Omega\to\mathbb{R},
\;
\mathbf{v}:\Omega\to\mathbb{R}^3,
\;
S:\Omega\to\mathbb{R}
\right\}.
\]

Each element of $\mathcal{M}$ represents a complete spatial
configuration of the system.

\section{Hilbert Space Structure}

To extend Euclidean geometry to this setting, we equip $\mathcal{M}$ with
an $L^2$ inner product.

Let
\[
X=(\Phi,\mathbf{v},S),
\qquad
Y=(\Phi',\mathbf{v}',S').
\]

Define
\[
\langle X,Y\rangle
=
\int_{\Omega}
\left(
\Phi\,\Phi'
+
\mathbf{v}\cdot\mathbf{v}'
+
S\,S'
\right)\,dx .
\]

The corresponding norm is
\[
\|X\|
=
\sqrt{\langle X,X\rangle}.
\]

With this structure, $\mathcal{M}$ becomes a Hilbert space.

\medskip

\textbf{Distances between configurations.}

The distance between two field configurations $X$ and $Y$ is

\[
d(X,Y)=\|X-Y\|
=
\left(
\int_{\Omega}
\big(
(\Phi-\Phi')^2
+
\|\mathbf{v}-\mathbf{v}'\|^2
+
(S-S')^2
\big)
dx
\right)^{1/2}.
\]

Thus two configurations are close if their scalar, vector, and entropy
fields differ only slightly across the domain.

\section{Geometric Interpretation}

The transition from $\mathbb{R}^n$ to the configuration space
$\mathcal{M}$ represents a shift from finite–dimensional geometry to the
geometry of function spaces.  In finite dimensions, a vector is an
ordered tuple of numbers describing a point or direction in space.
In the configuration space $\mathcal{M}$, the role of vectors is played
by fields defined over the spatial domain $\Omega$.  Thus a system state
is no longer specified by finitely many coordinates but by functions
that assign values to every point in the domain.

The algebraic operations familiar from Euclidean geometry extend
naturally to this setting.  The dot product between vectors becomes a
spatial integral of pointwise products of fields.  The length of a
vector becomes the norm of a field, obtained by integrating the square
of the field magnitude over the domain.  Distances between points in
$\mathbb{R}^n$ generalize to distances between configurations, measured
by the integral of squared differences between the corresponding
fields.

In this way the structure of Euclidean geometry is preserved while the
objects under consideration become functions rather than finite
coordinate vectors.  The resulting space retains the essential
geometric notions of length, distance, orthogonality, and projection,
but these notions now describe relationships between entire field
configurations.  Such infinite–dimensional geometric spaces provide the
natural mathematical setting for continuum physics, variational
principles, and field theories.

\medskip

The study of such configuration spaces forms the mathematical foundation
of modern field theories, variational principles, and continuum models.
Within this framework, physical dynamics can be interpreted as flows
through an infinite–dimensional geometric space whose structure is
determined by the underlying invariants of the system.

\chapter{Trigonometry as Oscillation, Rotation, and Recurrence}

\section{The Rotation Group}

Consider the linear transformation of the plane defined by the matrix
\[
R(\theta)=
\begin{pmatrix}
\cos\theta & -\sin\theta \\
\sin\theta & \cos\theta
\end{pmatrix},
\qquad \theta\in\mathbb{R}.
\]
For any vector $x=(x_1,x_2)\in\mathbb{R}^2$, the action of $R(\theta)$
rotates the vector counterclockwise through an angle $\theta$ about the
origin. Explicitly,
\[
R(\theta)x=
\begin{pmatrix}
x_1\cos\theta-x_2\sin\theta \\
x_1\sin\theta+x_2\cos\theta
\end{pmatrix}.
\]

The matrices $R(\theta)$ form a one–parameter subgroup of the general
linear group $\mathrm{GL}(2,\mathbb{R})$.  Under matrix multiplication
they satisfy the composition rule
\[
R(\theta_1)R(\theta_2)=R(\theta_1+\theta_2),
\]
which follows directly from the trigonometric addition formulas
\[
\cos(\theta_1+\theta_2)
=\cos\theta_1\cos\theta_2-\sin\theta_1\sin\theta_2,
\qquad
\sin(\theta_1+\theta_2)
=\sin\theta_1\cos\theta_2+\cos\theta_1\sin\theta_2.
\]

The identity element corresponds to $\theta=0$,
\[
R(0)=I,
\]
and the inverse of a rotation is given by
\[
R(\theta)^{-1}=R(-\theta).
\]

The set
\[
\mathrm{SO}(2)
=
\{R(\theta)\mid\theta\in\mathbb{R}\}
\]
is therefore a Lie group, known as the \textit{special orthogonal group}
in two dimensions.

\medskip

A defining property of $\mathrm{SO}(2)$ is the preservation of the
Euclidean inner product.  For any vectors $x,y\in\mathbb{R}^2$,
\[
\langle R(\theta)x,R(\theta)y\rangle
=
\langle x,y\rangle .
\]
Equivalently,
\[
R(\theta)^\top R(\theta)=I.
\]

This orthogonality condition ensures that rotations preserve both
lengths and angles.  In particular,
\[
\|R(\theta)x\|=\|x\|.
\]

The determinant of the rotation matrix is
\[
\det R(\theta)
=
\cos^2\theta+\sin^2\theta.
\]
Since $\det R(\theta)=1$ for all $\theta$, the trigonometric identity
\[
\sin^2\theta+\cos^2\theta=1
\]
can be interpreted as the algebraic condition ensuring that rotations
preserve oriented area.

Thus the fundamental identity of trigonometry expresses the invariance
property defining the rotation group.

\section{Harmonic Analysis}

Trigonometric functions also arise as solutions of fundamental
differential equations describing oscillatory motion.

Consider the second–order linear ordinary differential equation
\[
\ddot{x}+\omega^2 x=0,
\qquad \omega>0.
\]
This equation governs the dynamics of simple harmonic oscillators such
as springs, pendula (for small amplitudes), and many wave phenomena.

To solve the equation, assume a solution of exponential form
\[
x(t)=e^{\lambda t}.
\]
Substitution yields the characteristic equation
\[
\lambda^2+\omega^2=0,
\]
whose roots are
\[
\lambda=\pm i\omega.
\]

Therefore the general complex solution is
\[
x(t)=A e^{i\omega t}+B e^{-i\omega t}.
\]

Using Euler's formula
\[
e^{i\omega t}=\cos(\omega t)+i\sin(\omega t),
\]
the real solutions can be written in trigonometric form
\[
x(t)=C\cos(\omega t)+D\sin(\omega t).
\]

These functions describe periodic oscillations with angular frequency
$\omega$ and period
\[
T=\frac{2\pi}{\omega}.
\]

\medskip

The connection between rotations and oscillations becomes evident when
the system is expressed as a first–order dynamical system.  Define the
state vector
\[
\mathbf{z}(t)=
\begin{pmatrix}
x(t) \\
\dot{x}(t)
\end{pmatrix}.
\]
Then the harmonic oscillator equation can be written as
\[
\dot{\mathbf{z}}=
\begin{pmatrix}
0 & 1 \\
-\omega^2 & 0
\end{pmatrix}
\mathbf{z}.
\]

The solution of this linear system is
\[
\mathbf{z}(t)=e^{tA}\mathbf{z}(0),
\qquad
A=
\begin{pmatrix}
0 & 1 \\
-\omega^2 & 0
\end{pmatrix}.
\]

The matrix exponential $e^{tA}$ generates a continuous rotation in the
phase plane $(x,\dot{x})$, showing that harmonic motion corresponds
geometrically to uniform circular motion in phase space.

\section{Plane Waves and Higher–Dimensional Oscillation}

In spatially extended systems the same oscillatory structure appears in
the form of traveling waves.

Let $\mathbf{x}\in\mathbb{R}^3$ denote spatial position and $t$ denote
time.  A plane wave solution takes the form
\[
\psi(\mathbf{x},t)
=
e^{i(\mathbf{k}\cdot\mathbf{x}-\omega t)},
\]
where $\mathbf{k}\in\mathbb{R}^3$ is the wavevector and $\omega$ is the
angular frequency.

Taking real parts yields the familiar trigonometric representation
\[
\psi(\mathbf{x},t)
=
\cos(\mathbf{k}\cdot\mathbf{x}-\omega t).
\]

These solutions arise naturally when linearizing many physical systems
around a uniform equilibrium state.  Small perturbations decompose into
superpositions of plane waves, each characterized by a wavevector
$\mathbf{k}$ and frequency $\omega$.

Within the RSVP framework, linearization of the field equations around a
uniform background state produces perturbative modes of precisely this
form,
\[
e^{i(\mathbf{k}\cdot\mathbf{x}-\omega t)}.
\]
The trigonometric recurrence underlying elementary oscillations thus
extends naturally to higher–dimensional field configurations, where the
geometry of rotations and periodic motion governs the propagation of
waves through space.

\chapter{Stoichiometry as Conservation Algebra}

\section{The Stoichiometric Null Space}

Consider a chemical reaction network involving $s$ chemical species and
$r$ reactions. Let
\[
\mathbf{c}(t)=(c_1(t),\dots,c_s(t))^\top\in\mathbb{R}^s
\]
denote the vector of species concentrations.

The structure of the reaction network is encoded in the
\textbf{stoichiometric matrix}
\[
S\in\mathbb{Z}^{s\times r}.
\]

The entry $S_{ij}$ represents the net number of molecules of species $i$
produced by reaction $j$. Negative entries correspond to reactants and
positive entries correspond to products.

\medskip

Let $\boldsymbol{\gamma}\in\mathbb{R}^r$ denote a vector of reaction
coefficients. The condition for a balanced reaction is
\[
S\boldsymbol{\gamma}=\mathbf{0}.
\]

Thus admissible reaction coefficients lie in the
\textbf{null space} (kernel) of the stoichiometric matrix:
\[
\ker(S)
=
\{\boldsymbol{\gamma}\in\mathbb{R}^r\mid S\boldsymbol{\gamma}=0\}.
\]

Vectors in $\ker(S)$ correspond to combinations of reactions that
preserve the total count of each conserved element.

\medskip

If the matrix $S$ has rank $k$, then by the rank–nullity theorem
\[
\dim(\ker S)=r-k.
\]

This dimension represents the number of independent conservation
relations within the reaction network.

\medskip

Another important structure is the \textbf{left null space}
\[
\ker(S^\top)
=
\{\mathbf{w}\in\mathbb{R}^s\mid \mathbf{w}^\top S=0\}.
\]

Vectors $\mathbf{w}$ in $\ker(S^\top)$ correspond to conserved
quantities. If $\mathbf{w}^\top S=0$, then
\[
\mathbf{w}^\top\dot{\mathbf{c}}(t)
=
\mathbf{w}^\top S\,\mathbf{r}(\mathbf{c})
=
0.
\]

Thus the quantity $\mathbf{w}^\top\mathbf{c}(t)$ remains constant over
time. These conserved quantities correspond to elemental conservation
laws or other invariant totals.

\section{Dynamic Stoichiometry}

The time evolution of the system is governed by the reaction rate vector
\[
\mathbf{r}(\mathbf{c})
=
(r_1(\mathbf{c}),\dots,r_r(\mathbf{c}))^\top.
\]

Each component $r_j(\mathbf{c})$ specifies the rate at which reaction
$j$ occurs. Under deterministic mass–action kinetics,
\[
r_j(\mathbf{c})
=
k_j \prod_{i=1}^{s} c_i^{\nu_{ij}},
\]
where $\nu_{ij}$ denotes the stoichiometric coefficient of species $i$
in the reactant complex of reaction $j$, and $k_j$ is the reaction rate
constant.

The evolution of the concentration vector is therefore given by
\[
\dot{\mathbf{c}}=S\,\mathbf{r}(\mathbf{c}).
\]

This system of differential equations describes how the concentrations
change as reactions proceed.

\medskip

Each reaction contributes a vector $S_{\cdot j}$ in concentration space,
and the dynamics of the system are obtained by summing these
contributions weighted by their reaction rates.

Thus the stoichiometric matrix determines the directions in which the
system may evolve, while the reaction rate functions determine the
speed of motion along those directions.

\section{Stoichiometric Subspaces}

The trajectories of the dynamical system
\[
\dot{\mathbf{c}}=S\,\mathbf{r}(\mathbf{c})
\]
are constrained to remain within affine subspaces known as
\textbf{stoichiometric compatibility classes}.

Let $\mathbf{c}_0$ be an initial concentration vector. The associated
stoichiometric compatibility class is
\[
\mathbf{c}_0+\mathrm{im}(S)
=
\{\mathbf{c}_0+S\mathbf{y}\mid \mathbf{y}\in\mathbb{R}^r\}.
\]

All solutions starting from $\mathbf{c}_0$ remain within this affine
subspace.

This follows from the conservation relations. If
$\mathbf{w}^\top S=0$, then
\[
\mathbf{w}^\top\mathbf{c}(t)=\mathbf{w}^\top\mathbf{c}_0.
\]

Thus conserved quantities restrict the system to a lower–dimensional
region of concentration space.

\section{Continuum Limit and Field Conservation}

The algebraic structure of stoichiometry extends naturally to
continuum systems.

Let $\rho(x,t)$ denote the density of a conserved quantity within a
spatial domain $\Omega\subset\mathbb{R}^3$. The conservation principle
states that the rate of change of density within a region equals the
net flux entering or leaving that region.

This principle leads to the \textbf{continuity equation}
\[
\partial_t\rho+\nabla\cdot(\rho\mathbf{v})=0,
\]
where $\mathbf{v}(x,t)$ is the velocity field describing transport.

\medskip

The relationship between stoichiometric dynamics and field conservation
can be understood by interpreting the reaction terms as discrete sources
and sinks of matter.

In the discrete case, conservation is expressed by the algebraic
constraint
\[
S\boldsymbol{\gamma}=0.
\]

In the continuum case, conservation becomes a differential constraint on
the density field.

Thus the same structural principle appears in two different mathematical
forms: finite-dimensional linear algebra in reaction networks and
partial differential equations in field theory.

\chapter{RSVP Field Theory: Foundational Derivation}

\section{The Fundamental Fields}

Let $\Omega\subset\mathbb{R}^3$ be a spatial domain with coordinates
$x=(x_1,x_2,x_3)$ and time variable $t\in\mathbb{R}$.  The state of the
system at each spacetime point $(x,t)$ is described by three coupled
fields
\[
\Phi(x,t)\in\mathbb{R},\qquad
\mathbf{v}(x,t)\in\mathbb{R}^3,\qquad
S(x,t)\in\mathbb{R}.
\]

The scalar field $\Phi$ represents the density of the plenum.  The
vector field $\mathbf{v}$ represents the local flow velocity of the
plenum medium.  The scalar field $S$ represents the entropy density,
which governs dissipative relaxation processes within the system.

Mathematically the configuration space of the theory is therefore
\[
\mathcal{M}
=
\left\{
(\Phi,\mathbf{v},S)
\;\middle|\;
\Phi:\Omega\times\mathbb{R}\to\mathbb{R},
\;
\mathbf{v}:\Omega\times\mathbb{R}\to\mathbb{R}^3,
\;
S:\Omega\times\mathbb{R}\to\mathbb{R}
\right\}.
\]

The evolution of the system is governed by coupled partial differential
equations expressing conservation, transport, and thermodynamic
relaxation.

\section{Governing Equations}

The dynamics of the fields are described by the system
\begin{align}
\partial_t\Phi+\nabla\cdot(\Phi\mathbf{v})
&=D_\Phi\Delta\Phi+F_\Phi(\Phi,\mathbf{v},S), \tag{A}\\
\partial_t\mathbf{v}+(\mathbf{v}\cdot\nabla)\mathbf{v}
&=-\nabla P(\Phi,S)+D_v\Delta\mathbf{v}+F_v(\Phi,\mathbf{v},S), \tag{B}\\
\partial_t S+\nabla\cdot(S\mathbf{v})
&=D_S\Delta S+\sigma(\Phi,\mathbf{v},S). \tag{C}
\end{align}

The left–hand sides of these equations contain the transport operators
governing advection of the fields by the velocity $\mathbf{v}$.  The
operator
\[
\partial_t+\mathbf{v}\cdot\nabla
\]
is the material derivative describing the rate of change observed along
flow trajectories.

The diffusion terms
\[
D_\Phi\Delta\Phi,\qquad
D_v\Delta\mathbf{v},\qquad
D_S\Delta S
\]
represent smoothing processes that dissipate spatial gradients.

The nonlinear functions
\[
F_\Phi(\Phi,\mathbf{v},S),\qquad
F_v(\Phi,\mathbf{v},S),\qquad
\sigma(\Phi,\mathbf{v},S)
\]
encode the internal coupling between density, flow, and entropy.

Equation (A) expresses conservation of plenum density with diffusive and
source corrections.  Equation (B) governs the momentum balance of the
flow field and generalizes the Euler–Navier–Stokes structure.  Equation
(C) describes entropy transport and production.

\section{Lyapunov Functional and Entropic Descent}

The dissipative structure of the theory can be formulated using a
Lyapunov functional defined on the configuration space.

Let
\[
\mathcal{F}[\Phi,\mathbf{v},S]
=
\int_\Omega
\left[
\frac12|\mathbf{v}|^2
+
U(\Phi,S)
+
\frac{\kappa_\Phi}{2}|\nabla\Phi|^2
+
\frac{\kappa_S}{2}|\nabla S|^2
\right]dx .
\]

The first term represents kinetic energy associated with the flow
field.  The potential function $U(\Phi,S)$ encodes the thermodynamic
coupling between density and entropy.  The gradient terms penalize
spatial inhomogeneities and represent interfacial or coherence energy.

A dissipative dynamics is obtained by defining the evolution of each
field as a gradient descent in the functional $\mathcal{F}$:
\[
\partial_t X
=
-\frac{\delta\mathcal{F}}{\delta X},
\qquad
X\in\{\Phi,\mathbf{v},S\}.
\]

Under this evolution the functional decreases monotonically:
\[
\frac{d}{dt}\mathcal{F}
=
\int_\Omega
\frac{\delta\mathcal{F}}{\delta X}\,\partial_t X\,dx
=
-\int_\Omega
\left|
\frac{\delta\mathcal{F}}{\delta X}
\right|^2 dx
\le 0.
\]

Thus $\mathcal{F}$ acts as a Lyapunov functional for the dynamics.
Configurations evolve toward states that minimize the free–energy
functional.

Within the conceptual framework of RSVP cosmology this descent
represents the tendency of the plenum to reorganize toward smoother
field configurations of lower free energy.

\section{Linear Perturbation Theory}

To study stability, consider perturbations around a homogeneous
background state
\[
\Phi(x,t)=\Phi_0+\delta\Phi(x,t),
\qquad
\mathbf{v}(x,t)=\delta\mathbf{v}(x,t),
\qquad
S(x,t)=S_0+\delta S(x,t).
\]

Substituting these expressions into the governing equations and
retaining only first–order terms yields a linearized system for the
perturbations.

Solutions of the linearized equations can be expressed as Fourier modes
\[
\delta X(x,t)
=
\hat{X}\,e^{i(\mathbf{k}\cdot x-\omega t)},
\]
where $\mathbf{k}\in\mathbb{R}^3$ is the wavevector and $\omega$ is the
angular frequency.

Substituting this ansatz into the linearized equations produces a
dispersion relation
\[
\omega=\omega(\mathbf{k})
\]
that determines the growth, decay, or propagation of perturbations.

If $\operatorname{Im}(\omega)>0$ the perturbation grows exponentially,
indicating instability of the background state.  If
$\operatorname{Im}(\omega)<0$ the perturbation decays, indicating that
the state is linearly stable.

Thus the dispersion relation characterizes the propagation of waves and
the emergence of structure within the RSVP field system.

\chapter{Entropic Redshift and Non-Expansion Cosmology}

\section{Redshift Without Metric Expansion}

In the standard cosmological model, observed redshift is interpreted as
a consequence of the expansion of the spacetime metric.  Within the RSVP
framework, an alternative mechanism is proposed in which photon energy
attenuates gradually as radiation propagates through a structured
entropic medium.

Let $\mathcal{E}(x,t)$ denote an \textbf{entropic attenuation field}
defined over spacetime.  This field represents the cumulative effect of
microscopic interactions between propagating radiation and the
thermodynamic structure of the plenum.

Consider a photon trajectory parameterized by an affine parameter
$\lambda$.  Let $x(\lambda)$ denote the spacetime path of the photon.
The evolution of the photon frequency $\nu(\lambda)$ is assumed to obey
the differential equation
\[
\frac{d\nu}{d\lambda}
=
-\alpha\,\mathcal{E}(x(\lambda),t)\,\nu,
\]
where $\alpha$ is a coupling constant describing the strength of the
interaction between radiation and the entropic field.

This equation expresses the assumption that frequency attenuation is
proportional both to the instantaneous frequency and to the local
strength of the entropic field.

\section{Solution of the Frequency Transport Equation}

The differential equation
\[
\frac{d\nu}{d\lambda}=-\alpha\,\mathcal{E}(x(\lambda),t)\,\nu
\]
is a first-order linear ordinary differential equation.

Separating variables yields
\[
\frac{d\nu}{\nu}
=
-\alpha\,\mathcal{E}(x(\lambda),t)\,d\lambda.
\]

Integrating along the photon path from $\lambda=0$ to $\lambda$ gives
\[
\int_{\nu_0}^{\nu(\lambda)}\frac{d\nu'}{\nu'}
=
-\alpha\int_0^\lambda
\mathcal{E}(x(\lambda'),t)\,d\lambda'.
\]

Evaluating the integral on the left side yields
\[
\ln\!\left(\frac{\nu(\lambda)}{\nu_0}\right)
=
-\alpha\int_0^\lambda
\mathcal{E}(x(\lambda'),t)\,d\lambda'.
\]

Exponentiating both sides gives the solution
\[
\nu(\lambda)
=
\nu_0
\exp\!\left(
-\alpha
\int_0^\lambda
\mathcal{E}(x(\lambda'),t)\,d\lambda'
\right).
\]

Thus the photon frequency decays exponentially according to the
integrated entropic attenuation encountered along the path.

\section{Definition of Redshift}

Observed cosmological redshift is defined by the ratio of emitted and
observed frequencies:
\[
1+z=\frac{\nu_0}{\nu(\lambda)}.
\]

Substituting the solution for $\nu(\lambda)$ yields
\[
1+z
=
\exp\!\left(
\alpha
\int_0^\lambda
\mathcal{E}(x(\lambda'),t)\,d\lambda'
\right).
\]

Thus the redshift is determined by the accumulated entropic attenuation
along the photon trajectory.

\section{Weak-Field Approximation}

If the attenuation field is approximately homogeneous and weak along the
propagation path, one may approximate
\[
\mathcal{E}(x(\lambda),t)\approx\mathcal{E}_0.
\]

The integral then simplifies to
\[
\int_0^\lambda \mathcal{E}_0\,d\lambda'
=
\mathcal{E}_0\,\lambda.
\]

The redshift relation becomes
\[
1+z
=
\exp(\alpha\mathcal{E}_0\lambda).
\]

For small values of $\alpha\mathcal{E}_0\lambda$, the exponential may be
expanded to first order:
\[
1+z\approx 1+\alpha\mathcal{E}_0\lambda.
\]

Therefore
\[
z\approx\alpha\mathcal{E}_0\lambda.
\]

\section{Emergence of the Linear Distance–Redshift Law}

If the affine parameter $\lambda$ is proportional to physical distance
$d$, the previous relation implies
\[
z\propto d.
\]

This reproduces the empirically observed linear distance–redshift
relation traditionally attributed to cosmological expansion.

In the RSVP framework the same phenomenology emerges without requiring a
global expansion of the spacetime metric.  Instead, the redshift arises
from cumulative entropic attenuation experienced by photons as they
propagate through the structured plenum.

Thus the observed large-scale redshift–distance relation can be
interpreted as a transport phenomenon governed by the interaction
between radiation and the entropy field of the cosmological medium.

\chapter{TARTAN: Recursive Tiling and Annotated Noise}

\section{Multiscale Partition}

Let $\Omega\subset\mathbb{R}^d$ be a spatial domain.  The TARTAN framework
represents the system using a hierarchy of spatial partitions that
resolve the domain at progressively finer scales.

Formally, let
\[
\mathcal{T}_0 \supseteq \mathcal{T}_1 \supseteq \mathcal{T}_2 \supseteq \cdots
\]
be a nested sequence of partitions of $\Omega$.  Each $\mathcal{T}_k$
is a tiling of $\Omega$ into a finite or countable set of tiles
\[
\mathcal{T}_k=\{T_{k,i}\}_{i\in I_k},
\qquad
\Omega=\bigcup_{i\in I_k} T_{k,i},
\qquad
T_{k,i}\cap T_{k,j}=\varnothing\ \text{for }i\neq j.
\]

The refinement structure is encoded by projection maps
\[
\pi_k:\mathcal{T}_{k+1}\to\mathcal{T}_k,
\]
which associate each fine-scale tile with the unique coarser tile that
contains it.  Thus if $T'\in\mathcal{T}_{k+1}$ then
\[
T'\subseteq \pi_k(T').
\]

This hierarchical structure defines a multiscale representation of the
domain in which spatial resolution increases with the level index $k$.

\medskip

The collection
\[
(\mathcal{T}_k,\pi_k)
\]
therefore forms an inverse system of partitions.  In the limit
$k\to\infty$, the tiles converge toward infinitesimal regions of the
underlying continuum.

\section{Annotated Field States}

Each tile $T\in\mathcal{T}_k$ carries a local state vector
\[
X_k(T)
=
(\Phi_k(T),\mathbf{v}_k(T),S_k(T),\eta_k(T)).
\]

The first three components correspond to discretized values of the RSVP
fields: plenum density $\Phi$, vector flow $\mathbf{v}$, and entropy
density $S$.  The additional variable $\eta_k(T)$ represents a semantic
annotation attached to the tile.

The annotation variable records contextual information associated with
the state of the tile.  Examples include trajectory provenance,
classification labels, or metadata describing the origin of the local
configuration.  This component enables the framework to propagate not
only physical state variables but also higher-level descriptive
information through the simulation hierarchy.

\section{Recursive Field Updates}

State evolution occurs through recursive updates across scales.

Let $T'\in\mathcal{T}_{k+1}$ be a tile at level $k+1$.  Its parent tile at
the coarser level is $\pi_k(T')\in\mathcal{T}_k$.  The state of $T'$ is
computed by applying an update operator
\[
F_k
:
\mathcal{X}_k \times \Xi_k
\to
\mathcal{X}_{k+1},
\]
where $\mathcal{X}_k$ denotes the state space at level $k$ and
$\Xi_k$ denotes the space of stochastic perturbations.

The recursive update rule is therefore
\[
X_{k+1}(T')
=
F_k\!\left(
X_k(\pi_k(T')),
\,\xi_k(T')
\right).
\]

Thus the state of a fine-scale tile depends on the state of its parent
tile together with a stochastic perturbation term.

\section{Noise Decomposition}

The perturbation term
\[
\xi_k(T')
\]
represents stochastic fluctuations introduced during the refinement
process.  This noise is decomposed into two components,
\[
\xi
=
\xi_{\text{white}}
+
\xi_{\text{structured}}.
\]

The white component represents uncorrelated random fluctuations,
typically modeled as Gaussian noise.

The structured component represents correlated perturbations that
encode semantic or dynamical information.  In the TARTAN framework these
perturbations carry metadata describing the origin or interpretation of
the noise source.  Consequently the stochastic forcing term may contain
information about previously observed patterns or constraints within
the evolving system.

\section{Trajectory Memory}

In addition to spatial hierarchy, the system retains a record of
previous states through a trajectory buffer.

For each tile $T$ and time $t$, define the history buffer
\[
\mathcal{B}(T,t)
=
\{X(T,t-\tau_j)\}_{j=1}^{L},
\]
where $\tau_1,\dots,\tau_L$ denote discrete time offsets.

The evolution of the tile state is then given by the update operator
\[
X(T,t+\Delta t)
=
G\!\left(
X(T,t),
\mathcal{B}(T,t),
\xi(T,t)
\right),
\]
where $G$ represents the local dynamical rule governing the system.

This formulation introduces explicit history dependence into the update
process.  The future state of the system therefore depends not only on
its present configuration but also on a finite window of past states.

\section{Multiscale Dynamics}

The combination of recursive spatial refinement and trajectory memory
produces a dynamical system operating simultaneously across multiple
scales.

At coarse levels the model captures large-scale field structure,
whereas finer levels resolve local variations and noise-driven
fluctuations.  The refinement maps $\pi_k$ ensure consistency between
levels, while the update operators $F_k$ propagate information from
coarser to finer resolutions.

This architecture allows the TARTAN framework to represent complex
systems in which large-scale patterns emerge from locally stochastic
processes while preserving traceable provenance information throughout
the hierarchy.

\chapter{Simulated Agency as Sparse Projection}

\section{State Spaces and the Projection Operator}

Let $\mathcal{W}$ denote the space of possible world states and
$\mathcal{M}$ the space of internal model states available to an agent.
Elements $w\in\mathcal{W}$ represent configurations of the external
environment, whereas elements $m\in\mathcal{M}$ represent internal
representations constructed by the agent.

In general $\mathcal{W}$ is extremely high–dimensional and may describe
a complex physical or informational environment.  The internal model
space $\mathcal{M}$ is typically of much lower dimension, reflecting the
limited representational capacity of the agent.

The process of internal representation is therefore described by a
projection map
\[
P:\mathcal{W}\to\mathcal{M},
\qquad
m=P(w).
\]

This mapping compresses the full world state into a reduced internal
representation that retains only task-relevant structure.

\medskip

To formalize the notion of optimal representation, consider an objective
functional balancing two competing criteria: fidelity to the task and
economy of representation.

Let $L_{\text{task}}(m;w)$ denote a loss function measuring how well the
internal model $m$ captures the aspects of the world state $w$ relevant
to the agent's task.  Let $C(m)$ denote a complexity measure associated
with the representation itself.  The parameter $\beta>0$ controls the
trade–off between these two terms.

The optimal internal representation is therefore defined by
\[
m^*
=
\arg\min_{m\in\mathcal{M}}
\bigl[
L_{\text{task}}(m;w)
+
\beta\,C(m)
\bigr].
\]

This formulation expresses the principle of sparse projection: the agent
seeks the simplest internal model capable of performing the required
task with acceptable accuracy.

\medskip

The complexity term $C(m)$ may take several forms depending on the
representation scheme.  Examples include the dimensionality of the
representation, the number of active features, or an information–theoretic
measure such as description length.

\section{The Agency Loop}

The interaction between an agent and its environment forms a dynamical
feedback loop linking world states, internal representations, and
actions.

Let $F:\mathcal{W}\times\mathcal{A}\to\mathcal{W}$ denote the world
transition function, where $\mathcal{A}$ is the action space.  The
environment evolves according to
\[
w_{t+1}=F(w_t,a_t).
\]

At each time step the agent constructs an internal representation
\[
m_t=P(w_t).
\]

An action is then selected according to a policy
\[
\pi:\mathcal{M}\to\mathcal{A},
\qquad
a_t=\pi(m_t).
\]

Combining these relations yields the closed-loop system
\[
\begin{aligned}
w_{t+1} &= F(w_t,a_t), \\
m_t &= P(w_t), \\
a_t &= \pi(m_t).
\end{aligned}
\]

Substituting the last two equations into the first gives the composite
dynamical system
\[
w_{t+1}
=
F\bigl(w_t,\pi(P(w_t))\bigr).
\]

Thus the evolution of the world state depends on the composition of the
projection operator and the policy function.

\section{Dynamical Structure of the Agency Loop}

The coupled variables $(w_t,m_t)$ define a dynamical system on the
product space
\[
\mathcal{W}\times\mathcal{M}.
\]

A fixed point of this system is a pair $(w^*,m^*)$ satisfying
\[
w^*=F(w^*,\pi(m^*)),
\qquad
m^*=P(w^*).
\]

Such points correspond to stable configurations in which the agent's
actions maintain the environment in a steady state.

More generally, the system may admit periodic orbits
\[
(w_{t+k},m_{t+k})=(w_t,m_t),
\]
which correspond to recurring behavioural patterns or habits.

In this sense the agency loop forms a nonlinear dynamical system whose
qualitative behavior may be analyzed using the geometric methods
introduced earlier in the text.  Fixed points represent equilibria,
periodic trajectories correspond to cyclical behaviors, and more
complex attractors represent adaptive or exploratory dynamics within the
agent–environment interaction.

\chapter{Category-Theoretic Semantics of Constraint Systems}

\section{Categories}

Category theory provides a structural language for describing systems of
objects and the transformations between them.  Rather than focusing on
the internal composition of objects, category theory emphasizes the
relationships that connect them.

A \textbf{category} $\mathcal{C}$ consists of the following data.

First, there is a class of objects denoted $\mathrm{Ob}(\mathcal{C})$.
Objects represent the entities under consideration in the theory.

Second, for every ordered pair of objects $A,B\in\mathrm{Ob}(\mathcal{C})$
there is a set of morphisms
\[
\mathrm{Hom}_{\mathcal{C}}(A,B),
\]
whose elements are maps $f:A\to B$.  These morphisms represent
structure–preserving transformations between objects.

Third, morphisms admit a composition operation.  For any triple of
objects $A,B,C$ there is a map
\[
\circ:
\mathrm{Hom}(B,C)\times\mathrm{Hom}(A,B)
\to
\mathrm{Hom}(A,C)
\]
defined by
\[
(g,f)\mapsto g\circ f .
\]

This composition must satisfy the associativity condition
\[
h\circ(g\circ f)=(h\circ g)\circ f
\]
for all composable morphisms $f,g,h$.

Finally, for each object $A$ there exists an identity morphism
\[
\mathrm{id}_A:A\to A
\]
such that
\[
\mathrm{id}_B\circ f=f,
\qquad
f\circ \mathrm{id}_A=f
\]
for every morphism $f:A\to B$.

These axioms formalize the notion of composable transformations within a
structured system.

\medskip

Several standard mathematical structures form categories.

In the category $\mathbf{Set}$ the objects are sets and the morphisms
are functions between sets.

In the category $\mathbf{Vect}_k$ the objects are vector spaces over a
field $k$ and the morphisms are linear maps.

In the context of chemical reaction networks one may define a category
$\mathbf{RxnNet}$ in which objects represent configurations of chemical
species and morphisms represent stoichiometric transformations induced
by reactions.

\section{Functors and Structural Preservation}

Mappings between categories are described by functors.  A
\textbf{functor}
\[
F:\mathcal{C}\to\mathcal{D}
\]
assigns to each object $A\in\mathcal{C}$ an object $F(A)\in\mathcal{D}$
and to each morphism $f:A\to B$ a morphism
\[
F(f):F(A)\to F(B).
\]

This assignment preserves the categorical structure.  In particular,
functors satisfy the identities
\[
F(\mathrm{id}_A)=\mathrm{id}_{F(A)}
\]
and
\[
F(g\circ f)=F(g)\circ F(f).
\]

Thus functors transport the compositional structure of one category into
another.

Within constraint-based systems, functors provide a natural mechanism
for mapping between different levels of description.  For example, a
functor may map a physical system to a reduced model representation
while preserving the relationships between states.

\section{Symmetric Monoidal Categories}

Many systems involve the simultaneous composition of independent
components.  This structure is captured by the notion of a
\textbf{symmetric monoidal category}.

A symmetric monoidal category is a category $\mathcal{C}$ equipped with
a bifunctor
\[
\otimes:\mathcal{C}\times\mathcal{C}\to\mathcal{C}
\]
called the tensor product, together with a distinguished unit object
$I$.

The tensor product represents parallel composition.  If $A$ and $B$ are
objects, the object $A\otimes B$ represents the joint system obtained by
combining them.

The monoidal structure satisfies associativity and unit conditions up to
natural isomorphism.  In particular,
\[
(A\otimes B)\otimes C \cong A\otimes(B\otimes C)
\]
and
\[
A\otimes I \cong A.
\]

A symmetry isomorphism
\[
\sigma_{A,B}:A\otimes B\to B\otimes A
\]
ensures that the order of parallel composition is interchangeable.

\section{Constraint Systems as Monoidal Structures}

Constraint systems can be interpreted naturally within this categorical
framework.

Objects represent semantic modules, state spaces, or constraint
structures.  Morphisms represent consistency–preserving transformations
between these structures.  Sequential composition corresponds to the
composition of morphisms, while parallel composition of independent
modules is represented by the tensor product.

This interpretation allows complex systems to be assembled from simpler
components while maintaining explicit structural relationships.

\section{Colimits and Merging of Systems}

In many contexts one wishes to combine several compatible structures
into a single unified object.  Category theory formalizes this operation
using colimits.

Let $D:\mathcal{J}\to\mathcal{C}$ be a diagram in a category
$\mathcal{C}$.  A \textbf{colimit} of this diagram is an object
$\mathrm{colim}\,D$ together with morphisms
\[
\iota_j:D(j)\to \mathrm{colim}\,D
\]
such that the resulting family of maps is universal.

Universality means that for any object $X$ receiving compatible maps
\[
f_j:D(j)\to X
\]
there exists a unique morphism
\[
u:\mathrm{colim}\,D\to X
\]
satisfying
\[
f_j=u\circ\iota_j.
\]

In the context of constraint systems, colimits describe the operation of
merging multiple modules or subsystems that share compatible boundary
conditions.  The resulting object represents the minimal structure that
integrates all participating components while preserving their
relationships.

Thus categorical constructions provide a rigorous mathematical language
for describing how constraint systems combine, interact, and evolve
within larger structured environments.

\chapter{Sheaf-Theoretic Gluing and Local-to-Global Consistency}

\section{Presheaves}

Let $X$ be a topological space, interpreted as a \emph{context space}
whose open sets represent regions of local observation or computation.
Denote by $\mathrm{Open}(X)$ the partially ordered set of open subsets
of $X$, ordered by inclusion.

A \textbf{presheaf} $\mathcal{F}$ on $X$ assigns to each open set
$U\subseteq X$ a set $\mathcal{F}(U)$, called the set of \emph{sections}
over $U$.  Elements of $\mathcal{F}(U)$ represent pieces of information
defined locally on the region $U$.

For every inclusion of open sets $V\subseteq U$ there is a
\textbf{restriction map}
\[
\rho_{UV}:\mathcal{F}(U)\to\mathcal{F}(V).
\]

These maps must satisfy the compatibility conditions

\[
\rho_{UU}=\mathrm{id}_{\mathcal{F}(U)}
\]

and

\[
\rho_{UW}=\rho_{VW}\circ\rho_{UV}
\qquad
\text{whenever } W\subseteq V\subseteq U.
\]

Thus restricting a section in multiple steps produces the same result as
restricting it directly.

Equivalently, a presheaf may be defined as a contravariant functor

\[
\mathcal{F}:\mathrm{Open}(X)^{\mathrm{op}}\to\mathbf{Set}.
\]

This formulation emphasizes that presheaves encode how local data vary
over a topological space.

\section{Sheaf Condition}

A presheaf becomes a \textbf{sheaf} when local data determine global
data in a consistent way.

Let $\{U_i\}_{i\in I}$ be an open cover of a region
\[
U=\bigcup_{i\in I}U_i.
\]

Suppose we are given sections
\[
s_i\in\mathcal{F}(U_i)
\]
for each open set in the cover.

These sections are said to be \textbf{compatible} if they agree on all
pairwise overlaps.  Formally this means

\[
\rho_{U_i,U_i\cap U_j}(s_i)
=
\rho_{U_j,U_i\cap U_j}(s_j)
\qquad
\text{for all }i,j.
\]

A presheaf $\mathcal{F}$ is a \textbf{sheaf} if the following two
conditions hold.

First, whenever compatible sections $\{s_i\}$ are given on an open
cover, there exists a section

\[
s\in\mathcal{F}(U)
\]

such that

\[
\rho_{U,U_i}(s)=s_i
\qquad
\text{for all }i.
\]

Second, this global section is unique.

Thus a sheaf ensures that locally consistent pieces of information can
be uniquely glued together into a global description.

\section{Gluing and Consistency}

The sheaf condition formalizes a fundamental principle of many
scientific and mathematical systems: global structure emerges from
locally compatible data.

Local observations or computations often describe only a portion of a
system.  When these pieces agree on their regions of overlap, they can
be combined to produce a consistent global model.

Mathematically, the gluing operation is represented by the existence of
the section $s$ satisfying

\[
\rho_{U,U_i}(s)=s_i.
\]

The compatibility condition ensures that no contradictions arise on
overlapping regions.

\section{Applications to Constraint Systems}

Sheaf structures provide a natural framework for modeling systems in
which information is distributed across overlapping contexts.

In geometric problems, sections may represent functions defined on
local coordinate patches.  In physical field theories, sections may
represent local solutions of differential equations.  In networked
systems, sections may encode local constraint solutions associated with
subregions of the system.

The sheaf condition guarantees that if these local solutions agree where
their domains overlap, they determine a unique global solution.

\section{Local-to-Global Reasoning}

Many arguments throughout mathematics and science implicitly follow the
local-to-global paradigm captured by sheaf theory.

Local measurements constrain global physical models.  Local geometric
relations determine global shapes.  Local conservation laws determine
global system behavior.  Even in elementary contexts such as triangle
geometry or stoichiometric balance, reasoning often proceeds by
verifying consistency in small parts and then assembling the result into
a coherent whole.

\begin{principle}
Sheaf coherence provides the formal mathematical expression of
local-to-global consistency.  When local structures agree on their
overlaps, they uniquely determine a global structure.  This principle
underlies a wide range of scientific reasoning, from geometric
construction and field theory to network dynamics and constraint
systems.
\end{principle}

\chapter{Variational Principles and Entropic Descent}

\section{Variational Functionals}

Many mathematical and physical systems can be described by a functional
whose extrema determine the governing equations of the system.

Let $\Omega\subset\mathbb{R}^n$ be a spatial domain and let
$u:\Omega\to\mathbb{R}$ be a sufficiently smooth function.  Consider the
functional
\[
\mathcal{J}[u]
=
\int_\Omega
\mathcal{L}(u,\nabla u,x)\,dx,
\]
where $\mathcal{L}$ is called the \textbf{Lagrangian density}.  The
function $\mathcal{L}$ depends on the value of the field $u$, its
gradient $\nabla u$, and possibly explicitly on the spatial coordinate
$x$.

The central problem of the calculus of variations is to determine the
functions $u$ that make the functional $\mathcal{J}$ stationary.

\section{First Variation and Euler--Lagrange Equation}

To derive the governing equation, consider a perturbation of the field
\[
u_\epsilon(x)=u(x)+\epsilon\,\eta(x),
\]
where $\eta$ is an arbitrary smooth test function vanishing on the
boundary of $\Omega$.

The variation of the functional is
\[
\frac{d}{d\epsilon}\mathcal{J}[u_\epsilon]
=
\int_\Omega
\left(
\frac{\partial\mathcal{L}}{\partial u}\,\eta
+
\frac{\partial\mathcal{L}}{\partial(\nabla u)}
\cdot
\nabla\eta
\right)dx.
\]

Applying integration by parts to the second term yields
\[
\int_\Omega
\frac{\partial\mathcal{L}}{\partial(\nabla u)}\cdot\nabla\eta\,dx
=
-
\int_\Omega
\nabla\cdot
\left(
\frac{\partial\mathcal{L}}{\partial(\nabla u)}
\right)
\eta\,dx,
\]
where boundary terms vanish because $\eta=0$ on $\partial\Omega$.

Thus the first variation becomes
\[
\delta\mathcal{J}
=
\int_\Omega
\left(
\frac{\partial\mathcal{L}}{\partial u}
-
\nabla\cdot
\frac{\partial\mathcal{L}}{\partial(\nabla u)}
\right)
\eta\,dx.
\]

Since $\eta$ is arbitrary, the integrand must vanish.  Therefore
critical points of $\mathcal{J}$ satisfy the
\textbf{Euler--Lagrange equation}
\[
\frac{\partial\mathcal{L}}{\partial u}
-
\nabla\cdot
\left(
\frac{\partial\mathcal{L}}{\partial(\nabla u)}
\right)
=0.
\]

This equation determines the stationary configurations of the
functional.

\section{Gradient Flow Dynamics}

In many systems the dynamics evolve toward a configuration that
minimizes a functional rather than merely satisfying a stationary
condition.

Let $\mathcal{J}[u]$ be a functional defined on a suitable function
space.  The steepest descent dynamics associated with $\mathcal{J}$ are
given by the \textbf{gradient flow equation}
\[
\partial_t u
=
-\frac{\delta\mathcal{J}}{\delta u},
\]
where $\frac{\delta\mathcal{J}}{\delta u}$ denotes the functional
derivative.

The functional derivative is defined by the relation
\[
\delta\mathcal{J}
=
\int_\Omega
\frac{\delta\mathcal{J}}{\delta u}\,
\eta\,dx.
\]

Substituting the expression obtained above gives
\[
\frac{\delta\mathcal{J}}{\delta u}
=
\frac{\partial\mathcal{L}}{\partial u}
-
\nabla\cdot
\left(
\frac{\partial\mathcal{L}}{\partial(\nabla u)}
\right).
\]

Thus the gradient flow equation becomes
\[
\partial_t u
=
-
\left(
\frac{\partial\mathcal{L}}{\partial u}
-
\nabla\cdot
\frac{\partial\mathcal{L}}{\partial(\nabla u)}
\right).
\]

This evolution decreases the value of the functional over time.

Indeed,
\[
\frac{d}{dt}\mathcal{J}[u(t)]
=
\int_\Omega
\frac{\delta\mathcal{J}}{\delta u}\,
\partial_t u\,dx
=
-
\int_\Omega
\left|
\frac{\delta\mathcal{J}}{\delta u}
\right|^2 dx
\le 0.
\]

Therefore $\mathcal{J}$ acts as a Lyapunov functional for the dynamics.

\section{Examples Across Domains}

Variational principles appear in many different mathematical and
scientific contexts.

In linear algebra the least–squares problem
\[
\mathcal{J}(\mathbf{x})
=
\|A\mathbf{x}-\mathbf{b}\|^2
\]
defines a quadratic functional whose minimizer satisfies the normal
equations
\[
A^\top A\mathbf{x}=A^\top\mathbf{b}.
\]

In geometry the length of a curve
\[
\mathcal{J}[\gamma]
=
\int |\dot{\gamma}(t)|\,dt
\]
is minimized by geodesics.

In chemical systems free–energy functionals determine equilibrium
states of reaction networks.

In the RSVP framework the plenum fields evolve toward configurations
that minimize an entropic free–energy functional, producing the
smoothing dynamics discussed in earlier chapters.

\section{Unified Interpretation}

These examples illustrate a general principle.  Many physical,
geometric, and algebraic processes can be interpreted as gradient flows
on appropriately defined functionals.

Algebraic residual minimization, geometric shortest paths, chemical
equilibration, and entropic smoothing all arise from the same
variational structure.  The particular functional $\mathcal{J}$ encodes
the constraints and energies relevant to the system, while the gradient
flow equation determines how the system evolves toward its optimal
configuration.

\chapter{Worked Derivations: Bridges Between the Introductory Subjects and the Theory}

\section{From Linear Algebra to Stoichiometry to Field Theory}

The conceptual thread connecting linear algebra, chemical stoichiometry,
and field theory is the mathematical structure of conservation laws.
Each system expresses the requirement that certain quantities remain
balanced under allowed transformations.

\subsection*{Linear Constraint Formulation}

Consider a chemical reaction network involving a finite set of chemical
species.  Let $\gamma\in\mathbb{R}^r$ denote the vector of reaction
coefficients and let
\[
S\in\mathbb{Z}^{e\times r}
\]
be the \textbf{stoichiometric matrix}, where $e$ is the number of
elements and $r$ the number of species.

Each entry $S_{ij}$ records the number of atoms of element $i$
contained in species $j$.

Balancing a chemical reaction requires that the number of atoms of each
element be conserved.  This condition is expressed as the linear system
\[
S\gamma=\mathbf{0}.
\]

As a simple example consider the reaction
\[
2\mathrm{H_2}+ \mathrm{O_2}\rightarrow 2\mathrm{H_2O}.
\]

The stoichiometric matrix describing hydrogen and oxygen conservation
may be written
\[
S=
\begin{pmatrix}
2 & 0 & -2 \\
0 & 2 & -1
\end{pmatrix}.
\]

The coefficient vector
\[
\gamma=(1,\,\tfrac12,\,1)
\]
lies in the null space of $S$.  Multiplying by two yields the familiar
balanced integer equation
\[
2\mathrm{H_2}+ \mathrm{O_2}\rightarrow 2\mathrm{H_2O}.
\]

Thus chemical balancing is equivalent to computing a basis of the null
space $\ker(S)$.

\subsection*{Dynamic Stoichiometry}

The static conservation relation generalizes naturally to a dynamical
system.

Let
\[
\mathbf{c}(t)=(c_1(t),\dots,c_r(t))
\]
denote the vector of species concentrations and let
\[
\mathbf{r}(\mathbf{c})
\]
be the vector of reaction rates.

The evolution of concentrations is governed by
\[
\dot{\mathbf{c}}=S\,\mathbf{r}(\mathbf{c}).
\]

The columns of $S$ determine how each reaction changes the
concentrations of all species simultaneously.  The kernel of $S$
contains the conserved quantities of the reaction network.

\subsection*{Continuum Limit}

A continuum description emerges when the discrete species
concentrations are replaced by spatial density fields.

Let $\rho(x,t)$ represent the density of a conserved quantity in a
spatial domain $\Omega\subset\mathbb{R}^3$.  The total amount of this
quantity in a region $U\subset\Omega$ is
\[
\int_U \rho(x,t)\,dx.
\]

Conservation requires that the rate of change of this integral equals
the net flux across the boundary $\partial U$.

Applying the divergence theorem yields the
\textbf{continuity equation}
\[
\partial_t\rho+\nabla\cdot(\rho\mathbf{v})=0,
\]
where $\mathbf{v}(x,t)$ is the velocity field transporting the
quantity.

Thus the same conservation structure appears in three progressively
more general settings:

\[
S\gamma=\mathbf{0},
\qquad
\dot{\mathbf{c}}=S\mathbf{r}(\mathbf{c}),
\qquad
\partial_t\rho+\nabla\cdot(\rho\mathbf{v})=0.
\]

The first describes algebraic balance, the second describes reaction
dynamics, and the third describes continuous transport fields.

\section{From Euclidean Distance to the $L^2$ Metric}

Geometry provides another example of how elementary concepts extend
naturally to more advanced mathematical structures.

\subsection*{Euclidean Distance}

In finite-dimensional space $\mathbb{R}^n$, the distance between two
points $P,Q\in\mathbb{R}^n$ is defined by the Euclidean norm
\[
d(P,Q)=\|P-Q\|
=
\left(
\sum_{i=1}^n (P_i-Q_i)^2
\right)^{1/2}.
\]

This definition arises from the inner product
\[
\langle x,y\rangle=\sum_{i=1}^n x_i y_i.
\]

\subsection*{Extension to Function Spaces}

When the objects of interest are functions rather than vectors, the
inner product generalizes to an integral.

For functions $f,g:\Omega\to\mathbb{R}$ one defines the
$L^2$ inner product
\[
\langle f,g\rangle_{L^2}
=
\int_\Omega f(x)g(x)\,dx.
\]

The associated norm is
\[
\|f\|_{L^2}
=
\left(
\int_\Omega |f(x)|^2 dx
\right)^{1/2}.
\]

The distance between two functions therefore becomes
\[
d(f,g)=\|f-g\|_{L^2}
=
\left(
\int_\Omega |f(x)-g(x)|^2 dx
\right)^{1/2}.
\]

\subsection*{Application to Field Configurations}

In the RSVP framework, a system state is represented by a collection of
fields
\[
X=(\Phi,\mathbf{v},S).
\]

The natural metric on the configuration space is the $L^2$ distance
between fields:
\[
d(X_1,X_2)^2
=
\int_\Omega
\left(
|\Phi_1-\Phi_2|^2
+
|\mathbf{v}_1-\mathbf{v}_2|^2
+
|S_1-S_2|^2
\right)dx.
\]

This metric measures the magnitude of the difference between two
possible states of the system.

\section{From Trigonometric Oscillation to RSVP Eigenmodes}

Oscillatory solutions play a central role in both classical mechanics
and field theory.

\subsection*{Simple Harmonic Motion}

The differential equation
\[
\ddot{x}+\omega^2 x=0
\]
describes a harmonic oscillator.

Its general solution is
\[
x(t)=A\cos(\omega t)+B\sin(\omega t).
\]

This equation arises whenever a restoring force is proportional to
displacement.

\subsection*{Fourier Modes}

Many partial differential equations can be analyzed by decomposing
solutions into Fourier modes.

Consider a perturbation
\[
\delta\Phi(x,t).
\]

A Fourier mode has the form
\[
\delta\Phi(x,t)
=
A\,e^{i(\mathbf{k}\cdot x-\omega t)}.
\]

Substituting such modes into a linear partial differential equation
reduces the problem to an algebraic relation between $\omega$ and
$\mathbf{k}$.

This relation is called the \textbf{dispersion relation}.

\subsection*{Application to RSVP Fields}

Linearizing the RSVP equations around a homogeneous background state
produces equations of the form
\[
\partial_t^2(\delta\Phi)+\omega^2(\mathbf{k})\,\delta\Phi=0
\]
for each spatial Fourier mode.

Thus each perturbation mode behaves as a harmonic oscillator with
frequency $\omega(\mathbf{k})$ determined by the dispersion relation.

The spectrum of these oscillatory modes determines how disturbances
propagate through the plenum field.

\section{Unified Perspective}

These derivations illustrate how elementary mathematical concepts
extend naturally into the structures used in advanced theoretical
models.

Linear algebra describes static balance conditions.
Reaction network dynamics extend these balances into time-dependent
systems.
Continuum field equations represent the spatial limit of these same
conservation principles.

Similarly, Euclidean distance generalizes to metrics on function
spaces, and elementary trigonometric oscillations become the Fourier
modes that describe waves and field propagation.

The mathematical structures introduced in the early chapters therefore
form the conceptual foundation for the more advanced theories developed
later in the text.

% ============================================================
% BACK MATTER
% ============================================================
\backmatter

% ---- Glossary -----------------------------------------------
\chapter*{Glossary}
\addcontentsline{toc}{chapter}{Glossary}

\begin{multicols}{2}\small
\begin{description}[labelindent=0pt,leftmargin=0.8em,style=nextline]

\item[Amplitude] maximum displacement of a periodic function from its midline.

\item[Avogadro's number] $N_A=6.022\times10^{23}$; particles per mole.

\item[Category] a mathematical structure consisting of objects and morphisms with associative composition and identity maps.

\item[Configuration space] the set of all possible states of a system.

\item[Constraint] a condition restricting admissible values of variables.

\item[Continuity equation] conservation law $\partial_t\rho+\nabla\cdot(\rho\mathbf{v})=0$ describing transport of density.

\item[Conservation law] a relation asserting that a quantity is preserved under transformation.

\item[Dispersion relation] relation $\omega(\mathbf{k})$ connecting wave frequency and spatial wave number.

\item[Dimensional analysis] tracking units through calculations.

\item[Equilibrium] a fixed point of a dynamical system.

\item[Event-history] a causally ordered sequence of occurrences; primitive in the Spherepop formalism.

\item[Fourier mode] sinusoidal component $e^{i(\mathbf{k}\cdot x-\omega t)}$ used to analyze linear systems.

\item[Function] a rule assigning exactly one output to each input.

\item[Functional derivative] derivative of a functional with respect to a function.

\item[Functor] a structure-preserving map between categories.

\item[Gradient] vector of partial derivatives indicating direction of greatest increase.

\item[Gradient flow] dynamics $\partial_t u=-\delta\mathcal{J}/\delta u$ descending a functional.

\item[Identity (trig.)] an equation true for all valid angles.

\item[Inner product] bilinear operation defining angles and lengths in vector spaces.

\item[Invariant] a quantity unchanged by a transformation.

\item[Isometry] a transformation preserving distances.

\item[Lagrangian] function $\mathcal{L}(u,\nabla u,x)$ whose integral defines a variational functional.

\item[Limiting reagent] the reactant consumed first in a reaction.

\item[Linear system] a system of equations of the form $A\mathbf{x}=\mathbf{b}$.

\item[Lyapunov functional] scalar functional that decreases along system trajectories.

\item[Metric] rule defining distances between elements of a space.

\item[Molar mass] mass per mole of a substance (g/mol).

\item[Mole] SI unit of chemical amount.

\item[Morphism] structure-preserving map between objects in a category.

\item[Null space] $\ker(A)=\{\mathbf{x}\mid A\mathbf{x}=\mathbf{0}\}$.

\item[Period] smallest positive interval of repetition of a periodic function.

\item[Pythagorean theorem] $a^2+b^2=c^2$ for a right triangle.

\item[Radian] angle subtended by arc of length equal to radius.

\item[RSVP] Relativistic Scalar-Vector Plenum field theory.

\item[Sheaf] a presheaf satisfying the gluing axiom for local data.

\item[Similarity] geometric relation in which figures differ only by scale.

\item[Slope] rate of change of a linear function.

\item[Stoichiometric matrix] element-species incidence matrix $S$ encoding conservation.

\item[TARTAN] Recursive Tiling and Annotated Noise simulation framework.

\item[Unit circle] circle of radius 1; foundation of trigonometric definitions.

\item[Variable] a letter representing an unknown or varying quantity.

\item[Variational functional] a function on function spaces, $\mathcal{J}[u]$.

\end{description}
\end{multicols}


% ---- Selected answers ---------------------------------------
\chapter*{Selected Answers}
\addcontentsline{toc}{chapter}{Selected Answers}

\textbf{Chapter 2 (Equations and Balance).}
Practice 2.1: (a) $x=5$; (b) $x=15$; (c) $x=5$; (d) $x=7$.

\textbf{Chapter 7 (Pythagorean Theorem).}
Practice 7.1: diagonal $=\sqrt{49+576}=25$~cm.

\textbf{Chapter 10 (Atoms and Moles).}
Practice 10.1: (a) 58.44; (b) 44.01; (c) 180.16~g/mol.

\textbf{Chapter 17 (Balancing).}
Practice 17.1: (a) $4\mathrm{Fe}+3\mathrm{O_2}\to2\mathrm{Fe_2O_3}$;
(b) $\mathrm{C_3H_8}+5\mathrm{O_2}\to3\mathrm{CO_2}+4\mathrm{H_2O}$.

\textbf{Chapter 18 (Moles).}
Practice 18.1: $\approx33.6$~g CO$_2$.

% ---- Abbreviated Periodic Table ------------------------------
\chapter*{Abbreviated Periodic Table}
\addcontentsline{toc}{chapter}{Abbreviated Periodic Table}

\begin{center}
\renewcommand{\arraystretch}{1.45}
\begin{tabular}{clcc}
\toprule
\textbf{Symbol} & \textbf{Name} & \textbf{$Z$} & \textbf{$M$ (g/mol)} \\
\midrule
H  & Hydrogen  & 1  & 1.008  \\
C  & Carbon    & 6  & 12.011 \\
N  & Nitrogen  & 7  & 14.007 \\
O  & Oxygen    & 8  & 15.999 \\
Na & Sodium    & 11 & 22.990 \\
Mg & Magnesium & 12 & 24.305 \\
S  & Sulfur    & 16 & 32.060 \\
Cl & Chlorine  & 17 & 35.453 \\
K  & Potassium & 19 & 39.098 \\
Ca & Calcium   & 20 & 40.078 \\
Fe & Iron      & 26 & 55.845 \\
Cu & Copper    & 29 & 63.546 \\
\bottomrule
\end{tabular}
\end{center}

% ---- Bibliography -------------------------------------------
\bibliographystyle{plainnat}
\bibliography{constraint_refs}

\end{document}
